%----------------------------------------------------------------------------------------
% RIJEN: HANDBOEK (hoofdbestand). De oude cursus staat lokaal in ../oud/ (niet in git).
%
% Vanuit deze map:
%     py -3.13 -P maak-pdf.py handboek.tex     handboek + numerieke + uitgewerkte oplossingen
%     latexmk -pdf leerkracht.tex              leerkrachtenversie
%     latexmk -pdf zelfstudie.tex              zelfstudieversie
%----------------------------------------------------------------------------------------
\documentclass[11pt,fleqn,oneside,openany]{book}

\newcommand{\cursustitel}{Rijen}
\newcommand{\cursusondertitel}{Handboek}
\newcommand{\cursusauteur}{Heilige-Drievuldigheidscollege Leuven}

% De voorkant: gegevens, kunstwerk en motief. Gemaakt door Ontwerp/voorkanten/maak-voorkanten.py.
\def\cursusrubriek{Discrete wiskunde}
\def\cursusrichtingen{Economie-Wiskunde, Grieks-Wiskunde, Latijn-Wiskunde, Wetenschappen-Wiskunde}
\def\cursusgraad{3de graad doorstroomfinaliteit, 6de jaar}
\def\cursusuren{6 uur per week}
\def\omslagtitel{Rijen}
\def\omslagtitelgrootte{64}


%----------------------------------------------------------------------------------------
% OPMAAK VAN PROTOTYPE 2
% Compact handboek voor in de les. Uit dezelfde bron komen drie versies:
%   \def\versie{handboek}    kern, uitgewerkte voorbeelden, oefeningen zonder uitwerking
%   \def\versie{leerkracht}  handboek + didactische tips, antwoorden
%   \def\versie{zelfstudie}  handboek + extra uitleg, extra opgaven, antwoorden
%----------------------------------------------------------------------------------------
\providecommand{\versie}{handboek}

\usepackage{etoolbox}
\newif\ifhandboek \newif\ifleerkracht \newif\ifzelfstudie \newif\ifuitwerking
\ifdefstring{\versie}{handboek}{\handboektrue}{}
\ifdefstring{\versie}{leerkracht}{\leerkrachttrue\uitwerkingtrue}{}
\ifdefstring{\versie}{zelfstudie}{\zelfstudietrue\uitwerkingtrue}{}

\usepackage[dutch]{babel}
\usepackage[a4paper,left=2.5cm,right=2.5cm,top=2.4cm,bottom=2.2cm,headsep=6mm,footskip=8mm,headheight=14pt]{geometry}

\usepackage{amsmath}
% Source Sans voor tekst en koppen, newtxsf voor schreefloze wiskunde (beide onder de SIL Open Font License)
\usepackage[T1]{fontenc}
\usepackage[default]{sourcesanspro}
\usepackage{newtxsf}
\renewcommand{\familydefault}{\sfdefault}
\linespread{1.1}
\usepackage{icomma}
\usepackage{ragged2e}

\usepackage{array,booktabs}
\usepackage[inline]{enumitem}
\usepackage{multicol}
\usepackage{graphicx}
\graphicspath{{afbeeldingen/}}
\usepackage{xcolor}
\usepackage{tikz}
\usetikzlibrary{arrows.meta,decorations.pathreplacing,calc}
\usepackage[most]{tcolorbox}
\usepackage{fancyhdr}
\usepackage{answers}
\usepackage{needspace}
% Zoals \needspace, maar zonder voorkeur (\penalty-100) om bij elke kop te breken: anders blijven pagina's half leeg
\makeatletter
\renewcommand{\needspace}[1]{\par\penalty0\begingroup\setlength{\dimen@}{#1}%
  \dimen@ii\pagegoal\advance\dimen@ii-\pagetotal
  \ifdim\dimen@>\dimen@ii\break\fi\endgroup}
\makeatother
\usepackage{refcount}

% Elk hoofdstuk heeft een kleur. Donker (hd) = uitleg, fel (hf) = opdrachten.
% Contrast nagekeken: wit op donker minstens 5,3 en inkt op fel minstens 7,0.
\definecolor{hd1}{HTML}{1257B0}\definecolor{hf1}{HTML}{5CB8FF}% blauw
\definecolor{hd2}{HTML}{0A7A5C}\definecolor{hf2}{HTML}{2EDDA5}% groen
\definecolor{hd3}{HTML}{B24A06}\definecolor{hf3}{HTML}{FFA62B}% oranje
\definecolor{hd4}{HTML}{B0216F}\definecolor{hf4}{HTML}{FF7DBD}% roze
\definecolor{hd5}{HTML}{6A35B8}\definecolor{hf5}{HTML}{BC96FF}% paars
\definecolor{hd6}{HTML}{BF2626}\definecolor{hf6}{HTML}{FF8A7A}% rood
\definecolor{hd7}{HTML}{0A6E8A}\definecolor{hf7}{HTML}{3FD6F0}% cyaan
\definecolor{hd8}{HTML}{6E6300}\definecolor{hf8}{HTML}{FFD93D}% olijf
\definecolor{inkt}{HTML}{12202B}
% Rubrieken, voor de labels van Koppel
\colorlet{rubriek-analyse}{hd1}\colorlet{rubriek-algebra}{hd5}\colorlet{rubriek-meetkunde}{hd3}
\colorlet{rubriek-discreet}{hd2}\colorlet{rubriek-data}{hd4}
\definecolor{accent}{HTML}{0B6470}% kleur van de voorkaft en de oplossingen
\colorlet{fel}{hf7}
\definecolor{grijs}{HTML}{5F6368}
\definecolor{lijngrijs}{HTML}{9AA0A6}
\newcommand{\kopletter}{\sffamily\bfseries}
\newif\ifopdrachtdeel % vanaf de Opdrachten: tab, kop en sectienummers in de felle tint
\DeclareEmphSequence{\color{accent}\bfseries,\itshape,\upshape}
\newcommand{\eur}{\texteuro{}}

%----------------------------------------------------------------------------------------
% ALINEA'S
%----------------------------------------------------------------------------------------
\setlength{\parindent}{0pt}
\setlength{\parskip}{5pt plus 1pt}
\AtBeginDocument{\RaggedRight\frenchspacing
  \thickmuskip=5mu plus 1mu\relax\medmuskip=4mu plus 1mu minus 2mu\relax
  \setlength{\abovedisplayskip}{5pt plus 1pt}\setlength{\belowdisplayskip}{5pt plus 1pt}%
  \setlength{\abovedisplayshortskip}{2pt}\setlength{\belowdisplayshortskip}{2pt}}
\tolerance=800
\hyphenpenalty=2000
\clubpenalty=10000 \widowpenalty=10000
\setlength{\mathindent}{3mm}
\raggedbottom
\setlength{\columnsep}{7mm}
\setlength{\multicolsep}{3pt}
\makeatletter
\def\@minipagerestore{\RaggedRight\setlength{\parskip}{3pt}}
\makeatother

\setlist[itemize,1]{label=\textcolor{accent}{\textbullet},leftmargin=4mm,itemsep=.1ex,topsep=.3ex,parsep=0pt}
\setlist[enumerate,1]{label=\textbf{\alph*.},leftmargin=6mm,labelsep=1.4mm,itemsep=.5ex,topsep=.6ex,parsep=0pt,beginpenalty=10000,midpenalty=10000}% deelvragen blijven samen
% Korte deelvragen naast elkaar, uitgelijnd (bv. rijtjes getallen). Lange deelvragen: enumerate.
\newenvironment{reeks}[1][3]{\begin{multicols}{#1}\begin{enumerate}[midpenalty=0]}{\end{enumerate}\end{multicols}}% midpenalty=0: de kolommen moeten kunnen breken

%----------------------------------------------------------------------------------------
% KOPPEN
%----------------------------------------------------------------------------------------
\makeatletter
\def\@makechapterhead#1{%
  \begin{tikzpicture}[remember picture,overlay]
    \fill[accent] (current page.north west) rectangle ($(current page.north east)+(0,-3.5cm)$);
    \node[anchor=west,inner sep=0pt,text=fel,font=\kopletter\fontsize{54}{54}\selectfont]
      at ($(current page.north west)+(2.5cm,-1.95cm)$) {\thechapter};
    \node[anchor=south west,inner sep=0pt,text=white,font=\sffamily\small]
      at ($(current page.north west)+(4.4cm,-1.35cm)$) {\cursustitel};
    \node[anchor=north west,inner sep=0pt,text=white,font=\kopletter\fontsize{26}{31}\selectfont]
      at ($(current page.north west)+(4.4cm,-1.5cm)$) {#1};
  \end{tikzpicture}%
  \par\vspace*{1.05cm}}
\def\@makeschapterhead#1{%
  {\parindent\z@\RaggedRight\normalfont\vspace*{-6mm}%
   {\kopletter\fontsize{21}{25}\selectfont #1\par}
   \vspace{1mm}{\color{accent}\hrule height 1.4pt}\vspace{3mm}}}
\renewcommand{\section}{\@startsection{section}{1}{\z@}
  {-4ex \@plus -1ex \@minus -.4ex}{1.2ex \@plus .2ex}
  {\normalfont\kopletter\large}}
\DeclareRobustCommand{\secnummer}[1]{%
  \ifopdrachtdeel\def\secvul{fel}\def\sectekst{inkt}\else\def\secvul{accent}\def\sectekst{white}\fi
  \tikz[baseline=(c.base)]\node[fill=\secvul,rounded corners=2pt,text=\sectekst,inner xsep=1.2mm,inner ysep=.6mm,font=\sffamily\bfseries\small](c){#1};}
% Start van de Opdrachten: nieuwe pagina met een felle band
\newcommand{\opdrachtenkop}{%
  \clearpage\global\opdrachtdeeltrue\thispagestyle{plain}%
  \begin{tikzpicture}[remember picture,overlay]
    \fill[fel] (current page.north west) rectangle ($(current page.north east)+(0,-2.6cm)$);
    \node[anchor=west,inner sep=0pt,text=inkt,font=\kopletter\fontsize{26}{31}\selectfont]
      at ($(current page.north west)+(2.5cm,-1.45cm)$) {Opdrachten};
  \end{tikzpicture}%
  \par\vspace*{2mm}}
% Felle balk voor Koppel en Test jezelf
\newcommand{\opdrachtsectie}[1]{%
  \par\addvspace{3.5ex}\needspace{10\baselineskip}\noindent
  \tikz\node[fill=fel,rounded corners=2.5pt,inner xsep=3mm,inner ysep=1.5mm,text width=\dimexpr\linewidth-6mm\relax,align=left,text=inkt,font=\kopletter\large]{#1};%
  \par\nobreak\addvspace{.5ex}}
% Sprongenfiguur in een theoriekader; het argument staat boven de sprongen, bv. +v
\newcommand{\sprongfiguur}[1]{%
  \begin{tikzpicture}[x=.7cm,y=1cm,>={Stealth[length=1.8mm]},font=\small]
    \draw[lijngrijs] (-.2,0) -- (9.9,0);
    \foreach \x/\t in {0/u_1,1.6/u_2,3.2/u_3,7.8/u_{n-1},9.4/u_n}{
      \fill[accent] (\x,0) circle (2pt);
      \node[below=1pt] at (\x,0) {$\t$};}
    \node at (5.5,0) [fill=accent!9!white,inner sep=1pt] {$\cdots$};
    \foreach \a/\b in {0/1.6,1.6/3.2,7.8/9.4}{
      \draw[->,accent,thick] (\a,.08) to[bend left=45] node[above=-1pt,text=black,font=\scriptsize] {$#1$} (\b,.08);}
    \draw[decorate,decoration={brace,amplitude=4pt},grijs] (0,.62) -- (9.4,.62)
      node[midway,above=3pt,text=black] {$n-1$ sprongen};
  \end{tikzpicture}}
\renewcommand{\@seccntformat}[1]{\llap{\secnummer{\csname the#1\endcsname}\hspace{2mm}}}% nummer hangt in de linkermarge
\setcounter{secnumdepth}{1} % subsecties krijgen geen nummer
\renewcommand{\subsection}{\@startsection{subsection}{2}{\z@}
  {-3ex \@plus -.8ex \@minus -.3ex}{1ex \@plus .2ex}
  {\normalfont\kopletter\color{accent}\nobreak}}% \nobreak: \color maakt anders een breekpunt vlak na een sectiekop
\makeatother
\newcommand{\blokkop}[1]{\leavevmode{\kopletter\small\color{accent}#1}}
\pretocmd{\section}{\needspace{9\baselineskip}}{}{} % kop en kader niet los van de rest
% Niet meteen na een sectiekop: \needspace geeft daar een breekpunt, en dan blijft de sectiekop alleen onderaan staan
\makeatletter
\pretocmd{\subsection}{\if@nobreak\else\needspace{7\baselineskip}\fi}{}{}
\makeatother

\pagestyle{fancy}
\renewcommand{\chaptermark}[1]{\markboth{\thechapter.\ #1}{}}
\renewcommand{\sectionmark}[1]{}
\fancyhf{}
\fancyhead[L]{\kopletter\footnotesize\color{accent}\leftmark\ifopdrachtdeel\ \textperiodcentered\ Opdrachten\fi}
\fancyfoot[R]{\sffamily\footnotesize\color{grijs}\thepage}
\renewcommand{\headrulewidth}{0pt}
\fancypagestyle{plain}{\fancyhf{}\fancyfoot[R]{\sffamily\footnotesize\color{grijs}\thepage}\renewcommand{\headrulewidth}{0pt}}

% Tab aan de rand van elke pagina: kleur, nummer en plaats tonen het hoofdstuk.
% 14 mm breed, zodat er genoeg overblijft als een printer de rand niet bedrukt.
\newlength{\tabx}\newlength{\taby}
\AddToHook{shipout/background}{%
  \ifnum\value{chapter}>0
    \setlength{\tabx}{\paperwidth}\addtolength{\tabx}{-14mm}%
    \setlength{\taby}{\dimexpr -4.3cm - 3cm*\numexpr\value{chapter}-1\relax\relax}%
    \ifopdrachtdeel\def\tabvul{fel}\def\tabtekst{inkt}\else\def\tabvul{accent}\def\tabtekst{white}\fi
    \put(\tabx,\taby){\textcolor{\tabvul}{\rule[-2.7cm]{14mm}{2.7cm}}}%
    \put(\tabx,\taby){\raisebox{-1.6cm}[0pt][0pt]{\makebox[8.5mm]{\color{\tabtekst}\kopletter\LARGE\thechapter}}}%
  \fi}

%----------------------------------------------------------------------------------------
% BLOKKEN
%----------------------------------------------------------------------------------------
\tcbset{blok/.style={enhanced,sharp corners,boxrule=0pt,frame hidden,
    left=4mm,right=3mm,top=2mm,bottom=2mm,before skip=2ex,after skip=2ex}}

% Theorie: definities, eigenschappen, formules en hun bewijs. Volledig omkaderd en getint.
\newtcolorbox{theorie}{enhanced,sharp corners,boxrule=1pt,colframe=accent,colback=accent!9!white,
  left=4mm,right=3mm,top=2mm,bottom=2mm,before skip=2ex,after skip=2ex,
  before upper={\RaggedRight\setlength{\parskip}{3pt}}}
% Formule in een theoriekader: vet en in de kleur van het hoofdstuk. Gebruik in wiskundemodus: $\formule{...}$
\newcommand{\formule}[1]{\text{\boldmath\color{accent}$\displaystyle #1$}}
% Bewijs: staat in het theoriekader, onder de eigenschap. Binnenin meestal de omgeving stappen.
\newenvironment{bewijs}{\par\addvspace{1.5mm}\blokkop{Bewijs}\par\nobreak}{\par}
% Werkwijze: hoe je een soort opgave aanpakt. Enkel een balk links, geen tint, met een klein label.
\newtcolorbox{werkwijze}{blok,colback=white,
  borderline west={2.5pt}{0pt}{accent},
  before upper={\RaggedRight\setlength{\parskip}{3pt}\blokkop{Werkwijze}\quad}}

% Voorbeeld: altijd uitgewerkt, ook in het handboek. Oefeningen zijn dat niet.
\newcounter{voorbeeld}[chapter]
\providecommand{\theHvoorbeeld}{\thechapter.\arabic{voorbeeld}}
\newcommand{\voorbeeld}[3]{%
  \par\addvspace{2ex}\refstepcounter{voorbeeld}%
  \vbox{\noindent\tikz[baseline=(c.base)]\node[fill=accent,rounded corners=2pt,text=white,inner xsep=1.6mm,inner ysep=.6mm,font=\sffamily\bfseries\footnotesize](c){Voorbeeld \thevoorbeeld\ifstrempty{#1}{}{: #1}};\quad #2\par
  \uitwerkingblok{#3}}
  \addvspace{1.5ex}}

\newcommand{\uitwerkingblok}[1]{%
  \begin{tcolorbox}[blok,colback=white,borderline west={1pt}{0pt}{accent},before skip=.5ex,after skip=1ex,
      before upper=\RaggedRight\setlength{\parskip}{3pt}]
    #1
  \end{tcolorbox}}

\newenvironment{stappen}{%
  \par\renewcommand{\arraystretch}{1.25}%
  \noindent\begin{tabular}{@{}>{$\displaystyle}p{7cm}<{$}@{\hspace{4mm}}>{\RaggedRight\small\arraybackslash}p{7.6cm}@{}}%
}{\end{tabular}\par}

% Alleen voor de leerkracht
\newcommand{\lkr}[1]{\ifleerkracht
  \begin{tcolorbox}[blok,colback=black!5!white,borderline west={1pt}{0pt}{grijs},before skip=.6ex,after skip=.8ex,
      before upper=\RaggedRight\sffamily\footnotesize]
    {\bfseries Voor de leerkracht.} #1
  \end{tcolorbox}\fi}

% Alleen in de zelfstudieversie: uitleg in lopende tekst
\newcommand{\thuis}[1]{\ifzelfstudie\par #1\par\fi}

% Alleen in de zelfstudieversie: extra opgave, met het antwoord onderaan de pagina
\newcounter{nujij}[chapter]
\providecommand{\theHnujij}{\thechapter.\arabic{nujij}}
\makeatletter
\newcommand{\nujij}[2]{\ifzelfstudie
  \refstepcounter{nujij}%
  \begin{tcolorbox}[blok,colback=white,frame hidden,borderline={1.2pt}{0pt}{fel},left=2.5mm,right=2.5mm,
      before upper=\RaggedRight]
    \blokkop{Nu jij \thenujij}\quad #1
  \end{tcolorbox}%
  \begingroup\long\def\@makefntext##1{\noindent\footnotesize ##1}%
  \footnotetext{\RaggedRight\textbf{Nu jij \thenujij:} #2}\endgroup\fi}
\makeatother

\newcommand{\vak}{\mbox{\fboxsep=0pt\fboxrule=.5pt\fbox{\rule{0pt}{1.8ex}\hspace{1.4em}}}}

%----------------------------------------------------------------------------------------
% OEFENINGEN: een kolom, nummer links, balkjes voor de moeilijkheid (bij de Opdrachten)
%----------------------------------------------------------------------------------------
\newcounter{oefening}[chapter] % nummering herstart per hoofdstuk
\providecommand{\theHoefening}{\thechapter.\arabic{oefening}}
\newcommand{\oefstrook}[1]{%
  \needspace{10\baselineskip}\par\noindent
  \tikz[baseline=(s.base)]\node[fill=fel,rounded corners=2.5pt,inner xsep=2.2mm,inner ysep=1mm,text=inkt,font=\sffamily\bfseries\small](s){#1};%
  \hspace{1.5mm}\textcolor{fel}{\leaders\hrule height 2.9pt depth -1.2pt\hfill\kern0pt}\par\nobreak}
\newcommand{\oefnummer}[1]{%
  \tikz[baseline=(c.base)]\node[fill=fel,rounded corners=2pt,inner xsep=1.2mm,inner ysep=.95mm,minimum width=7mm,text=inkt,font=\sffamily\bfseries\small](c){#1};}
\newenvironment{oefeningen}[1][]{\par\addvspace{1.2ex}\ifstrempty{#1}{}{\addvspace{1ex}\oefstrook{#1}}}{\par\addvspace{1ex}}
\newenvironment{oefening}[1][1]{%
  \needspace{5\baselineskip}% opgave niet los van de deelvragen
  \refstepcounter{oefening}%
  \renewcommand{\Currentlabel}{\theoefening}%
  \begin{list}{}{\setlength{\leftmargin}{15mm}\setlength{\labelwidth}{15mm}\setlength{\labelsep}{0pt}%
      \setlength{\topsep}{1.7ex}\setlength{\partopsep}{0pt}\setlength{\parsep}{3pt}\setlength{\itemsep}{0pt}}%
  \interlinepenalty=10000 % een oefening breekt niet over twee pagina's
  \item[{\makebox[15mm][l]{\oefnummer{\theoefening}\ifopdrachtdeel\hspace{.9mm}\niveau{#1}\fi}}]\ignorespaces
}{\end{list}}

% Kleine prent rechts naast een opdracht: \begin{metprent}{naam} opgave en deelvragen \end{metprent}
% De prent staat in afbeeldingen/prenten/, de solution komt na \end{metprent}. Optioneel: [breedte].
\newenvironment{metprent}[2][2.1cm]{%
  \def\prentnaam{#2}\def\prentbreedte{#1}%
  \begin{minipage}[t]{\dimexpr\linewidth-#1-4mm\relax}%
}{%
  \end{minipage}\hfill
  \begin{minipage}[t]{\prentbreedte}
    \vspace{-3mm}\raggedleft
    \includegraphics[width=\linewidth,height=2.1cm,keepaspectratio]{prenten/\prentnaam}
  \end{minipage}\par}

% Moeilijkheid: drie oplopende balkjes, niveau 1, 2 of 3. Ook duidelijk in zwart-wit.
\newcommand{\niveau}[1]{%
  \begin{tikzpicture}[baseline=.1mm,x=1mm,y=1mm]
    \foreach \b/\h in {1/1.6,2/2.8,3/4}{
      \ifnum\b>#1\relax\def\balk{black!20}\else\def\balk{inkt}\fi
      \fill[\balk] ({(\b-1)*1.4},0) rectangle ++(1.05,\h);}
  \end{tikzpicture}}
% Onderwerp van een opdracht, enkel in de leerkrachtenversie
\newcommand{\thema}[1]{\ifleerkracht{\sffamily\scriptsize\color{grijs}[#1]}\ \fi\ignorespaces}
% Label naar andere leerstof: de kleur toont de rubriek (analyse, algebra, meetkunde, discreet, data)
\newcommand{\koppel}[2]{%
  \tikz[baseline=(k.base)]\node[fill=rubriek-#1,rounded corners=2mm,text=white,inner xsep=2mm,inner ysep=.7mm,font=\sffamily\bfseries\footnotesize](k){%
    \tikz[baseline=-.2mm,x=1mm,y=1mm,line width=.5pt,white]{\draw[rounded corners=.5mm] (0,0) rectangle (1.9,1.1);\draw[rounded corners=.5mm] (1.3,0) rectangle (3.2,1.1);}\ #2};\ \ignorespaces}
\newcommand{\voorbeeldpil}[1]{%
  \tikz[baseline=(c.base)]\node[fill=accent,rounded corners=2pt,text=white,inner xsep=1.6mm,inner ysep=.6mm,font=\sffamily\bfseries\footnotesize](c){#1};}
\newcommand{\oefstrookje}[1]{%
  \tikz[baseline=(s.base)]\node[fill=fel,rounded corners=2.5pt,inner xsep=2.2mm,inner ysep=1mm,text=inkt,font=\sffamily\bfseries\small](s){#1};}

% Schema 'Fout? Herhaal'. In het hoofdstuk: \herhaalschema{\fv{test}{voorbeeld}{oefening} \fzv{test}{oefening} ...}
% met labels als argument. Het komt ook in de lijsten met oplossingen, met de nummers al ingevuld.
\protected\def\foutrij#1#2#3{\oefnummer{#1} & $\rightarrow$ & \ifstrempty{#2}{}{\voorbeeldpil{Voorbeeld #2}\ }\oefnummer{#3}\\}
\protected\def\foutschema#1{%
  \par\addvspace{2ex}\needspace{8\baselineskip}\noindent\blokkop{Fout? Herhaal}\par\nobreak\vspace{.6ex}%
  {\renewcommand{\arraystretch}{1.5}\noindent\begin{tabular}{@{}l@{\hspace{2.5mm}}c@{\hspace{2.5mm}}l@{}}#1\end{tabular}\par}\addvspace{1.5ex}}
\newcommand{\herhaalschema}[1]{%
  \begingroup
  \def\fv##1##2##3{\foutrij{\getrefnumber{##1}}{\getrefnumber{##2}}{\getrefnumber{##3}}}%
  \def\fzv##1##2{\foutrij{\getrefnumber{##1}}{}{\getrefnumber{##2}}}%
  \edef\schemaregels{#1}%
  \foutschema{\schemaregels}%
  \Writetofile{oplossingen}{\foutschema{\schemaregels}}%
  \Writetofile{uitgewerkteoplossingen}{\foutschema{\schemaregels}}%
  \endgroup}

%----------------------------------------------------------------------------------------
% OPLOSSINGEN
%----------------------------------------------------------------------------------------
\Newassociation{solution}{Sol}{oplossingen}
\Newassociation{fullsolution}{Fullsol}{uitgewerkteoplossingen}
\renewenvironment{Sol}[1]{\par\noindent\textcolor{accent}{\textbf{#1.}}~\ignorespaces}{\par\vspace{.4ex}}
\renewenvironment{Fullsol}[1]{\par\medskip\noindent\textcolor{accent}{\textbf{Oefening #1}}\par\nobreak}{\par}
\newcommand{\oplossingenkop}[2]{%
  \colorlet{accent}{hd#1}\colorlet{fel}{hf#1}%
  \par\vspace{2mm}\noindent{\kopletter\color{accent}#1.\ \ #2}\par\nobreak\vspace{1mm}}
\newcommand{\mychapter}[2]{%
  \clearpage % eerst de vorige pagina afsluiten: die houdt de kleur van het vorige hoofdstuk
  \global\opdrachtdeelfalse
  \colorlet{accent}{hd\number\numexpr\value{chapter}+1\relax}%
  \colorlet{fel}{hf\number\numexpr\value{chapter}+1\relax}%
  \Writetofile{oplossingen}{\protect\oplossingenkop{#1}{#2}}%
  \Writetofile{uitgewerkteoplossingen}{\protect\oplossingenkop{#1}{#2}}%
  \chapter{#2}}

%----------------------------------------------------------------------------------------
% VOORKAFT
%----------------------------------------------------------------------------------------
%----------------------------------------------------------------------------------------
% VOORKANT: dezelfde indeling voor alle handboeken. Dit bestand wordt gekopieerd door
% Ontwerp/voorkanten/maak-voorkanten.py: pas het daar aan, niet in een cursus.
%
% Per cursus:  motief.tex             kunstwerk, bovenrand van de band en vormen
%              voorkant-gegevens.tex  rubriek, richtingen, graad, uren en titel
% Coördinaten in cm, met de oorsprong linksonder op de pagina (A4: 21 x 29,7).
%----------------------------------------------------------------------------------------
\providecommand{\omslagsamples}{120}

% Hoogte van de bovenrand op plaats #1, in macro #2. Handig om iets op de rand te zetten.
\newcommand{\randwaarde}[2]{\def\x{#1}\pgfmathsetmacro{#2}{\omslagrand}}
\newcommand{\randlijn}{%
  \draw[white,line width=2pt,line cap=round,line join=round]
    plot[domain=0:21,samples=\omslagsamples,variable=\x] (\x,{\omslagrand});}
\newcommand{\standaardband}{%
  \fill[accent] (0,0) -- plot[domain=0:21,samples=\omslagsamples,variable=\x] (\x,{\omslagrand}) -- (21,0) -- cycle;}

% Witte lijnen met een rand in de bandkleur, zodat ze ook op een licht schilderij zichtbaar blijven.
\tikzset{
  omslaglijn/.style={draw=white,line width=#1,preaction={draw=accent,line width=#1+1.8pt}},
  omslaglijn/.default=1pt,
  omslagstip/.style={fill=white,draw=accent,line width=.7pt},
  omslaglabel/.style={fill=accent,text=white,inner sep=2pt,rounded corners=1.5pt,font=\small},
}

\InputIfFileExists{motief}{}{}

\newcommand{\maakomslag}{%
  \begingroup\thispagestyle{empty}%
  \begin{tikzpicture}[remember picture,overlay]
  \begin{scope}[shift={(current page.south west)}]
    \node[anchor=north west,inner sep=0pt] at (0,29.7)
      {\includegraphics[width=\paperwidth]{\omslagbestand}};
    \ifcsname omslagband\endcsname\omslagband\else\standaardband\fi
    \omslagmotief
    % titel, met de rubriek erboven en de versie eronder
    \node[name=omslagtitelnode,anchor=south west,inner sep=0pt,text=white,align=left,
      font=\kopletter\fontsize{\omslagtitelgrootte}{\numexpr\omslagtitelgrootte*107/100\relax}\selectfont]
      at (2,5.2) {\strut\omslagtitel\strut};
    \node[anchor=south west,inner sep=0pt,text=fel,font=\sffamily\bfseries\normalsize]
      at ([yshift=2mm]omslagtitelnode.north west)
      {Wiskunde\ifdefempty{\cursusrubriek}{}{ \textbullet\ \cursusrubriek}};
    \node[anchor=north west,fill=fel,text=inkt,rounded corners=2.5pt,inner xsep=2.5mm,inner ysep=1.2mm,
      font=\kopletter\large] at (2,5.15) {\cursusondertitel};
    % de gegevens
    \node[anchor=south west,inner sep=0pt,text=white] at (2,1.45)
      {\sffamily\small\renewcommand{\arraystretch}{1.15}\hyphenpenalty=10000\exhyphenpenalty=10000%
       \begin{tabular}{@{}>{\bfseries}l@{\hspace{4mm}}>{\RaggedRight\arraybackslash}p{9.2cm}@{}}
         Richtingen & \cursusrichtingen \\
         Graad      & \cursusgraad \\
         Wiskunde   & \cursusuren
       \end{tabular}};
    \node[anchor=south west,inner sep=0pt,text=white,font=\sffamily\tiny] at (2,.75) {Omslag: \omslagbron};
    \node[anchor=south east,inner sep=0pt] at (19,1.5)
      {\includegraphics[width=4.6cm]{axioma-logo-wit.pdf}};
  \end{scope}
  \end{tikzpicture}%
  \vfill\endgroup\clearpage}


\usepackage[hyperfootnotes=false]{hyperref}
\hypersetup{hidelinks,pdftitle={\cursustitel},pdfauthor={\cursusauteur},pdflang={nl}}


\ifleerkracht\renewcommand{\cursusondertitel}{Leerkrachtenversie}\fi
\ifzelfstudie\renewcommand{\cursusondertitel}{Zelfstudie voor thuis}\fi

\providecommand{\uitvoer}{cursus}
\newif\ifcursus
\ifdefstring{\uitvoer}{cursus}{\cursustrue}{\cursusfalse}

\begin{document}

\ifcursus
\Opensolutionfile{oplossingen}
\Opensolutionfile{uitgewerkteoplossingen}
\fi

\ifcursus\else
\renewcommand{\cursusondertitel}{\ifdefstring{\uitvoer}{numeriek}{Numerieke oplossingen}{Uitgewerkte oplossingen}}
\fi
\maakomslag

\ifcursus

%----------------------------------------------------------------------------------------
% ZO WERK JE MET DIT BOEK: komt vooraan in elk handboek. De uitleg over hoe de onderdelen
% werken staat alleen hier, niet meer in de hoofdstukken.
%----------------------------------------------------------------------------------------
\chapter*{Zo werk je met dit boek}

\ifhandboek
Dit boek gebruik je in de les. Je vindt er de belangrijkste leerstof, uitgewerkte voorbeelden en vooral veel opgaven. De oplossing van een opgave noteer je zelf in je schrift.
\fi
\ifzelfstudie
Deze versie is voor thuis: om een les in te halen of iets rustig te herhalen. Ze bevat hetzelfde als het handboek, met extra uitleg, extra opgaven en de antwoorden. De oefeningen hebben dezelfde nummers als in het handboek.
\fi
\ifleerkracht
Dit is de leerkrachtenversie: het handboek met didactische tips en de antwoorden. De oefeningen hebben dezelfde nummers als in het handboek.
\fi

\section*{Kleuren}

Elk hoofdstuk heeft een eigen kleur en een nummer aan de rand van de pagina. De \textbf{donkere} tint gebruiken we voor uitleg, de \textbf{felle} tint voor opdrachten.

\section*{De onderdelen}

{\renewcommand{\arraystretch}{1.75}
\begin{tabular}{@{}>{\RaggedRight\arraybackslash}p{3.3cm}>{\RaggedRight\arraybackslash}p{12.3cm}@{}}
\tikz[baseline=(k.base)]\node[fill=accent!9!white,draw=accent,line width=1pt,inner xsep=2.2mm,inner ysep=1mm,font=\sffamily\bfseries\small,text=accent](k){Theorie}; & Het omkaderde blok: definities, eigenschappen en formules, met hun bewijs. Dit moet je kennen, en hier zoek je het terug. \\
\tikz[baseline=(k.base)]{\node[inner xsep=2.2mm,inner ysep=1mm,font=\sffamily\bfseries\small,text=accent](k){Werkwijze};\draw[accent,line width=2.5pt] (k.north west) -- (k.south west);} & Het blok met een balk links: de stappen om een bepaald soort opgave aan te pakken. \\
\voorbeeldpil{Voorbeeld 1} & Een uitgewerkt voorbeeld: zo pak je dit soort opgave aan. \\
\oefstrookje{Start} & \textbf{Opwarmer:} leerstof die je nodig hebt, uit het hoofd. \textbf{Verken:} probeer dit zelf vóór de uitleg; fouten horen erbij. \\
\oefnummer{8} & Een oefening. Na elk stuk leerstof staat een korte oefening om het meteen te proberen. \\
\oefstrookje{Opdrachten} & Het grootste deel van de opgaven, na alle leerstof van het hoofdstuk. Ze staan bewust door elkaar: welke methode je nodig hebt, moet je zelf herkennen. Dit deel start op een nieuwe pagina, en vanaf daar is de tab aan de rand fel. \\
\niveau{1}\ \niveau{2}\ \niveau{3} & De moeilijkheid van een opdracht: basis, gevorderd of uitdagend. \\
\oefstrookje{Koppel} & Een kleine reeks die dit hoofdstuk verbindt met andere leerstof. De kleur van het label toont de rubriek: \koppel{analyse}{Analyse}\koppel{algebra}{Algebra}\koppel{meetkunde}{Meetkunde}\koppel{discreet}{Discrete wiskunde}\koppel{data}{Data en onzekerheid} \\
\oefstrookje{Test jezelf} & Zonder boek of notities. Verbeter daarna met de oplossingen. \\
\blokkop{Fout? Herhaal} & Dit schema staat bij de test en bij de antwoorden. Het toont wat je moet herhalen bij een fout: \oefnummer{45} $\rightarrow$ \voorbeeldpil{Voorbeeld 1} \oefnummer{9}
\end{tabular}}

\vfill
{\footnotesize\color{grijs} De kleine prenten bij de opdrachten komen uit het publiek domein, via Wikimedia Commons: tekeningen van Pearson Scott Foresman, een gravure uit \textit{Popular Science Monthly} en een houtsnede uit de 16de eeuw.\par}

\clearpage % voor de teller van het hoofdstuk verandert: deze pagina krijgt geen tab

%----------------------------------------------------------------------------------------
% RIJEN: BEGRIPPEN. Alle opgaven zijn eigen werk.
%
% Opbouw: Start, leerstof (kern, uitgewerkt voorbeeld, korte oefening), en daarna op een
% nieuwe pagina de Opdrachten (door elkaar, stijgende moeilijkheid), Koppel en Test jezelf.
% Labels van dit hoofdstuk beginnen met h1-.
%----------------------------------------------------------------------------------------

% Kleine grafiek van een rij: {naam}{ymin}{ymax}{ticks op de y-as}{punten n/u_n}
\newcommand{\minigrafiek}[5]{%
  \pgfmathsetmacro{\eenheid}{2.3/(#3-(#2))}%
  \begin{tikzpicture}[x=.38cm,y=\eenheid cm,>={Stealth[length=1.4mm]},font=\scriptsize]
    \draw[->,lijngrijs] (-.2,0) -- (6.8,0) node[right,text=black] {$n$};
    \draw[->,lijngrijs] (0,#2) -- (0,#3) node[above,text=black] {$u_n$};
    \foreach \x in {1,...,6} \draw[lijngrijs] (\x,0) ++(0,1.2pt) -- ++(0,-2.4pt) node[below=-1pt,text=grijs] {\x};
    \foreach \y in {#4} \draw[lijngrijs] (0,\y) ++(1.2pt,0) -- ++(-2.4pt,0) node[left=-1pt,text=grijs] {$\y$};
    \foreach \x/\y in {#5} \fill[accent] (\x,\y) circle (1.6pt);
    \node[anchor=south,font=\sffamily\bfseries\footnotesize] at (3.4,#3) {#1};
  \end{tikzpicture}}

\mychapter{1}{Rijen: begrippen}

\lkr{Leerplan B+S'': LPD 44 (recursief en expliciet voorschrift, rij als discrete functie). Dit hoofdstuk legt de begrippen vast voor hoofdstuk 2; rekenkundige en meetkundige rijen krijgen hier nog geen naam. Drie rode draden: het verschil tussen volgnummer en term, het verschil tussen expliciet en recursief, en de grafiek met losse punten.}

\begin{oefeningen}[Start]
\begin{oefening}
\blokkop{Opwarmer}
\begin{enumerate}
  \item Gegeven: $f(x) = x^2 - 3x$. Bereken $f(4)$ en $f(-1)$.
  \item Bereken $(-1)^5$, $(-1)^{12}$, $-2^4$ en $(-2)^4$.
  \item Werk uit en vereenvoudig: $(n+1)^2 - n^2$.
  \item Los op: $n^2 + n - 12 = 0$.
  \item Bepaal het domein van $f(x) = \dfrac{3}{x-2}$.
\end{enumerate}
\begin{solution}
\textbf{a.}~4 en 4 \quad \textbf{b.}~$-1$; 1; $-16$; 16 \quad \textbf{c.}~$2n+1$ \quad \textbf{d.}~$n=3$ of $n=-4$ \quad \textbf{e.}~$\mathbb{R}\setminus\{2\}$
\end{solution}
\begin{fullsolution}
\textbf{a.}~$f(4)=16-12=4$ en $f(-1)=1+3=4$.\\
\textbf{b.}~Oneven exponent: $(-1)^5=-1$. Even exponent: $(-1)^{12}=1$. Zonder haakjes hoort de macht alleen bij 2: $-2^4=-16$. Met haakjes: $(-2)^4=16$.\\
\textbf{c.}~$n^2+2n+1-n^2=2n+1$.\\
\textbf{d.}~$(n-3)(n+4)=0$, dus $n=3$ of $n=-4$.\\
\textbf{e.}~De noemer mag niet nul zijn, dus $x\neq 2$.
\end{fullsolution}
\end{oefening}

\begin{oefening}\label{oef:h1-verken}
\blokkop{Verken}\quad Met blokjes bouw je trappen. Trap 1 telt 1 blokje, trap 2 telt 3 blokjes, trap 3 telt 6 blokjes.

\begin{tikzpicture}[x=3.1mm,y=3.1mm]
  \foreach \n/\s in {1/0,2/3,3/7,4/12}{
    \foreach \k in {1,...,\n}{
      \foreach \h in {1,...,\k}{
        \draw[accent,fill=accent!15] (\s+\k-1,\h-1) rectangle ++(1,1);}}
    \node[below,font=\scriptsize\sffamily] at (\s+\n/2,0) {trap \n};}
\end{tikzpicture}
\begin{enumerate}
  \item Hoeveel blokjes telt trap 4? En trap 5?
  \item Trap 9 telt 45 blokjes. Hoeveel blokjes telt trap 10? Hoe zie je dat zonder te tekenen?
  \item Leg twee gelijke trappen tegen elkaar, één ervan omgekeerd. Welke rechthoek krijg je met twee keer trap 4? Hoeveel blokjes telt trap 100 dan?
\end{enumerate}
\begin{solution}
\textbf{a.}~10 en 15 \quad \textbf{b.}~55 \quad \textbf{c.}~4 op 5; 5050
\end{solution}
\begin{fullsolution}
\textbf{a.}~Er komt telkens een kolom bij: $6+4=10$ en $10+5=15$.\\
\textbf{b.}~Trap 10 is trap 9 plus een kolom van 10 blokjes: $45+10=55$.\\
\textbf{c.}~Twee keer trap 4 geeft een rechthoek van 4 op 5, dus trap 4 telt $\dfrac{4\cdot 5}{2}=10$ blokjes. Zo telt trap 100 er $\dfrac{100\cdot 101}{2}=5050$.
\end{fullsolution}
\end{oefening}
\end{oefeningen}

\lkr{Laat oefening \ref{oef:h1-verken} maken vóór de uitleg. In b redeneren leerlingen recursief (de vorige trap plus een kolom), in c expliciet (rechtstreeks uit het nummer van de trap). Kom hierop terug bij de theorie \emph{Expliciet of recursief?}.}

%========================================================================================
\section{Termen en volgnummers}

\begin{theorie}
Een \emph{rij} is een opeenvolging van getallen in een vaste volgorde. De getallen zijn de \emph{termen} van de rij. De index is het \emph{volgnummer}: $u_1$ is de eerste term, $u_2$ de tweede, en $u_n$ de $n$-de term of \emph{algemene term}, met $n \in \mathbb{N}_0$.

\vspace{1mm}
{\renewcommand{\arraystretch}{1.25}%
\begin{tabular}{@{}l*{7}{>{\centering\arraybackslash}p{11mm}}@{}}
volgnummer & 1 & 2 & 3 & $\cdots$ & $n-1$ & $n$ & $n+1$ \\
term & $u_1$ & $u_2$ & $u_3$ & $\cdots$ & $u_{n-1}$ & $u_n$ & $u_{n+1}$
\end{tabular}}

\vspace{1mm}
Let op: $u_{n-1}$ is de term vóór $u_n$. Dat is iets anders dan $u_n - 1$.
\end{theorie}

\thuis{Het volgnummer zegt \emph{waar} een term staat, de term zelf zegt \emph{welk getal} daar staat. In de rij $7,\ 2,\ 9,\ \ldots$ is $u_3 = 9$: op plaats 3 staat het getal 9. Drie puntjes betekenen dat de rij verdergaat.}

\voorbeeld{}{Het aantal dagen per maand in een gewoon jaar geeft een rij:\\[.5ex]
$31,\ 28,\ 31,\ 30,\ 31,\ 30,\ 31,\ 31,\ 30,\ 31,\ 30,\ 31$\\[.5ex]
Geef $u_2$ en $u_7$. Voor welke volgnummers is $u_n = 30$?\label{vb:h1-maanden}}{%
\begin{stappen}
u_2 = 28 & De tweede term: februari. \\
u_7 = 31 & De zevende term: juli. \\
n = 4,\ 6,\ 9 \text{ of } 11 & Eenzelfde getal mag meer dan eens voorkomen. Bij een volgnummer hoort wel maar één term.
\end{stappen}}

\lkr{Deze rij heeft twaalf termen en geen handig voorschrift. Zo zien leerlingen meteen dat een rij meer is dan een formule.}

\nujij{In de rij $4,\ 7,\ 7,\ 1,\ 0,\ 7$: wat is $u_3$? Voor welke volgnummers is $u_n = 7$?}{$u_3 = 7$; $n = 2$, 3 of 6.}

\begin{oefeningen}
\begin{oefening}
Gegeven is de rij $8,\ 3,\ 8,\ 0,\ 5,\ 3,\ 1$.
\begin{enumerate}
  \item Wat is $u_4$?
  \item Voor welke volgnummers is $u_n = 3$?
  \item Neem $n = 5$. Wat zijn $u_{n-1}$ en $u_{n+1}$?
\end{enumerate}
\begin{solution}
\textbf{a.}~0 \quad \textbf{b.}~$n=2$ of $n=6$ \quad \textbf{c.}~0 en 3
\end{solution}
\begin{fullsolution}
\textbf{a.}~De vierde term is 0.\\
\textbf{b.}~Het getal 3 staat op plaats 2 en plaats 6.\\
\textbf{c.}~$u_{n-1}=u_4=0$ en $u_{n+1}=u_6=3$.
\end{fullsolution}
\end{oefening}
\end{oefeningen}

%========================================================================================
\section{Expliciet voorschrift}

\subsection{Termen berekenen}

\begin{theorie}
Een \emph{expliciet voorschrift} geeft $u_n$ rechtstreeks uit het volgnummer $n$. Je kunt elke term berekenen zonder de vorige termen te kennen.
\end{theorie}

\voorbeeld{}{Gegeven: $u_n = n^2 - 5n + 7$. Bereken $u_1$, $u_2$, $u_3$ en $u_{10}$. Schrijf ook $u_{n+1}$ zo eenvoudig mogelijk.\label{vb:h1-termen}}{%
\begin{stappen}
u_1 = 1 - 5 + 7 = 3 & Vervang $n$ door 1. \\
u_2 = 4 - 10 + 7 = 1, \quad u_3 = 9 - 15 + 7 = 1 & Twee termen kunnen gelijk zijn. \\
u_{10} = 100 - 50 + 7 = 57 & De vorige termen heb je niet nodig. \\
u_{n+1} = (n+1)^2 - 5(n+1) + 7 & Vervang $n$ door $n+1$, met haakjes. \\
\phantom{u_{n+1}} = n^2 - 3n + 3 & Dit is de volgende term, en dus niet $u_n + 1$.
\end{stappen}}

\lkr{Laat leerlingen $u_{n+1}$ en $u_n + 1$ naast elkaar uitrekenen voor $n = 2$: 1 en 2. De indexnotatie is voor veel leerlingen het eerste struikelblok.}

\nujij{Gegeven: $u_n = 3n - n^2$. Bereken $u_1$, $u_4$ en $u_{n-1}$.}{$u_1 = 2$; $u_4 = -4$; $u_{n-1} = 3(n-1) - (n-1)^2 = -n^2 + 5n - 4$.}

\begin{oefeningen}
\begin{oefening}
Gegeven: $u_n = 2n^2 - n$. Bereken $u_1$, $u_2$, $u_3$ en $u_{20}$. Schrijf ook $u_{n-1}$ zo eenvoudig mogelijk.
\begin{solution}
1; 6; 15; 780; $u_{n-1}=2n^2-5n+3$
\end{solution}
\begin{fullsolution}
$u_1=2-1=1$, $u_2=8-2=6$, $u_3=18-3=15$, $u_{20}=800-20=780$.\\
$u_{n-1}=2(n-1)^2-(n-1)=2n^2-4n+2-n+1=2n^2-5n+3$.
\end{fullsolution}
\end{oefening}
\end{oefeningen}

\subsection{Voorschrift opstellen}

\begin{werkwijze}
Zoek het verband tussen volgnummer en term.
\begin{itemize}
  \item Vergelijk de termen met $n$, met $n^2$ en met machten zoals $2^n$. Tel er zo nodig iets bij of trek er iets van af.
  \item Bij een breuk bekijk je teller en noemer apart.
  \item Wisselend teken: $(-1)^n$ geeft $-1,\ 1,\ {-1},\ 1,\ \ldots$ en $(-1)^{n+1}$ geeft $1,\ {-1},\ 1,\ {-1},\ \ldots$
  \item Controleer je voorschrift met een term die je nog niet gebruikt hebt.
\end{itemize}
\end{werkwijze}

\voorbeeld{}{Bepaal een expliciet voorschrift van $\tfrac13,\ {-\tfrac25},\ \tfrac37,\ {-\tfrac49},\ \ldots$\label{vb:h1-opstellen}}{%
\begin{stappen}
\text{teken: } (-1)^{n+1} & Positief bij oneven $n$, negatief bij even $n$. \\
\text{teller: } 1,\ 2,\ 3,\ 4 \ \Rightarrow \ n & De teller is het volgnummer. \\
\text{noemer: } 3,\ 5,\ 7,\ 9 \ \Rightarrow \ 2n+1 & Telkens 2 meer: het dubbele van $n$, plus 1. \\[1ex]
u_n = (-1)^{n+1} \cdot \dfrac{n}{2n+1} & Controle: $u_4 = (-1)^5 \cdot \dfrac{4}{9} = -\dfrac{4}{9}$.
\end{stappen}}

\thuis{\blokkop{Half uitgewerkt}\quad Rij $3,\ 6,\ 11,\ 18,\ \ldots$ Vergelijk met de kwadraten $1,\ 4,\ 9,\ 16$. Schrijf over in je schrift en vul aan: \ elke term is $\vak$ meer dan $n^2$, dus $u_n = \vak$ en $u_5 = \vak$. \quad {\footnotesize(Antwoord: 2 meer, $u_n = n^2 + 2$ en $u_5 = 27$.)}}

\nujij{Bepaal een expliciet voorschrift van $1,\ {-\tfrac12},\ \tfrac13,\ {-\tfrac14},\ \ldots$}{$u_n = (-1)^{n+1} \cdot \tfrac1n$.}

\begin{oefeningen}
\begin{oefening}
Bepaal een expliciet voorschrift.
\begin{reeks}
  \item $6,\ 11,\ 16,\ 21,\ \ldots$
  \item $2,\ 5,\ 10,\ 17,\ \ldots$
  \item $-\tfrac12,\ \tfrac23,\ {-\tfrac34},\ \tfrac45,\ \ldots$
\end{reeks}
\begin{solution}
\textbf{a.}~$u_n=5n+1$ \quad \textbf{b.}~$u_n=n^2+1$ \quad \textbf{c.}~$u_n=(-1)^n\cdot\dfrac{n}{n+1}$
\end{solution}
\begin{fullsolution}
\textbf{a.}~Telkens 5 meer, dus vergelijk met $5n$: $5,\ 10,\ 15$. Elke term is 1 meer: $u_n=5n+1$.\\
\textbf{b.}~Elke term is 1 meer dan een kwadraat: $u_n=n^2+1$.\\
\textbf{c.}~Negatief bij oneven $n$, teller $n$, noemer $n+1$: $u_n=(-1)^n\cdot\dfrac{n}{n+1}$.
\end{fullsolution}
\end{oefening}
\end{oefeningen}

\subsection{Term of geen term}

\begin{werkwijze}
Is $x$ een term van de rij? Los $u_n = x$ op. Is er een oplossing $n \in \mathbb{N}_0$, dan is $x$ de $n$-de term. Anders is $x$ geen term.
\end{werkwijze}

\voorbeeld{}{Gegeven: $u_n = n^2 + 2n$. Is 48 een term van de rij? En 50?\label{vb:h1-isterm}}{%
\begin{stappen}
n^2 + 2n - 48 = 0 \ \Rightarrow \ (n-6)(n+8) = 0 & $n = 6$ of $n = -8$. Enkel $n = 6$ is een volgnummer: \textbf{ja}, 48 is de 6de term. \\
n^2 + 2n - 50 = 0 \ \Rightarrow \ n = -1 \pm \sqrt{51} & $n \approx 6{,}14$ of $n \approx -8{,}14$. Geen volgnummer: \textbf{nee}.
\end{stappen}}

\nujij{Gegeven: $u_n = n^2 + n$. Is 90 een term van de rij? En 100?}{90 is de 9de term ($n = 9$ of $n = -10$). Bij 100 is $n \approx 9{,}51$, dus geen term.}

\begin{oefeningen}
\begin{oefening}
Gegeven: $u_n = n^2 - 4n$. Is 21 een term van de rij? En 20?
\begin{solution}
21 is de 7de term; 20 is geen term.
\end{solution}
\begin{fullsolution}
$n^2-4n-21=0$ geeft $(n-7)(n+3)=0$, dus $n=7$ (en $n=-3$ valt weg).\\
$n^2-4n-20=0$ geeft $n=2\pm\sqrt{24}$, dus $n\approx 6{,}90$ of $n\approx -2{,}90$: geen volgnummer.
\end{fullsolution}
\end{oefening}
\end{oefeningen}

\subsection{Starten bij volgnummer 0}

Bij iets dat in de tijd verandert, start je vaak bij $u_0$: de beginwaarde, voor er een stap gebeurd is. Dan is $u_n$ de waarde na $n$ stappen, en is $u_1$ de \emph{tweede} term.

\voorbeeld{}{Een lift stijgt: de hoogte in meter is elke seconde $12,\ 15,\ 18,\ \ldots$ Geef een expliciet voorschrift als de eerste term $u_1$ heet, en als hij $u_0$ heet. Welk volgnummer heeft de hoogte 30 in elk geval?\label{vb:h1-nul}}{%
\begin{stappen}
u_1 = 12: \quad u_n = 3n + 9 & Controle: $3 \cdot 1 + 9 = 12$. \\
u_0 = 12: \quad u_n = 3n + 12 & Controle: $3 \cdot 0 + 12 = 12$. \\
3n + 9 = 30 \ \Rightarrow \ n = 7 & Bij start $u_1$ is 30 de term $u_7$. \\
3n + 12 = 30 \ \Rightarrow \ n = 6 & Bij start $u_0$ is 30 de term $u_6$: de hoogte na 6 seconden.
\end{stappen}}

\lkr{Hetzelfde getal, een ander volgnummer. In hoofdstuk 2 komt dit terug bij lineaire en exponentiële groei.}

\begin{oefeningen}
\begin{oefening}
Een vijver telt bij het begin 40 kikkers. $u_n$ is het aantal kikkers na $n$ jaar. Na een jaar zijn er 52, na twee jaar 61.
\begin{enumerate}
  \item Wat is $u_0$? Wat is $u_2$?
  \item Hoeveel termen tel je van $u_0$ tot en met $u_{10}$?
\end{enumerate}
\begin{solution}
\textbf{a.}~$u_0=40$ en $u_2=61$ \quad \textbf{b.}~11
\end{solution}
\begin{fullsolution}
\textbf{a.}~$u_0$ is het begin: 40. $u_2$ is het aantal na twee jaar: 61.\\
\textbf{b.}~De termen $u_1$ tot en met $u_{10}$ zijn er 10, en $u_0$ komt erbij: 11.
\end{fullsolution}
\end{oefening}
\end{oefeningen}

%========================================================================================
\section{Recursief voorschrift}

\subsection{Termen berekenen}

\begin{theorie}
Een \emph{recursief voorschrift} geeft $u_n$ uit een of meer vorige termen. Je hebt ook de \emph{begintermen} nodig: zonder die ligt de rij niet vast.
\end{theorie}

Bijvoorbeeld: $u_1 = 3$ en $u_n = 2u_{n-1} - 1$ voor $n \geqslant 2$. Hetzelfde voorschrift, met de index een plaats verschoven: $u_{n+1} = 2u_n - 1$ voor $n \geqslant 1$.

\thuis{\emph{Recursief} komt van het Latijnse \emph{recurrere}, teruglopen: je grijpt telkens terug naar de vorige term. Lees $u_n = 2u_{n-1} - 1$ dus als: een term is het dubbele van de vorige term, min 1.}

\voorbeeld{}{Gegeven: $u_1 = 3$ en $u_n = 2u_{n-1} - 1$ voor $n \geqslant 2$. Bereken $u_2$, $u_3$ en $u_4$.\label{vb:h1-rec}}{%
\begin{stappen}
u_2 = 2u_1 - 1 = 2 \cdot 3 - 1 = 5 & Vervang $n$ door 2: de vorige term is $u_1$. \\
u_3 = 2u_2 - 1 = 2 \cdot 5 - 1 = 9 & Elke stap gebruikt de uitkomst van de vorige stap. \\
u_4 = 2u_3 - 1 = 2 \cdot 9 - 1 = 17 & Voor $u_{100}$ moet je zo 99 stappen zetten.
\end{stappen}}

\nujij{Gegeven: $u_1 = 10$ en $u_n = u_{n-1} - 2n$ voor $n \geqslant 2$. Bereken $u_2$, $u_3$ en $u_4$.}{$u_2 = 6$; $u_3 = 0$; $u_4 = -8$.}

\begin{oefeningen}
\begin{oefening}
Gegeven: $u_1 = 1$ en $u_n = 3u_{n-1} + 2$ voor $n \geqslant 2$. Bereken $u_2$, $u_3$ en $u_4$.
\begin{solution}
5; 17; 53
\end{solution}
\begin{fullsolution}
$u_2=3\cdot 1+2=5$, $u_3=3\cdot 5+2=17$, $u_4=3\cdot 17+2=53$.
\end{fullsolution}
\end{oefening}
\end{oefeningen}

\subsection{Twee begintermen}

Gebruikt het voorschrift twee vorige termen, dan heb je twee begintermen nodig. De bekendste is de \emph{rij van Fibonacci}: $u_1 = u_2 = 1$ en $u_n = u_{n-1} + u_{n-2}$ voor $n \geqslant 3$. Die geeft $1,\ 1,\ 2,\ 3,\ 5,\ 8,\ 13,\ \ldots$

\voorbeeld{}{Gegeven: $u_1 = 2$, $u_2 = 3$ en $u_n = u_{n-1} \cdot u_{n-2} - 1$ voor $n \geqslant 3$. Bereken $u_3$, $u_4$ en $u_5$.\label{vb:h1-twee}}{%
\begin{stappen}
u_3 = u_2 \cdot u_1 - 1 = 3 \cdot 2 - 1 = 5 & Bij $n = 3$ zijn de twee vorige termen $u_2$ en $u_1$. \\
u_4 = u_3 \cdot u_2 - 1 = 5 \cdot 3 - 1 = 14 & \\
u_5 = u_4 \cdot u_3 - 1 = 14 \cdot 5 - 1 = 69 &
\end{stappen}}

\begin{oefeningen}
\begin{oefening}
Bepaal de eerste vijf termen.
\begin{enumerate}
  \item $u_1 = 5$, $u_2 = 2$ en $u_n = u_{n-1} + u_{n-2}$ voor $n \geqslant 3$
  \item $u_1 = 1$ en $u_n = u_{n-1} + 2n$ voor $n \geqslant 2$
\end{enumerate}
\begin{solution}
\textbf{a.}~5, 2, 7, 9, 16 \quad \textbf{b.}~1, 5, 11, 19, 29
\end{solution}
\begin{fullsolution}
\textbf{a.}~$u_3=2+5=7$, $u_4=7+2=9$, $u_5=9+7=16$.\\
\textbf{b.}~$u_2=1+4=5$, $u_3=5+6=11$, $u_4=11+8=19$, $u_5=19+10=29$. Het volgnummer staat hier mee in het voorschrift.
\end{fullsolution}
\end{oefening}
\end{oefeningen}

\subsection{Voorschrift opstellen}

\begin{werkwijze}
Kijk hoe je van een term naar de volgende gaat: wat komt erbij, waarmee vermenigvuldig je? Hangt de stap af van het volgnummer, schrijf ze dan met $n$ en controleer ze bij de eerste stap. Vergeet de begintermen niet.
\end{werkwijze}

\voorbeeld{}{Bepaal een recursief voorschrift van $1,\ 2,\ 4,\ 7,\ 11,\ 16,\ \ldots$\label{vb:h1-recopst}}{%
\begin{stappen}
\text{erbij: } 1,\ 2,\ 3,\ 4,\ 5 & De stap wordt telkens 1 groter. \\
u_2 - u_1 = 1,\quad u_3 - u_2 = 2 & Van $u_{n-1}$ naar $u_n$ komt er dus $n - 1$ bij, en niet $n$. \\
u_n = u_{n-1} + (n-1),\ u_1 = 1 & Geldig voor $n \geqslant 2$. Controle: $u_5 = 7 + 4 = 11$.
\end{stappen}}

\lkr{De fout $u_n = u_{n-1} + n$ is hier heel gewoon. Laat leerlingen die fout zelf ontdekken door de eerste stap te controleren: dan is $u_2 = 3$.}

\begin{oefeningen}
\begin{oefening}
Bepaal een recursief voorschrift.
\begin{reeks}
  \item $3,\ 6,\ 12,\ 24,\ \ldots$
  \item $2,\ 3,\ 5,\ 8,\ 12,\ \ldots$
  \item $1,\ 4,\ 13,\ 40,\ \ldots$
\end{reeks}
\begin{solution}
\textbf{a.}~$u_n=2u_{n-1}$, $u_1=3$ \quad \textbf{b.}~$u_n=u_{n-1}+(n-1)$, $u_1=2$ \quad \textbf{c.}~$u_n=3u_{n-1}+1$, $u_1=1$
\end{solution}
\begin{fullsolution}
\textbf{a.}~Telkens maal 2.\\
\textbf{b.}~Er komt 1, 2, 3, 4 bij: bij de stap naar $u_n$ is dat $n-1$.\\
\textbf{c.}~Probeer een vermenigvuldiging: $3\cdot 1+1=4$, $3\cdot 4+1=13$, $3\cdot 13+1=40$.
\end{fullsolution}
\end{oefening}
\end{oefeningen}

\subsection{Expliciet of recursief?}

\begin{tabular}{@{}>{\RaggedRight\arraybackslash}p{7cm}>{\RaggedRight\arraybackslash}p{7.6cm}@{}}
\textbf{Expliciet:} rechtstreeks uit het volgnummer. Handig om meteen een verre term te berekenen. & \textbf{Recursief:} stap per stap uit de vorige termen. Handig om te beschrijven hoe een rij ontstaat. \\[2mm]
\begin{tikzpicture}[>={Stealth[length=1.8mm]},font=\small]
  \node (a) at (0,0) {$100$};
  \node (b) at (2.4,0) {$u_{100}$};
  \draw[->,accent,thick] (a) -- (b);
\end{tikzpicture} &
\begin{tikzpicture}[>={Stealth[length=1.8mm]},font=\small]
  \node (a) at (0,0) {$u_1$};
  \node (b) at (1.1,0) {$u_2$};
  \node (c) at (2.2,0) {$u_3$};
  \node (d) at (3.3,0) {$\cdots$};
  \node (e) at (4.5,0) {$u_{100}$};
  \draw[->,accent,thick] (a) -- (b);
  \draw[->,accent,thick] (b) -- (c);
  \draw[->,accent,thick] (c) -- (d);
  \draw[->,accent,thick] (d) -- (e);
\end{tikzpicture}
\end{tabular}

\vspace{1mm}
Niet elke rij heeft een eenvoudig voorschrift. Voor de priemgetallen $2,\ 3,\ 5,\ 7,\ 11,\ \ldots$ kent niemand een eenvoudige formule.

\voorbeeld{}{Gegeven: $u_n = n^2$. Bepaal een recursief voorschrift.\label{vb:h1-omzet}}{%
\begin{stappen}
u_{n-1} = (n-1)^2 = n^2 - 2n + 1 & Vervang $n$ door $n - 1$. \\
u_n - u_{n-1} = 2n - 1 & Bereken wat erbij komt. \\
u_n = u_{n-1} + 2n - 1,\ u_1 = 1 & Geldig voor $n \geqslant 2$. Controle: $u_2 = 1 + 3 = 4$ en $u_3 = 4 + 5 = 9$.
\end{stappen}}

\lkr{Kom terug op oefening \ref{oef:h1-verken}: welke redenering was recursief en welke expliciet? Het aantal blokjes van trap $n$ is $u_n = u_{n-1} + n$ met $u_1 = 1$, of $u_n = \tfrac{n(n+1)}{2}$.}

\begin{oefeningen}
\begin{oefening}
Gegeven: $u_n = n^2 + 3n$. Bepaal een recursief voorschrift.
\begin{solution}
$u_n=u_{n-1}+2n+2$, $u_1=4$
\end{solution}
\begin{fullsolution}
$u_{n-1}=(n-1)^2+3(n-1)=n^2+n-2$, dus $u_n-u_{n-1}=2n+2$. Verder is $u_1=4$. Controle: $u_2=4+6=10=2^2+3\cdot 2$.
\end{fullsolution}
\end{oefening}
\end{oefeningen}

\thuis{\blokkop{Met GeoGebra (thuis)}\quad Controleer je rekenwerk met twee commando's.
\begin{itemize}
  \item Expliciet: \texttt{Rij(3n+1,n,1,10)} geeft $u_1$ tot en met $u_{10}$ van $u_n = 3n + 1$.
  \item Recursief: \texttt{IteratieLijst(2x+4,-2,9)} geeft $u_1 = -2$ en daarna nog 9 termen van $u_n = 2u_{n-1} + 4$. De $x$ staat voor de vorige term.
\end{itemize}}

%========================================================================================
\section{De grafiek van een rij}

\begin{theorie}
\begin{minipage}[t]{.6\linewidth}
Bij elk volgnummer hoort precies één term: een rij is een \emph{functie met domein} $\mathbb{N}_0$.

\vspace{1.5mm}
In de grafiek staan de volgnummers horizontaal en de termen verticaal. De grafiek bestaat uit de losse punten $(n, u_n)$. Verbind ze niet: $u_{2{,}5}$ bestaat niet.

\vspace{1.5mm}
De figuur toont $u_n = n^2 - 6n + 10$. De stippellijn is de functie $f(x) = x^2 - 6x + 10$, met domein $\mathbb{R}$. Die hoort niet bij de rij.
\end{minipage}\hfill
\begin{minipage}[t]{.36\linewidth}
\vspace{-2mm}
\begin{tikzpicture}[x=.62cm,y=.26cm,>={Stealth[length=1.6mm]},font=\scriptsize]
  \draw[->,lijngrijs] (-.1,0) -- (6.9,0) node[right,text=black] {$n$};
  \draw[->,lijngrijs] (0,-.3) -- (0,12) node[right,text=black] {$u_n$};
  \foreach \x in {1,...,6} \draw[lijngrijs] (\x,.15) -- (\x,-.15) node[below=-1pt,text=grijs] {\x};
  \foreach \y in {5,10} \draw[lijngrijs] (.07,\y) -- (-.07,\y) node[left=-1pt,text=grijs] {\y};
  \draw[lijngrijs,dashed,domain=.25:6.2,samples=40] plot (\x,{\x*\x-6*\x+10});
  \foreach \x/\y in {1/5,2/2,3/1,4/2,5/5,6/10} \fill[accent] (\x,\y) circle (1.9pt);
\end{tikzpicture}
\end{minipage}
\end{theorie}

\voorbeeld{}{Teken de grafiek van $u_n = (-1)^n \cdot \dfrac{12}{n}$ voor $n = 1$ tot en met 6.\label{vb:h1-grafiek}}{%
\begin{minipage}[t]{.5\linewidth}
\begin{tabular}{@{}l*{6}{>{\centering\arraybackslash}p{7.5mm}}@{}}
$n$ & 1 & 2 & 3 & 4 & 5 & 6 \\
$u_n$ & $-12$ & $6$ & $-4$ & $3$ & $-2{,}4$ & $2$
\end{tabular}

\vspace{2mm}
Bereken eerst de termen en zet ze dan als losse punten uit. De punten liggen om en om onder en boven de horizontale as. Een lijn tussen de punten zou iets tonen dat er niet is.
\end{minipage}\hfill
\begin{minipage}[t]{.45\linewidth}
\vspace{-2mm}
\begin{tikzpicture}[x=.68cm,y=.13cm,>={Stealth[length=1.6mm]},font=\scriptsize]
  \draw[->,lijngrijs] (-.1,0) -- (6.9,0) node[right,text=black] {$n$};
  \draw[->,lijngrijs] (0,-13.5) -- (0,8.5) node[right,text=black] {$u_n$};
  % oneven n: punt onder de as, dus label erboven; even n omgekeerd
  \foreach \x in {1,3,5} \draw[lijngrijs] (\x,-.3) -- (\x,.3) node[above=-1pt,text=grijs] {\x};
  \foreach \x in {2,4,6} \draw[lijngrijs] (\x,.3) -- (\x,-.3) node[below=-1pt,text=grijs] {\x};
  \foreach \y in {-10,-5,5} \draw[lijngrijs] (.07,\y) -- (-.07,\y) node[left=-1pt,text=grijs] {$\y$};
  \foreach \x/\y in {1/-12,2/6,3/-4,4/3,5/-2.4,6/2} \fill[accent] (\x,\y) circle (1.9pt);
\end{tikzpicture}
\end{minipage}}

\begin{oefeningen}
\begin{oefening}
Teken de grafiek van $u_n = (n-3)^2$ voor $n = 1$ tot en met 6.
\begin{solution}
De punten $(1,4)$, $(2,1)$, $(3,0)$, $(4,1)$, $(5,4)$ en $(6,9)$.
\end{solution}
\begin{fullsolution}
De termen zijn 4, 1, 0, 1, 4 en 9. Zet de punten $(1,4)$, $(2,1)$, $(3,0)$, $(4,1)$, $(5,4)$ en $(6,9)$ los uit, zonder ze te verbinden.
\end{fullsolution}
\end{oefening}
\end{oefeningen}

%----------------------------------------------------------------------------------------
% OPDRACHTEN: door elkaar. Niveau 1, dan 2, dan 3.
%----------------------------------------------------------------------------------------
\opdrachtenkop

\lkr{De opdrachten staan bewust door elkaar, met stijgende moeilijkheid. In deze versie staat bij elke opdracht het onderwerp, zodat je kunt kiezen.}

\begin{oefeningen}
\begin{oefening}[1]\label{oef:h1-termen}
\thema{termen en volgnummers}Gegeven is de rij $6,\ 2,\ 9,\ 2,\ 7,\ 0,\ 4$.
\begin{enumerate}
  \item Wat is $u_3$?
  \item Welk volgnummer heeft de term 7?
  \item Voor welke volgnummers is $u_n = 2$?
  \item Neem $n = 6$. Wat zijn $u_{n-1}$ en $u_{n+1}$?
\end{enumerate}
\begin{solution}
\textbf{a.}~9 \quad \textbf{b.}~5 \quad \textbf{c.}~$n=2$ of $n=4$ \quad \textbf{d.}~7 en 4
\end{solution}
\begin{fullsolution}
\textbf{a.}~De derde term is 9.\\
\textbf{b.}~Het getal 7 staat op plaats 5.\\
\textbf{c.}~Het getal 2 staat op plaats 2 en plaats 4.\\
\textbf{d.}~$u_{n-1}=u_5=7$ en $u_{n+1}=u_7=4$.
\end{fullsolution}
\end{oefening}

\begin{oefening}[1]\label{oef:h1-expl}
\thema{expliciet: termen}Bepaal de eerste vier termen.
\begin{reeks}
  \item $u_n = 4n - 7$
  \item $u_n = n^2 + n$
  \item $u_n = \tfrac{6}{n}$
  \item $u_n = (-1)^n \cdot 3$
  \item $u_n = 2^n - 1$
  \item $u_n = \tfrac{n-1}{n+1}$
\end{reeks}
\begin{solution}
\textbf{a.}~$-3,\ 1,\ 5,\ 9$ \quad \textbf{b.}~2, 6, 12, 20 \quad \textbf{c.}~6, 3, 2, $\tfrac32$ \quad \textbf{d.}~$-3,\ 3,\ {-3},\ 3$ \quad \textbf{e.}~1, 3, 7, 15 \quad \textbf{f.}~0, $\tfrac13$, $\tfrac12$, $\tfrac35$
\end{solution}
\begin{fullsolution}
Vervang telkens $n$ door 1, 2, 3 en 4.\\
\textbf{a.}~$4-7=-3$, $8-7=1$, $12-7=5$, $16-7=9$.\\
\textbf{b.}~$1+1=2$, $4+2=6$, $9+3=12$, $16+4=20$.\\
\textbf{c.}~$\tfrac61=6$, $\tfrac62=3$, $\tfrac63=2$, $\tfrac64=\tfrac32$.\\
\textbf{d.}~Oneven $n$ geeft $-3$, even $n$ geeft 3.\\
\textbf{e.}~$2-1=1$, $4-1=3$, $8-1=7$, $16-1=15$.\\
\textbf{f.}~$\tfrac02=0$, $\tfrac13$, $\tfrac24=\tfrac12$, $\tfrac35$.
\end{fullsolution}
\end{oefening}

\begin{oefening}[1]\label{oef:h1-recterm}
\thema{recursief: termen}Bepaal de eerste vijf termen. Het voorschrift geldt telkens vanaf de eerste term die niet gegeven is.
\begin{enumerate}
  \item $u_1 = 2$ en $u_n = u_{n-1} + 7$
  \item $u_1 = 81$ en $u_n = \dfrac{u_{n-1}}{3}$
  \item $u_1 = 1$ en $u_n = 2u_{n-1} + 3$
  \item $u_1 = 5$ en $u_n = u_{n-1} - n$
  \item $u_1 = 2$, $u_2 = -1$ en $u_n = u_{n-1} + u_{n-2}$
  \item $u_1 = 1$, $u_2 = 3$ en $u_n = u_{n-1} \cdot u_{n-2}$
\end{enumerate}
\begin{solution}
\textbf{a.}~2, 9, 16, 23, 30 \quad \textbf{b.}~81, 27, 9, 3, 1 \quad \textbf{c.}~1, 5, 13, 29, 61 \quad \textbf{d.}~5, 3, 0, $-4$, $-9$ \quad \textbf{e.}~2, $-1$, 1, 0, 1 \quad \textbf{f.}~1, 3, 3, 9, 27
\end{solution}
\begin{fullsolution}
\textbf{a.}~Telkens 7 erbij.\\
\textbf{b.}~Telkens delen door 3.\\
\textbf{c.}~$2\cdot 1+3=5$, $2\cdot 5+3=13$, $2\cdot 13+3=29$, $2\cdot 29+3=61$.\\
\textbf{d.}~$5-2=3$, $3-3=0$, $0-4=-4$, $-4-5=-9$.\\
\textbf{e.}~$-1+2=1$, $1+(-1)=0$, $0+1=1$.\\
\textbf{f.}~$3\cdot 1=3$, $3\cdot 3=9$, $9\cdot 3=27$.
\end{fullsolution}
\end{oefening}

\begin{oefening}[1]\label{oef:h1-opstellen}
\thema{expliciet: opstellen}Bepaal een expliciet voorschrift.
\begin{reeks}
  \item $3,\ 6,\ 9,\ 12,\ \ldots$
  \item $5,\ 7,\ 9,\ 11,\ \ldots$
  \item $1,\ 8,\ 27,\ 64,\ \ldots$
  \item $2,\ 4,\ 8,\ 16,\ \ldots$
  \item $\tfrac12,\ \tfrac13,\ \tfrac14,\ \tfrac15,\ \ldots$
  \item $1,\ {-1},\ 1,\ {-1},\ \ldots$
\end{reeks}
\begin{solution}
\textbf{a.}~$u_n=3n$ \quad \textbf{b.}~$u_n=2n+3$ \quad \textbf{c.}~$u_n=n^3$ \quad \textbf{d.}~$u_n=2^n$ \quad \textbf{e.}~$u_n=\dfrac{1}{n+1}$ \quad \textbf{f.}~$u_n=(-1)^{n+1}$
\end{solution}
\begin{fullsolution}
\textbf{a.}~De drievouden: $3n$.\\
\textbf{b.}~Telkens 2 meer: vergelijk met $2n$. Elke term is 3 meer, dus $2n+3$.\\
\textbf{c.}~De derdemachten: $n^3$.\\
\textbf{d.}~De machten van 2: $2^n$.\\
\textbf{e.}~De noemer is 1 meer dan het volgnummer: $\dfrac{1}{n+1}$.\\
\textbf{f.}~Positief bij oneven $n$: $(-1)^{n+1}$.
\end{fullsolution}
\end{oefening}

\begin{oefening}[1]\label{oef:h1-plant}
\thema{starten bij 0}Een plant is bij het planten 6 cm hoog. $u_n$ is de hoogte in cm na $n$ weken. Er geldt $u_1 = 8$, $u_2 = 11$ en $u_3 = 13$.
\begin{enumerate}
  \item Wat is $u_0$?
  \item Hoeveel groeit de plant in de derde week?
  \item Hoeveel termen tel je van $u_0$ tot en met $u_{20}$?
\end{enumerate}
\begin{solution}
\textbf{a.}~6 \quad \textbf{b.}~2 cm \quad \textbf{c.}~21
\end{solution}
\begin{fullsolution}
\textbf{a.}~De beginhoogte: $u_0=6$.\\
\textbf{b.}~De derde week loopt van $u_2$ tot $u_3$: $13-11=2$ cm.\\
\textbf{c.}~$u_1$ tot en met $u_{20}$ zijn 20 termen, en $u_0$ komt erbij: 21.
\end{fullsolution}
\end{oefening}

\begin{oefening}[1]\label{oef:h1-recopst}
\thema{recursief: opstellen}Bepaal een recursief voorschrift.
\begin{reeks}
  \item $5,\ 8,\ 11,\ 14,\ \ldots$
  \item $3,\ {-6},\ 12,\ {-24},\ \ldots$
  \item $48,\ 24,\ 12,\ 6,\ \ldots$
  \item $1,\ 3,\ 6,\ 10,\ \ldots$
\end{reeks}
\begin{solution}
\textbf{a.}~$u_n=u_{n-1}+3$, $u_1=5$ \quad \textbf{b.}~$u_n=-2u_{n-1}$, $u_1=3$ \quad \textbf{c.}~$u_n=\tfrac12 u_{n-1}$, $u_1=48$ \quad \textbf{d.}~$u_n=u_{n-1}+n$, $u_1=1$
\end{solution}
\begin{fullsolution}
\textbf{a.}~Telkens 3 erbij.\\
\textbf{b.}~Telkens maal $-2$.\\
\textbf{c.}~Telkens de helft.\\
\textbf{d.}~Er komt 2, 3, 4 bij: bij de stap naar $u_n$ is dat precies $n$. Controle: $u_2=1+2=3$.
\end{fullsolution}
\end{oefening}

\begin{oefening}[1]\label{oef:h1-isterm}
\thema{term of geen term}Is het getal een term van de rij?
\begin{enumerate}
  \item $u_n = 3n + 1$: het getal 40, en het getal 50
  \item $u_n = n^2 - n$: het getal 30, en het getal 40
\end{enumerate}
\begin{solution}
\textbf{a.}~40 is de 13de term; 50 is geen term \quad \textbf{b.}~30 is de 6de term; 40 is geen term
\end{solution}
\begin{fullsolution}
\textbf{a.}~$3n+1=40$ geeft $n=13$. $3n+1=50$ geeft $n=\tfrac{49}{3}$: geen volgnummer.\\
\textbf{b.}~$n^2-n-30=0$ geeft $(n-6)(n+5)=0$, dus $n=6$. $n^2-n-40=0$ geeft $n=\dfrac{1\pm\sqrt{161}}{2}$, dus $n\approx 6{,}84$ of $n\approx -5{,}84$: geen volgnummer.
\end{fullsolution}
\end{oefening}

\begin{oefening}[1]
\thema{fouten: teken}Iemand schrijft bij de rij $-1,\ 2,\ {-3},\ 4,\ \ldots$ het voorschrift $u_n = (-1)^{n+1} \cdot n$. Wat is er fout? Geef het juiste voorschrift.
\begin{solution}
Het voorschrift begint met $+1$; juist is $u_n=(-1)^n\cdot n$.
\end{solution}
\begin{fullsolution}
Voor $n=1$ geeft het voorschrift $(-1)^2\cdot 1=1$, maar de eerste term is $-1$. De termen zijn negatief bij oneven $n$, dus $u_n=(-1)^n\cdot n$.
\end{fullsolution}
\end{oefening}

\begin{oefening}[1]\label{oef:h1-grafiek}
\thema{grafiek}Welke grafiek hoort bij welk voorschrift?

\vspace{1mm}
\noindent\textbf{1.}~$u_n = \dfrac{6}{n}$ \qquad \textbf{2.}~$u_n = 7 - n$ \qquad \textbf{3.}~$u_n = (-1)^n \cdot 3$ \qquad \textbf{4.}~$u_n = (n-3)^2$

\vspace{1mm}
\noindent\minigrafiek{A}{-1}{10}{5,10}{1/4,2/1,3/0,4/1,5/4,6/9}\hfill
\minigrafiek{B}{-.5}{7}{3,6}{1/6,2/3,3/2,4/1.5,5/1.2,6/1}\hfill
\minigrafiek{C}{-4}{4}{-3,3}{1/-3,2/3,3/-3,4/3,5/-3,6/3}\hfill
\minigrafiek{D}{-.5}{7}{3,6}{1/6,2/5,3/4,4/3,5/2,6/1}
\begin{solution}
1 bij B; 2 bij D; 3 bij C; 4 bij A
\end{solution}
\begin{fullsolution}
Bereken telkens enkele termen.\\
\textbf{1.}~$6,\ 3,\ 2,\ 1{,}5,\ \ldots$: grafiek B.\\
\textbf{2.}~$6,\ 5,\ 4,\ 3,\ \ldots$: grafiek D. Ze begint ook bij 6, dus de tweede term beslist.\\
\textbf{3.}~$-3,\ 3,\ {-3},\ \ldots$: grafiek C.\\
\textbf{4.}~$4,\ 1,\ 0,\ 1,\ \ldots$: grafiek A.
\end{fullsolution}
\end{oefening}

\begin{oefening}[1]
\thema{recursief: begintermen}Twee leerlingen gebruiken het voorschrift $u_n = 3u_{n-1} - 2$. De ene vindt $2,\ 4,\ 10,\ \ldots$ De andere vindt $1,\ 1,\ 1,\ \ldots$ Wie heeft gelijk?
\begin{solution}
Beiden: ze gebruiken een andere beginterm.
\end{solution}
\begin{fullsolution}
Beide rijen voldoen: $3\cdot 2-2=4$ en $3\cdot 4-2=10$, en ook $3\cdot 1-2=1$. Ze starten met een andere beginterm. Zonder beginterm ligt de rij niet vast.
\end{fullsolution}
\end{oefening}

\begin{oefening}[1]\label{oef:h1-subst}
\thema{expliciet: index invullen}Gegeven: $u_n = n^2 + 1$. Schrijf zo eenvoudig mogelijk.
\begin{reeks}[4]
  \item $u_{n+1}$
  \item $u_n + 1$
  \item $u_{2n}$
  \item $2u_n$
\end{reeks}
\begin{solution}
\textbf{a.}~$n^2+2n+2$ \quad \textbf{b.}~$n^2+2$ \quad \textbf{c.}~$4n^2+1$ \quad \textbf{d.}~$2n^2+2$
\end{solution}
\begin{fullsolution}
\textbf{a.}~$(n+1)^2+1=n^2+2n+2$.\\
\textbf{b.}~$n^2+1+1=n^2+2$.\\
\textbf{c.}~$(2n)^2+1=4n^2+1$.\\
\textbf{d.}~$2(n^2+1)=2n^2+2$.\\
Bij a en c verandert het volgnummer, bij b en d de term.
\end{fullsolution}
\end{oefening}

\begin{oefening}[1]
\thema{fouten: recursief}Iemand schrijft bij de rij $3,\ 4,\ 6,\ 9,\ 13,\ \ldots$ het voorschrift $u_n = u_{n-1} + n$ met $u_1 = 3$. Wat is er fout? Geef het juiste voorschrift.
\begin{solution}
Dat geeft $u_2=5$; juist is $u_n=u_{n-1}+(n-1)$ met $u_1=3$.
\end{solution}
\begin{fullsolution}
Controleer de eerste stap: het voorschrift geeft $u_2=3+2=5$, maar de tweede term is 4. Er komt 1, 2, 3, 4 bij: bij de stap naar $u_n$ is dat $n-1$. Dus $u_n=u_{n-1}+(n-1)$ met $u_1=3$.
\end{fullsolution}
\end{oefening}

\begin{oefening}[1]
\thema{expliciet: teken}Bepaal de eerste vier termen. Welke rijen hebben een wisselend teken?
\begin{reeks}
  \item $u_n = (-2)^n$
  \item $u_n = -2^n$
  \item $u_n = -(-2)^n$
\end{reeks}
\begin{solution}
\textbf{a.}~$-2,\ 4,\ {-8},\ 16$ \quad \textbf{b.}~$-2,\ {-4},\ {-8},\ {-16}$ \quad \textbf{c.}~$2,\ {-4},\ 8,\ {-16}$; a en c wisselen
\end{solution}
\begin{fullsolution}
\textbf{a.}~De macht hoort bij $-2$: $-2,\ 4,\ {-8},\ 16$.\\
\textbf{b.}~De macht hoort alleen bij 2: $-2,\ {-4},\ {-8},\ {-16}$.\\
\textbf{c.}~Het tegengestelde van a: $2,\ {-4},\ 8,\ {-16}$.\\
Bij a en c wisselt het teken.
\end{fullsolution}
\end{oefening}

\begin{oefening}[2]\label{oef:h1-nuleen}
\thema{starten bij 0}Gegeven is de rij $10,\ 14,\ 18,\ 22,\ \ldots$
\begin{enumerate}
  \item Geef een expliciet voorschrift als de eerste term $u_1$ heet.
  \item Geef een expliciet voorschrift als de eerste term $u_0$ heet.
  \item Welk volgnummer heeft de term 50 in elk geval?
\end{enumerate}
\begin{solution}
\textbf{a.}~$u_n=4n+6$ \quad \textbf{b.}~$u_n=4n+10$ \quad \textbf{c.}~$n=11$; $n=10$
\end{solution}
\begin{fullsolution}
\textbf{a.}~$u_n=4n+6$. Controle: $u_1=10$.\\
\textbf{b.}~$u_n=4n+10$. Controle: $u_0=10$.\\
\textbf{c.}~$4n+6=50$ geeft $n=11$. $4n+10=50$ geeft $n=10$.
\end{fullsolution}
\end{oefening}

\begin{oefening}[2]
\begin{metprent}{schaakbord}
\thema{context: expliciet en recursief}In een toernooi speelt elke speler één keer tegen elke andere speler. $u_n$ is het aantal wedstrijden met $n$ spelers.
\begin{enumerate}
  \item Bereken $u_1$, $u_2$, $u_3$ en $u_4$.
  \item Een nieuwe speler sluit aan. Leg daarmee uit waarom \mbox{$u_n = u_{n-1} + (n-1)$}.
  \item Leg uit waarom $u_n = \dfrac{n(n-1)}{2}$.
  \item Hoeveel spelers doen mee als er 45 wedstrijden zijn?
\end{enumerate}
\end{metprent}
\begin{solution}
\textbf{a.}~0, 1, 3, 6 \quad \textbf{d.}~10
\end{solution}
\begin{fullsolution}
\textbf{a.}~Eén speler speelt niet. Daarna: 1, 3 en 6 wedstrijden.\\
\textbf{b.}~De $n$-de speler speelt tegen de $n-1$ spelers die er al waren. De andere wedstrijden blijven dezelfde.\\
\textbf{c.}~Elk van de $n$ spelers speelt $n-1$ wedstrijden. Zo tel je elke wedstrijd twee keer, dus deel je door 2.\\
\textbf{d.}~$n(n-1)=90$ geeft $(n-10)(n+9)=0$, dus $n=10$.
\end{fullsolution}
\end{oefening}

\begin{oefening}[2]
\thema{expliciet: opstellen}Bepaal een expliciet voorschrift.
\begin{reeks}
  \item $\tfrac23,\ \tfrac34,\ \tfrac45,\ \tfrac56,\ \ldots$
  \item $\tfrac12,\ {-\tfrac14},\ \tfrac16,\ {-\tfrac18},\ \ldots$
  \item $3,\ 6,\ 11,\ 18,\ \ldots$
  \item $2,\ 6,\ 12,\ 20,\ \ldots$
  \item $4,\ 0,\ 4,\ 0,\ \ldots$
  \item $2,\ \tfrac43,\ \tfrac85,\ \tfrac{16}{7},\ \ldots$
\end{reeks}
\begin{solution}
\textbf{a.}~$u_n=\dfrac{n+1}{n+2}$ \quad \textbf{b.}~$u_n=\dfrac{(-1)^{n+1}}{2n}$ \quad \textbf{c.}~$u_n=n^2+2$ \quad \textbf{d.}~$u_n=n(n+1)$ \quad \textbf{e.}~$u_n=2+2\cdot(-1)^{n+1}$ \quad \textbf{f.}~$u_n=\dfrac{2^n}{2n-1}$
\end{solution}
\begin{fullsolution}
\textbf{a.}~Teller $n+1$, noemer $n+2$.\\
\textbf{b.}~Positief bij oneven $n$, noemer $2n$.\\
\textbf{c.}~Elke term is 2 meer dan een kwadraat.\\
\textbf{d.}~$1\cdot 2$, $2\cdot 3$, $3\cdot 4$, $4\cdot 5$: dus $n(n+1)$.\\
\textbf{e.}~Het gemiddelde van 4 en 0 is 2, en daarrond schommelt de rij met 2: $2+2\cdot(-1)^{n+1}$.\\
\textbf{f.}~Schrijf $2=\tfrac21$. De tellers zijn $2^n$ en de noemers de oneven getallen $2n-1$.
\end{fullsolution}
\end{oefening}

\begin{oefening}[2]
\thema{term: ongelijkheid}Vanaf welk volgnummer is de term groter dan 1000?
\begin{reeks}[2]
  \item $u_n = n^2 + 5n$
  \item $u_n = 5 \cdot 2^n$
\end{reeks}
\begin{solution}
\textbf{a.}~$n=30$ \quad \textbf{b.}~$n=8$
\end{solution}
\begin{fullsolution}
\textbf{a.}~$n^2+5n-1000=0$ geeft $n\approx 29{,}22$, de negatieve oplossing valt weg. Dus vanaf $n=30$. Controle: $u_{29}=986$ en $u_{30}=1050$.\\
\textbf{b.}~$5\cdot 2^n>1000$ geeft $2^n>200$. Omdat $2^7=128$ en $2^8=256$, is dat vanaf $n=8$.
\end{fullsolution}
\end{oefening}

\begin{oefening}[2]\label{oef:h1-grafvs}
\begin{minipage}[t]{.55\linewidth}
\thema{grafiek en voorschrift}De grafiek toont de eerste termen van een rij. Geef een expliciet en een recursief voorschrift. Bereken $u_{50}$.
\end{minipage}\hfill
\begin{minipage}[t]{.38\linewidth}
\vspace{-3mm}
\minigrafiek{}{-.5}{9}{2,4,6,8}{1/8,2/6,3/4,4/2,5/0}
\end{minipage}
\begin{solution}
$u_n=10-2n$; $u_n=u_{n-1}-2$, $u_1=8$; $u_{50}=-90$
\end{solution}
\begin{fullsolution}
De termen zijn $8,\ 6,\ 4,\ 2,\ 0$. Telkens 2 minder, dus $u_n=u_{n-1}-2$ met $u_1=8$. Vergelijk met $-2n$: $-2,\ {-4},\ {-6}$. Elke term is 10 meer, dus $u_n=10-2n$. Dan is $u_{50}=10-100=-90$.
\end{fullsolution}
\end{oefening}

\begin{oefening}[2]
\thema{recursief: index verschoven}Gegeven: $u_1 = 3$ en $u_{n+1} = u_n^{\,2} - 2$ voor $n \geqslant 1$.
\begin{enumerate}
  \item Bereken $u_2$, $u_3$ en $u_4$.
  \item Schrijf hetzelfde voorschrift met $u_n$ en $u_{n-1}$. Vanaf welke $n$ geldt het dan?
\end{enumerate}
\begin{solution}
\textbf{a.}~7, 47, 2207 \quad \textbf{b.}~$u_n=u_{n-1}^{\,2}-2$ voor $n\geqslant 2$
\end{solution}
\begin{fullsolution}
\textbf{a.}~$u_2=9-2=7$, $u_3=49-2=47$, $u_4=2209-2=2207$.\\
\textbf{b.}~Vervang $n$ door $n-1$: $u_n=u_{n-1}^{\,2}-2$. De voorwaarde $n-1\geqslant 1$ wordt $n\geqslant 2$.
\end{fullsolution}
\end{oefening}

\begin{oefening}[2]\label{oef:h1-omzet}
\thema{van expliciet naar recursief}Bepaal een recursief voorschrift.
\begin{reeks}
  \item $u_n = 5n - 2$
  \item $u_n = 3 \cdot 2^n$
  \item $u_n = n^2 + n$
\end{reeks}
\begin{solution}
\textbf{a.}~$u_n=u_{n-1}+5$, $u_1=3$ \quad \textbf{b.}~$u_n=2u_{n-1}$, $u_1=6$ \quad \textbf{c.}~$u_n=u_{n-1}+2n$, $u_1=2$
\end{solution}
\begin{fullsolution}
\textbf{a.}~$u_n-u_{n-1}=5n-2-(5n-7)=5$ en $u_1=3$.\\
\textbf{b.}~$\dfrac{u_n}{u_{n-1}}=\dfrac{3\cdot 2^n}{3\cdot 2^{n-1}}=2$ en $u_1=6$.\\
\textbf{c.}~$u_{n-1}=(n-1)^2+(n-1)=n^2-n$, dus $u_n-u_{n-1}=2n$. Verder is $u_1=2$.
\end{fullsolution}
\end{oefening}

\begin{oefening}[2]
\thema{voorschrift in woorden}Zeg elke term hardop en je vindt de volgende. De 21 lees je als: één 2, één 1. Wat zijn de volgende twee termen van deze rij?\\[.5ex]
$1,\ 11,\ 21,\ 1211,\ 111221,\ \ldots$
\begin{solution}
312211 en 13112221
\end{solution}
\begin{fullsolution}
111221 lees je als: drie 1, twee 2, één 1, dus 312211. Dat lees je als: één 3, één 1, twee 2, twee 1, dus 13112221. Dit is een recursief voorschrift in woorden, zonder formule.
\end{fullsolution}
\end{oefening}

\begin{oefening}[2]
\thema{term of geen term}Gegeven: $u_n = \dfrac{n+1}{2n}$. Zijn de volgende getallen termen van de rij?
\begin{reeks}
  \item $\tfrac35$
  \item $\tfrac32$
  \item $\tfrac12$
\end{reeks}
\begin{solution}
\textbf{a.}~ja, de 5de term \quad \textbf{b.}~nee \quad \textbf{c.}~nee
\end{solution}
\begin{fullsolution}
\textbf{a.}~$5(n+1)=6n$ geeft $n=5$.\\
\textbf{b.}~$2(n+1)=6n$ geeft $n=\tfrac12$: geen volgnummer.\\
\textbf{c.}~$2(n+1)=2n$ geeft $2=0$: geen oplossing. De termen komen wel steeds dichter bij $\tfrac12$.
\end{fullsolution}
\end{oefening}

\begin{oefening}[2]\label{oef:h1-recexpl}
\thema{van recursief naar expliciet}Bereken de eerste vijf termen en vermoed een expliciet voorschrift.
\begin{enumerate}
  \item $u_1 = 3$ en $u_n = u_{n-1} + 4$
  \item $u_1 = 2$ en $u_n = 3u_{n-1}$
  \item $u_1 = 1$ en $u_n = u_{n-1} + 2n - 1$
  \item $u_1 = 1$ en $u_n = 2u_{n-1} + 1$
\end{enumerate}
\begin{solution}
\textbf{a.}~$u_n=4n-1$ \quad \textbf{b.}~$u_n=2\cdot 3^{n-1}$ \quad \textbf{c.}~$u_n=n^2$ \quad \textbf{d.}~$u_n=2^n-1$
\end{solution}
\begin{fullsolution}
\textbf{a.}~3, 7, 11, 15, 19: telkens 1 minder dan $4n$.\\
\textbf{b.}~2, 6, 18, 54, 162: telkens maal 3, en bij $u_n$ is dat $n-1$ keer gebeurd.\\
\textbf{c.}~1, 4, 9, 16, 25: de kwadraten.\\
\textbf{d.}~1, 3, 7, 15, 31: telkens 1 minder dan een macht van 2.
\end{fullsolution}
\end{oefening}

\begin{oefening}[2]
\begin{minipage}[t]{.72\linewidth}
\thema{context: expliciet en recursief}Een kaartenhuis met $n$ verdiepingen telt $u_n$ kaarten: $u_1 = 2$, $u_2 = 7$ en $u_3 = 15$. De figuur toont het kaartenhuis met 3 verdiepingen.
\end{minipage}\hfill
\begin{minipage}[t]{.24\linewidth}
\vspace{-3mm}
\begin{tikzpicture}[x=4.2mm,y=5.6mm,line width=.7pt]
  \foreach \r in {0,1,2}{
    \pgfmathtruncatemacro{\k}{3-\r}
    \foreach \i in {1,...,\k}{\draw (\i-1+\r/2,\r) -- (\i-.5+\r/2,\r+1) -- (\i+\r/2,\r);}
    \ifnum\k>1 \draw (.5+\r/2,\r+1) -- (\k-.5+\r/2,\r+1);\fi}
\end{tikzpicture}
\end{minipage}
\begin{enumerate}
  \item Controleer dat $u_n = \dfrac{n(3n+1)}{2}$ klopt voor deze drie termen, en bereken $u_{20}$.
  \item Een nieuwe onderste verdieping telt $n$ paar schuine kaarten en $n-1$ platte kaarten. Leid daaruit een recursief voorschrift af.
  \item Hoeveel volledige verdiepingen bouw je met een spel van 52 kaarten?
\end{enumerate}
\begin{solution}
\textbf{a.}~$u_{20}=610$ \quad \textbf{b.}~$u_n=u_{n-1}+3n-1$, $u_1=2$ \quad \textbf{c.}~5
\end{solution}
\begin{fullsolution}
\textbf{a.}~$\tfrac{1\cdot 4}{2}=2$, $\tfrac{2\cdot 7}{2}=7$, $\tfrac{3\cdot 10}{2}=15$, en $u_{20}=\tfrac{20\cdot 61}{2}=610$.\\
\textbf{b.}~Er komen $2n+(n-1)=3n-1$ kaarten bij: $u_n=u_{n-1}+3n-1$ met $u_1=2$.\\
\textbf{c.}~$u_5=40$ en $u_6=57$, dus 5 verdiepingen.
\end{fullsolution}
\end{oefening}

\begin{oefening}[2]
\thema{recursief: terugrekenen}Gegeven: $u_n = 2u_{n-1} + 2$ en $u_4 = 20$. Bereken $u_3$, $u_2$ en $u_1$.
\begin{solution}
9; 3,5; 0,75
\end{solution}
\begin{fullsolution}
Uit $u_n=2u_{n-1}+2$ volgt $u_{n-1}=\dfrac{u_n-2}{2}$. Dus $u_3=\tfrac{18}{2}=9$, $u_2=\tfrac72=3{,}5$ en $u_1=\tfrac{1{,}5}{2}=0{,}75$.
\end{fullsolution}
\end{oefening}

\begin{oefening}[2]
\thema{expliciet: een begin bepaalt geen rij}Gegeven: $u_n = 2^{n-1}$ en $v_n = 2^{n-1} + (n-1)(n-2)(n-3)(n-4)$.
\begin{enumerate}
  \item Bereken van beide rijen de eerste vijf termen.
  \item Iemand zegt: `De rij $1,\ 2,\ 4,\ 8,\ \ldots$ gaat zeker verder met 16.' Wat antwoord je?
\end{enumerate}
\begin{solution}
\textbf{a.}~1, 2, 4, 8, 16 en 1, 2, 4, 8, 40 \quad \textbf{b.}~Niet zeker: de eerste termen leggen de rij niet vast.
\end{solution}
\begin{fullsolution}
\textbf{a.}~Voor $n=1$, 2, 3 en 4 is het product nul, dus beide rijen zijn gelijk. Voor $n=5$ is het product $4\cdot 3\cdot 2\cdot 1=24$, dus $v_5=16+24=40$.\\
\textbf{b.}~Een eindig aantal termen legt de rij niet vast. De puntjes zijn enkel duidelijk als het voorschrift gegeven is of de regelmaat afgesproken is.
\end{fullsolution}
\end{oefening}

\begin{oefening}[2]
\thema{context: recursief}Een zwembad bevat op een maandag 400 g chloor. Elke nacht verdwijnt 20\% van het chloor. Elke ochtend voegt de beheerder 50 g toe. $u_n$ is de hoeveelheid chloor na $n$ dagen.
\begin{enumerate}
  \item Waarom is $u_n = 0{,}8\,u_{n-1} + 50$? Wat is $u_0$?
  \item Bereken $u_1$, $u_2$ en $u_3$.
  \item Bij welke hoeveelheid chloor verandert er niets meer?
\end{enumerate}
\begin{solution}
\textbf{a.}~$u_0=400$ \quad \textbf{b.}~370; 346; 326,8 g \quad \textbf{c.}~250 g
\end{solution}
\begin{fullsolution}
\textbf{a.}~Na het verlies blijft 80\% over, dat is $0{,}8\,u_{n-1}$. Daarna komt er 50 g bij. De beginwaarde is $u_0=400$.\\
\textbf{b.}~$u_1=320+50=370$, $u_2=296+50=346$, $u_3=276{,}8+50=326{,}8$.\\
\textbf{c.}~Er verandert niets als $x=0{,}8\,x+50$, dus $0{,}2\,x=50$ en $x=250$.
\end{fullsolution}
\end{oefening}

\begin{oefening}[2]
\thema{recursief met gevallen}Gegeven: $u_1 = 6$. Is $u_{n-1}$ even, dan is $u_n = \dfrac{u_{n-1}}{2}$. Is $u_{n-1}$ oneven, dan is $u_n = 3u_{n-1} + 1$.
\begin{enumerate}
  \item Schrijf de termen op tot je 1 bereikt. Welk volgnummer heeft die term?
  \item Wat gebeurt er als je daarna verder rekent?
  \item Begin nu met $u_1 = 7$. Welk volgnummer heeft de eerste term 1?
\end{enumerate}
Niemand heeft al kunnen bewijzen dat je vanuit elk natuurlijk getal ooit bij 1 komt.
\begin{solution}
\textbf{a.}~6, 3, 10, 5, 16, 8, 4, 2, 1; volgnummer 9 \quad \textbf{b.}~4, 2, 1 herhaalt zich \quad \textbf{c.}~17
\end{solution}
\begin{fullsolution}
\textbf{a.}~6, 3, 10, 5, 16, 8, 4, 2, 1: de negende term.\\
\textbf{b.}~Na 1 volgt $3\cdot 1+1=4$, en dan 2 en 1: de rij blijft 4, 2, 1 herhalen.\\
\textbf{c.}~7, 22, 11, 34, 17, 52, 26, 13, 40, 20, 10, 5, 16, 8, 4, 2, 1: de zeventiende term.
\end{fullsolution}
\end{oefening}

\begin{oefening}[2]
\thema{context: expliciet en recursief}$u_n$ is het aantal diagonalen van een $n$-hoek. Er geldt $u_3 = 0$, $u_4 = 2$ en $u_5 = 5$.
\begin{enumerate}
  \item Waarom start deze rij niet bij $u_1$?
  \item Leg uit waarom $u_n = \dfrac{n(n-3)}{2}$.
  \item Welke veelhoek heeft 35 diagonalen?
  \item Geef een recursief voorschrift.
\end{enumerate}
\begin{solution}
\textbf{a.}~Een veelhoek heeft minstens 3 hoekpunten \quad \textbf{c.}~een tienhoek \quad \textbf{d.}~$u_n=u_{n-1}+n-2$, $u_3=0$
\end{solution}
\begin{fullsolution}
\textbf{a.}~Een veelhoek heeft minstens 3 hoekpunten.\\
\textbf{b.}~Elk hoekpunt verbind je met alle andere hoekpunten, behalve met zijn twee buren: dat zijn $n-3$ diagonalen per hoekpunt. Zo tel je elke diagonaal twee keer.\\
\textbf{c.}~$n(n-3)=70$ geeft $(n-10)(n+7)=0$, dus $n=10$.\\
\textbf{d.}~$u_{n-1}=\dfrac{(n-1)(n-4)}{2}=\dfrac{n^2-5n+4}{2}$ en $u_n=\dfrac{n^2-3n}{2}$, dus $u_n-u_{n-1}=\dfrac{2n-4}{2}=n-2$.
\end{fullsolution}
\end{oefening}

\begin{oefening}[3]
\thema{periodieke rij}Gegeven: $u_1 = 2$, $u_2 = 3$ en $u_n = \dfrac{u_{n-1}}{u_{n-2}}$ voor $n \geqslant 3$.
\begin{enumerate}
  \item Bereken de eerste acht termen. Wat valt op?
  \item Bereken $u_{100}$.
  \item Toon aan dat hetzelfde gebeurt voor elke $u_1 = a$ en $u_2 = b$ met $a, b \neq 0$.
\end{enumerate}
\begin{solution}
\textbf{a.}~2, 3, $\tfrac32$, $\tfrac12$, $\tfrac13$, $\tfrac23$, 2, 3: om de 6 termen dezelfde \quad \textbf{b.}~$\tfrac12$
\end{solution}
\begin{fullsolution}
\textbf{a.}~$u_3=\tfrac32$, $u_4=\tfrac{3/2}{3}=\tfrac12$, $u_5=\tfrac{1/2}{3/2}=\tfrac13$, $u_6=\tfrac{1/3}{1/2}=\tfrac23$, $u_7=\tfrac{2/3}{1/3}=2$, $u_8=\tfrac{2}{2/3}=3$. Elke term hangt enkel af van de twee vorige, dus vanaf hier herhaalt alles zich om de 6 termen.\\
\textbf{b.}~$100=16\cdot 6+4$, dus $u_{100}=u_4=\tfrac12$.\\
\textbf{c.}~$a,\ b,\ \dfrac{b}{a},\ \dfrac{1}{a},\ \dfrac{1}{b},\ \dfrac{a}{b},\ a,\ b$: na zes termen staan $a$ en $b$ er opnieuw.
\end{fullsolution}
\end{oefening}

\begin{oefening}[3]
\thema{bewijs: voorschrift controleren}Gegeven: $u_1 = 1$ en $u_n = 2u_{n-1} + 1$. Toon aan dat $v_n = 2^n - 1$ aan dit recursieve voorschrift voldoet: controleer de beginterm, en vul $v_{n-1}$ in. Waarom zijn $u_n$ en $v_n$ dan dezelfde rij?
\begin{solution}
Beginterm en stap kloppen, dus de rijen zijn gelijk.
\end{solution}
\begin{fullsolution}
$v_1=2-1=1$ en $2v_{n-1}+1=2\left(2^{n-1}-1\right)+1=2^n-2+1=2^n-1=v_n$.\\
De beginterm en het voorschrift bepalen de rij volledig: uit $u_1$ volgt $u_2$, daaruit $u_3$, enzovoort. Twee rijen met dezelfde beginterm en hetzelfde voorschrift zijn dus gelijk.
\end{fullsolution}
\end{oefening}

\begin{oefening}[3]
\begin{metprent}{bij}
\thema{context: Fibonacci}Bij bijen heeft een mannetje alleen een moeder, en een vrouwtje een moeder en een vader. Vertrek van één mannetje. $u_n$ is het aantal voorouders $n$ generaties terug. Zo is $u_1 = 1$, de moeder, en $u_2 = 2$, haar ouders.
\begin{enumerate}
  \item Bereken $u_3$, $u_4$ en $u_5$.
  \item Leg uit waarom $u_n = u_{n-1} + u_{n-2}$ voor $n \geqslant 3$. Tel daarvoor apart de moeders en de vaders in generatie $n$.
\end{enumerate}
\end{metprent}
\begin{solution}
\textbf{a.}~3, 5, 8
\end{solution}
\begin{fullsolution}
\textbf{a.}~Generatie 2 telt 1 vrouwtje en 1 mannetje. Het vrouwtje heeft 2 ouders en het mannetje 1, dus $u_3=3$: 2 vrouwtjes en 1 mannetje. Zo vind je ook $u_4=2\cdot 2+1=5$: 3 vrouwtjes en 2 mannetjes, en $u_5=2\cdot 3+2=8$.\\
\textbf{b.}~Elke bij in generatie $n-1$ heeft één moeder, dus generatie $n$ telt $u_{n-1}$ vrouwtjes. Alleen de vrouwtjes in generatie $n-1$ hebben een vader. Dat aantal vrouwtjes is $u_{n-2}$, volgens dezelfde redenering. Samen: $u_n=u_{n-1}+u_{n-2}$.
\end{fullsolution}
\end{oefening}

\begin{oefening}[3]
\thema{bewijs: expliciet en recursief}Toon aan dat de rij $u_n = \dfrac{1}{n}$ voldoet aan $u_n = \dfrac{u_{n-1}}{1 + u_{n-1}}$ voor $n \geqslant 2$.
\begin{solution}
Vul $u_{n-1}=\dfrac{1}{n-1}$ in; dat geeft $\dfrac1n$.
\end{solution}
\begin{fullsolution}
$\dfrac{u_{n-1}}{1+u_{n-1}}=\dfrac{\dfrac{1}{n-1}}{1+\dfrac{1}{n-1}}=\dfrac{\dfrac{1}{n-1}}{\dfrac{n}{n-1}}=\dfrac{1}{n}=u_n$.
\end{fullsolution}
\end{oefening}

\begin{oefening}[3]
\thema{hulprij}Gegeven: $u_1 = k$ en $u_n = 2u_{n-1} - 3$.
\begin{enumerate}
  \item Voor welke $k$ is de rij constant?
  \item Neem $k = 5$ en stel $v_n = u_n - 3$. Toon aan dat $v_n = 2v_{n-1}$.
  \item Geef een expliciet voorschrift van $u_n$ voor $k = 5$.
\end{enumerate}
\begin{solution}
\textbf{a.}~$k=3$ \quad \textbf{c.}~$u_n=2^n+3$
\end{solution}
\begin{fullsolution}
\textbf{a.}~Constant betekent $u_2=u_1$, dus $k=2k-3$ en $k=3$. Dan blijft elke term 3.\\
\textbf{b.}~$v_n=u_n-3=2u_{n-1}-3-3=2(u_{n-1}-3)=2v_{n-1}$.\\
\textbf{c.}~$v_1=5-3=2$ en elke term van $v_n$ is het dubbele van de vorige, dus $v_n=2^n$. Dan is $u_n=v_n+3=2^n+3$. Controle: 5, 7, 11, 19.
\end{fullsolution}
\end{oefening}

\begin{oefening}[3]
\thema{Fibonacci: vermoeden en bewijs}Neem de rij van Fibonacci $1,\ 1,\ 2,\ 3,\ 5,\ 8,\ \ldots$ en stel $w_n = u_{n+1} \cdot u_{n-1} - u_n^{\,2}$.
\begin{enumerate}
  \item Bereken $w_2$, $w_3$, $w_4$ en $w_5$. Wat vermoed je?
  \item Toon aan: $w_{n+1} = -w_n$. Vervang daarvoor in $w_{n+1}$ eerst $u_{n+2}$ door $u_{n+1} + u_n$.
\end{enumerate}
\begin{solution}
\textbf{a.}~$1,\ {-1},\ 1,\ {-1}$: vermoeden $w_n=(-1)^n$
\end{solution}
\begin{fullsolution}
\textbf{a.}~$w_2=2\cdot 1-1=1$, $w_3=3\cdot 1-4=-1$, $w_4=5\cdot 2-9=1$, $w_5=8\cdot 3-25=-1$.\\
\textbf{b.}~$w_{n+1}=u_{n+2}u_n-u_{n+1}^{\,2}=(u_{n+1}+u_n)u_n-u_{n+1}^{\,2}=u_n^{\,2}+u_{n+1}(u_n-u_{n+1})$.\\
Omdat $u_{n+1}=u_n+u_{n-1}$ is $u_n-u_{n+1}=-u_{n-1}$. Dus $w_{n+1}=u_n^{\,2}-u_{n+1}u_{n-1}=-w_n$.
\end{fullsolution}
\end{oefening}
\end{oefeningen}

%----------------------------------------------------------------------------------------
% KOPPEL: verbanden met andere leerstof
%----------------------------------------------------------------------------------------
\opdrachtsectie{Koppel}

\lkr{Rationale functies, goniometrie en matrices zijn gezien in het 5de jaar. Telproblemen en logica staan ook in het 6de jaar: kijk na of ze al aan bod kwamen, of laat die twee oefeningen weg.}

\begin{oefeningen}
\begin{oefening}[2]
\koppel{analyse}{Rationale functies} Gegeven: $f(x) = \dfrac{x+1}{x-3}$ en $u_n = f(n)$.
\begin{enumerate}
  \item Bepaal het domein van $f$.
  \item Bereken $u_1$, $u_2$, $u_4$, $u_5$ en $u_6$. Waarom ontbreekt $u_3$?
  \item Hoe los je dat op, zodat er een rij ontstaat met losse termen zonder gat?
\end{enumerate}
\begin{solution}
\textbf{a.}~$\mathbb{R}\setminus\{3\}$ \quad \textbf{b.}~$-1$, $-3$, 5, 3, $\tfrac73$ \quad \textbf{c.}~Start bij $n=4$.
\end{solution}
\begin{fullsolution}
\textbf{a.}~De noemer is nul bij $x=3$.\\
\textbf{b.}~$\tfrac{2}{-2}=-1$, $\tfrac{3}{-1}=-3$, $\tfrac51=5$, $\tfrac62=3$, $\tfrac73$. De term $u_3$ bestaat niet omdat 3 niet in het domein ligt.\\
\textbf{c.}~Neem de rij pas vanaf $n=4$, dus met domein $\{4,5,6,\ldots\}$.
\end{fullsolution}
\end{oefening}

\begin{oefening}[2]
\koppel{analyse}{Goniometrische functies} Gegeven: $u_n = \cos\left(n \cdot \dfrac{\pi}{3}\right)$.
\begin{enumerate}
  \item Bereken de eerste zeven termen. Om de hoeveel termen herhaalt de rij zich?
  \item Wat is de periode van de functie $f(x) = \cos\left(x \cdot \dfrac{\pi}{3}\right)$? Waarom is dat hetzelfde getal?
  \item Schrijf $(-1)^n$ als een cosinus van een veelvoud van $n$.
\end{enumerate}
\begin{solution}
\textbf{a.}~$\tfrac12$, $-\tfrac12$, $-1$, $-\tfrac12$, $\tfrac12$, 1, $\tfrac12$; om de 6 \quad \textbf{b.}~6 \quad \textbf{c.}~$\cos(n\pi)$
\end{solution}
\begin{fullsolution}
\textbf{a.}~De hoeken zijn $\tfrac{\pi}{3},\ \tfrac{2\pi}{3},\ \pi,\ \tfrac{4\pi}{3},\ \tfrac{5\pi}{3},\ 2\pi,\ \tfrac{7\pi}{3}$. Vanaf daar herhaalt de cosinus zich.\\
\textbf{b.}~De periode is $\dfrac{2\pi}{\pi/3}=6$. De rij neemt de functiewaarden in de gehele getallen, en de periode is hier zelf een geheel getal.\\
\textbf{c.}~$\cos(n\pi)$ is $-1$ bij oneven $n$ en 1 bij even $n$.
\end{fullsolution}
\end{oefening}

\begin{oefening}[2]
\koppel{algebra}{Matrices} Gegeven: $A = \begin{pmatrix} 1 & 1 \\ 1 & 0 \end{pmatrix}$.
\begin{enumerate}
  \item Bereken $A^2$, $A^3$ en $A^4$. Welke rij herken je in de elementen?
  \item Toon aan: $A \begin{pmatrix} u_{n-1} \\ u_{n-2} \end{pmatrix} = \begin{pmatrix} u_n \\ u_{n-1} \end{pmatrix}$ is precies het voorschrift van die rij.
\end{enumerate}
\begin{solution}
\textbf{a.}~$\begin{pmatrix} 2 & 1 \\ 1 & 1 \end{pmatrix}$, $\begin{pmatrix} 3 & 2 \\ 2 & 1 \end{pmatrix}$, $\begin{pmatrix} 5 & 3 \\ 3 & 2 \end{pmatrix}$; de rij van Fibonacci
\end{solution}
\begin{fullsolution}
\textbf{a.}~$A^2=\begin{pmatrix} 2 & 1 \\ 1 & 1 \end{pmatrix}$, $A^3=\begin{pmatrix} 3 & 2 \\ 2 & 1 \end{pmatrix}$, $A^4=\begin{pmatrix} 5 & 3 \\ 3 & 2 \end{pmatrix}$. De elementen zijn termen van de rij van Fibonacci.\\
\textbf{b.}~Het product is $\begin{pmatrix} u_{n-1}+u_{n-2} \\ u_{n-1} \end{pmatrix}$. De eerste rij zegt $u_n=u_{n-1}+u_{n-2}$. De tweede rij zegt alleen dat $u_{n-1}$ gelijk blijft aan zichzelf.
\end{fullsolution}
\end{oefening}

\begin{oefening}[2]
\koppel{discreet}{Telproblemen} $u_n$ is het aantal manieren om $n$ verschillende boeken naast elkaar te zetten.
\begin{enumerate}
  \item Bereken $u_1$, $u_2$ en $u_3$.
  \item Leg uit waarom $u_n = n \cdot u_{n-1}$. Denk aan de plaatsen voor het laatste boek.
  \item Bereken $u_6$. Hoe noteren we $u_n$ met een speciaal symbool?
\end{enumerate}
\begin{solution}
\textbf{a.}~1, 2, 6 \quad \textbf{c.}~720; $u_n=n!$
\end{solution}
\begin{fullsolution}
\textbf{a.}~Eén boek: 1 manier. Twee boeken: 2 manieren. Drie boeken: 6 manieren.\\
\textbf{b.}~Zet eerst $n-1$ boeken: dat kan op $u_{n-1}$ manieren. Het laatste boek kan dan op $n$ plaatsen: vooraan, achteraan of ertussen.\\
\textbf{c.}~$u_4=24$, $u_5=120$, $u_6=720$. Dit is de faculteit: $u_n=n!$.
\end{fullsolution}
\end{oefening}

\begin{oefening}[2]
\koppel{discreet}{Logica en bewijzen} Gegeven: $u_n = n^2 + n + 41$.
\begin{enumerate}
  \item Bereken $u_1$, $u_2$, $u_3$ en $u_4$. Wat valt op?
  \item Iemand besluit: `Alle termen van deze rij zijn priemgetallen.' Weerleg dit met een tegenvoorbeeld. Bekijk daarvoor $u_{41}$ zonder rekenmachine.
\end{enumerate}
\begin{solution}
\textbf{a.}~43, 47, 53, 61: priemgetallen \quad \textbf{b.}~$u_{41}=41\cdot 43$
\end{solution}
\begin{fullsolution}
\textbf{a.}~43, 47, 53 en 61 zijn allemaal priemgetallen.\\
\textbf{b.}~$u_{41}=41^2+41+41=41(41+1+1)=41\cdot 43$: geen priemgetal. Enkele voorbeelden bewijzen dus niets, één tegenvoorbeeld weerlegt wel.
\end{fullsolution}
\end{oefening}
\end{oefeningen}

%----------------------------------------------------------------------------------------
\opdrachtsectie{Test jezelf}

\begin{oefeningen}
\begin{oefening}[1]\label{oef:h1-t1}
Gegeven is de rij $5,\ 1,\ 8,\ 1,\ 3$. Wat is $u_3$? Welk volgnummer heeft de term 3? Voor welke volgnummers is $u_n = 1$?
\begin{solution}
8; 5; $n=2$ of $n=4$
\end{solution}
\begin{fullsolution}
De derde term is 8. Het getal 3 staat op plaats 5. Het getal 1 staat op plaats 2 en plaats 4.
\end{fullsolution}
\end{oefening}

\begin{oefening}[1]\label{oef:h1-t2}
Bepaal de eerste vier termen van $u_n = (-1)^n \cdot (n+2)$.
\begin{solution}
$-3,\ 4,\ {-5},\ 6$
\end{solution}
\begin{fullsolution}
$(-1)^1\cdot 3=-3$, $(-1)^2\cdot 4=4$, $(-1)^3\cdot 5=-5$, $(-1)^4\cdot 6=6$.
\end{fullsolution}
\end{oefening}

\begin{oefening}[2]\label{oef:h1-t3}
Bepaal een expliciet voorschrift.
\begin{reeks}[2]
  \item $4,\ 7,\ 10,\ 13,\ \ldots$
  \item $\tfrac13,\ {-\tfrac24},\ \tfrac35,\ {-\tfrac46},\ \ldots$
\end{reeks}
\begin{solution}
\textbf{a.}~$u_n=3n+1$ \quad \textbf{b.}~$u_n=(-1)^{n+1}\cdot\dfrac{n}{n+2}$
\end{solution}
\begin{fullsolution}
\textbf{a.}~Telkens 3 meer: vergelijk met $3n$. Elke term is 1 meer, dus $u_n=3n+1$.\\
\textbf{b.}~Positief bij oneven $n$, teller $n$, noemer $n+2$.
\end{fullsolution}
\end{oefening}

\begin{oefening}[1]\label{oef:h1-t4}
Gegeven: $u_n = n^2 - n$. Is 56 een term van de rij? En 60?
\begin{solution}
56 is de 8ste term; 60 is geen term.
\end{solution}
\begin{fullsolution}
$n^2-n-56=0$ geeft $(n-8)(n+7)=0$, dus $n=8$. $n^2-n-60=0$ geeft $n=\dfrac{1\pm\sqrt{241}}{2}$, dus $n\approx 8{,}26$ of $n\approx -7{,}26$: geen volgnummer.
\end{fullsolution}
\end{oefening}

\begin{oefening}[1]\label{oef:h1-t5}
Gegeven: $u_1 = 4$, $u_2 = 1$ en $u_n = u_{n-1} + 2u_{n-2}$ voor $n \geqslant 3$. Bepaal de eerste vijf termen.
\begin{solution}
4, 1, 9, 11, 29
\end{solution}
\begin{fullsolution}
$u_3=1+8=9$, $u_4=9+2=11$, $u_5=11+18=29$.
\end{fullsolution}
\end{oefening}

\begin{oefening}[2]\label{oef:h1-t6}
Bepaal een recursief voorschrift van $5,\ 6,\ 8,\ 11,\ 15,\ \ldots$
\begin{solution}
$u_n=u_{n-1}+(n-1)$, $u_1=5$
\end{solution}
\begin{fullsolution}
Er komt 1, 2, 3, 4 bij: bij de stap naar $u_n$ is dat $n-1$. Dus $u_n=u_{n-1}+(n-1)$ met $u_1=5$.
\end{fullsolution}
\end{oefening}

\begin{oefening}[2]\label{oef:h1-t7}
Gegeven: $u_n = n^2 + 2n$. Bepaal een recursief voorschrift.
\begin{solution}
$u_n=u_{n-1}+2n+1$, $u_1=3$
\end{solution}
\begin{fullsolution}
$u_{n-1}=(n-1)^2+2(n-1)=n^2-1$, dus $u_n-u_{n-1}=2n+1$. Verder is $u_1=3$. Controle: $u_2=3+5=8=4+4$.
\end{fullsolution}
\end{oefening}

\begin{oefening}[1]\label{oef:h1-t8}
Teken de grafiek van $u_n = (-1)^n \cdot 2n$ voor $n = 1$ tot en met 5. Waarom verbind je de punten niet?
\begin{solution}
$(1,\ {-2})$, $(2,\ 4)$, $(3,\ {-6})$, $(4,\ 8)$, $(5,\ {-10})$; tussen de volgnummers is er geen term.
\end{solution}
\begin{fullsolution}
De termen zijn $-2,\ 4,\ {-6},\ 8,\ {-10}$. Zet de punten los uit. Een rij heeft enkel termen bij natuurlijke volgnummers: tussen $n=1$ en $n=2$ is er niets.
\end{fullsolution}
\end{oefening}

\begin{oefening}[2]\label{oef:h1-t9}
Gegeven is de rij $20,\ 17,\ 14,\ \ldots$ Geef een expliciet voorschrift als de eerste term $u_1$ heet, en als hij $u_0$ heet.
\begin{solution}
$u_n=23-3n$; $u_n=20-3n$
\end{solution}
\begin{fullsolution}
Telkens 3 minder. Met $u_1=20$: $u_n=23-3n$. Met $u_0=20$: $u_n=20-3n$.
\end{fullsolution}
\end{oefening}
\end{oefeningen}

\herhaalschema{%
  \fv{oef:h1-t1}{vb:h1-maanden}{oef:h1-termen}%
  \fv{oef:h1-t2}{vb:h1-termen}{oef:h1-expl}%
  \fv{oef:h1-t3}{vb:h1-opstellen}{oef:h1-opstellen}%
  \fv{oef:h1-t4}{vb:h1-isterm}{oef:h1-isterm}%
  \fv{oef:h1-t5}{vb:h1-twee}{oef:h1-recterm}%
  \fv{oef:h1-t6}{vb:h1-recopst}{oef:h1-recopst}%
  \fv{oef:h1-t7}{vb:h1-omzet}{oef:h1-omzet}%
  \fv{oef:h1-t8}{vb:h1-grafiek}{oef:h1-grafiek}%
  \fv{oef:h1-t9}{vb:h1-nul}{oef:h1-nuleen}}

\ifzelfstudie
\section*{Maak je eigen samenvatting}
Sluit dit boek en schrijf in je schrift: wat is een term, een volgnummer en de algemene term, wat is het verschil tussen een expliciet en een recursief voorschrift, hoe stel je ze op, hoe test je of een getal een term is, wanneer start je bij $u_0$, hoe ziet de grafiek van een rij eruit, en welke fouten wil je vermijden? Vergelijk daarna met de kaders.
\fi

%----------------------------------------------------------------------------------------
% REKENKUNDIGE EN MEETKUNDIGE RIJEN. Alle opgaven zijn eigen werk.
%
% Opbouw: Start, leerstof (kern, uitgewerkt voorbeeld, korte oefening), en daarna op een
% nieuwe pagina de Opdrachten (het grootste deel, door elkaar), Koppel en Test jezelf.
%
%   \voorbeeld{titel}{opgave}{uitwerking}   uitgewerkt voorbeeld, in alle versies
%   \begin{oefeningen}                      oefeningen zonder titel, tussen de leerstof
%   \begin{oefeningen}[Titel]               oefeningen met een strook (Start)
%   \begin{oefening}[2]                     niveau 1, 2 of 3: balkjes bij de Opdrachten
%   \thema{...}                             onderwerp, enkel in de leerkrachtenversie
%   \koppel{rubriek}{leerstof}              label naar andere leerstof, in de kleur van de rubriek
%   \opdrachtenkop                          start van de Opdrachten: nieuwe pagina, felle band en tab
%   \opdrachtsectie{Titel}                  felle balk voor Koppel en Test jezelf
%   \lkr{...}                               enkel in de leerkrachtenversie
%   \thuis{...}, \nujij{opgave}{antwoord}   enkel in de zelfstudieversie
%----------------------------------------------------------------------------------------

\mychapter{2}{Rekenkundige en meetkundige rijen}

\lkr{Leerplan B+S'': LPD 44, zonder de somformules (die staan in hoofdstuk 3). De voorbeelden staan uitgewerkt in het handboek; de oefeningen noteren leerlingen in hun schrift. Twee rode draden: de sprongenfiguur ($n-1$ sprongen, zowel bij $+v$ als bij $\cdot\, q$) en het onderscheid tussen rekenkundig en meetkundig, dat in de gemengde Opdrachten steeds terugkomt.}

\begin{oefeningen}[Start]
\begin{oefening}
\blokkop{Opwarmer}
\begin{enumerate}
  \item $u_n = 3n - 1$. Bereken $u_5$.
  \item $u_1 = 2$ en $u_n = u_{n-1} + 5$. Bereken $u_4$.
  \item Los op: $3n - 7 = 50$.
  \item Los op: $a + 2b = 11$ en $a + 8b = 35$.
  \item Bepaal de rico van de rechte door $(1, 2)$ en $(4, 11)$.
  \item Los op: $2^x = 32$ en $3^x = 50$.
\end{enumerate}
\begin{solution}
\textbf{a.}~14 \quad \textbf{b.}~17 \quad \textbf{c.}~$n=19$ \quad \textbf{d.}~$a=3$, $b=4$ \quad \textbf{e.}~3 \quad \textbf{f.}~$x=5$; $x\approx 3{,}56$
\end{solution}
\begin{fullsolution}
\textbf{a.}~$u_5=3\cdot 5-1=14$.\\
\textbf{b.}~$u_2=7$, $u_3=12$, $u_4=17$.\\
\textbf{c.}~$3n=57$, dus $n=19$.\\
\textbf{d.}~Aftrekken: $6b=24$, dus $b=4$ en $a=3$.\\
\textbf{e.}~$\dfrac{11-2}{4-1}=3$.\\
\textbf{f.}~$2^5=32$, dus $x=5$. $x=\log_3 50=\dfrac{\log 50}{\log 3}\approx 3{,}56$.
\end{fullsolution}
\end{oefening}

\begin{oefening}\label{oef:verken}
\blokkop{Verken}\quad Eén plastic stoel is 80 cm hoog. Elke stoel die je erbovenop stapelt, maakt de stapel 5 cm hoger.
\begin{enumerate}
  \item Hoe hoog is een stapel van 2, 3 en 4 stoelen?
  \item Hoe hoog is een stapel van 25 stoelen?
  \item Hoeveel stoelen mogen er maximaal op een stapel om onder een deur van 2,12~m door te kunnen?
\end{enumerate}
\begin{solution}
\textbf{a.}~85, 90 en 95 cm \quad \textbf{b.}~200 cm \quad \textbf{c.}~27
\end{solution}
\begin{fullsolution}
\textbf{a.}~80 cm is de eerste stoel, dus 85, 90 en 95 cm.\\
\textbf{b.}~$80+24\cdot 5=200$ cm: de eerste stoel is geen sprong.\\
\textbf{c.}~$80+(n-1)\cdot 5\leqslant 212$ geeft $n-1\leqslant 26{,}4$, dus maximaal 27 stoelen.
\end{fullsolution}
\end{oefening}
\end{oefeningen}

\lkr{Laat oefening \ref{oef:verken} maken vóór de uitleg, zonder formule. Vergelijk daarna de antwoorden op b: wie rekent $25\cdot 5$ en wie $24\cdot 5$, en waarom?}

%========================================================================================
\section{Rekenkundige rijen}

\subsection{Definitie en voorschriften}

\begin{theorie}
\begin{minipage}[t]{.5\linewidth}
Elke term is de vorige term plus hetzelfde getal: het \emph{verschil} $v$.

\vspace{1mm}
Recursief: \ $\formule{u_n = u_{n-1} + v}$, met $u_1$ gegeven\\[1mm]
Expliciet: \ \,$\formule{u_n = u_1 + (n-1)\,v}$
\end{minipage}\hfill
\begin{minipage}[t]{.47\linewidth}
\vspace{-3mm}
\sprongfiguur{+v}
\end{minipage}
\end{theorie}

\thuis{Bij de stoelen uit oefening \ref{oef:verken} is $u_1 = 80$ en $v = 5$. Voor 25 stoelen tel je geen 25 keer 5 cm op, want de eerste stoel is geen sprong. Van $u_1$ naar $u_n$ zijn het dus altijd $n-1$ sprongen, en dat geeft het expliciet voorschrift.}

\voorbeeld{}{Bepaal het expliciet voorschrift van $7,\ 11,\ 15,\ \ldots$ en bereken $u_{30}$.\label{vb:term}}{%
\begin{stappen}
v = 11 - 7 = 4 & Verschil bepalen. Controle: $15 - 11 = 4$. \\
u_n = 7 + (n-1) \cdot 4 = 4n + 3 & Invullen in het voorschrift en vereenvoudigen. \\
u_{30} = 4 \cdot 30 + 3 = 123 & $n = 30$ invullen.
\end{stappen}}

\lkr{Laat leerlingen $u_{30}$ ook narekenen met $7 + 29 \cdot 4$: zo controleren ze het vereenvoudigde voorschrift.}

\nujij{Bepaal het expliciet voorschrift van $12,\ 9,\ 6,\ \ldots$ en bereken $u_{25}$.}{$v = -3$, $u_n = 15 - 3n$, $u_{25} = -60$.}

\thuis{\blokkop{Half uitgewerkt}\quad Rij $3;\ 3{,}5;\ 4;\ \ldots$ Bereken $u_{41}$. Schrijf over in je schrift en vul aan: \ $v = \vak$, \ $u_{41} = 3 + \vak \cdot \vak = \vak$. \quad {\footnotesize(Antwoord: $v = 0{,}5$ en $u_{41} = 3 + 40 \cdot 0{,}5 = 23$.)}}

\begin{oefeningen}
\begin{oefening}
Bepaal het expliciet voorschrift van $2,\ 9,\ 16,\ \ldots$ en bereken $u_{20}$.
\begin{solution}
$u_n=7n-5$; $u_{20}=135$
\end{solution}
\begin{fullsolution}
$v=7$, dus $u_n=2+(n-1)\cdot 7=7n-5$ en $u_{20}=7\cdot 20-5=135$.
\end{fullsolution}
\end{oefening}
\end{oefeningen}

\subsection{Term of geen term}

\begin{werkwijze}
Los $u_n = x$ op. Is $n$ een natuurlijk getal (minstens 1), dan is $x$ de $n$-de term. Anders is $x$ geen term.
\end{werkwijze}

\voorbeeld{}{Is 203 een term van $7,\ 11,\ 15,\ \ldots$? En 200?\label{vb:isterm}}{%
\begin{stappen}
7 + (n-1) \cdot 4 = 203 \ \Rightarrow \ n = 50 & Natuurlijk getal: \textbf{ja}, 203 is de 50ste term. \\
7 + (n-1) \cdot 4 = 200 \ \Rightarrow \ n = 49{,}25 & Geen natuurlijk getal: \textbf{nee}.
\end{stappen}}

\nujij{Is 10 een term van $100,\ 94,\ 88,\ \ldots$? En 1?}{10 is de 16de term; bij 1 is $n = 17{,}5$, dus geen term.}

\begin{oefeningen}
\begin{oefening}
Is 88 een term van $4,\ 11,\ 18,\ \ldots$? En 90?
\begin{solution}
88 is de 13de term; 90 is geen term.
\end{solution}
\begin{fullsolution}
$4+(n-1)\cdot 7=88$ geeft $n-1=12$, dus $n=13$. $4+(n-1)\cdot 7=90$ geeft $n-1=\tfrac{86}{7}$: geen natuurlijk getal.
\end{fullsolution}
\end{oefening}
\end{oefeningen}

\subsection{De rij vinden uit twee termen}

\begin{theorie}
Tel de sprongen: \quad $\formule{u_n = u_m + (n-m)\,v}$
\end{theorie}

\voorbeeld{}{$u_3 = 11$ en $u_9 = 35$. Bepaal $u_1$ en $v$ op twee manieren: met een stelsel en met sprongen.\label{vb:tweetermen}}{%
\begin{minipage}[t]{.47\linewidth}
\textbf{Stelsel.} \ $u_1 + 2v = 11$ en $u_1 + 8v = 35$.\\
Tweede min eerste: $6v = 24$, dus $v = 4$.\\
Invullen: $u_1 = 11 - 2 \cdot 4 = 3$.
\end{minipage}\hfill
\begin{minipage}[t]{.47\linewidth}
\textbf{Sprongen.} \ Van $u_3$ naar $u_9$: 6 sprongen.\\
$35 - 11 = 24 = 6v$, dus $v = 4$.\\
Twee sprongen terug: $u_1 = 11 - 2 \cdot 4 = 3$.
\end{minipage}}

\lkr{Laat leerlingen de twee methodes vergelijken: waar zit dezelfde berekening?}

\nujij{$u_4 = -2$ en $u_{10} = 16$. Bepaal $u_1$ en $v$.}{$v = 3$ en $u_1 = -11$.}

\begin{oefeningen}
\begin{oefening}
$u_3 = 10$ en $u_8 = 35$. Bepaal $u_1$ en $v$.
\begin{solution}
$v=5$, $u_1=0$
\end{solution}
\begin{fullsolution}
5 sprongen: $5v=25$, dus $v=5$. Twee sprongen terug: $u_1=10-2\cdot 5=0$.
\end{fullsolution}
\end{oefening}
\end{oefeningen}

\subsection{Rekenkundig midden}

\begin{theorie}
Elke term (behalve de eerste) is het gemiddelde van de twee termen ernaast:

\qquad $\formule{u_n = \dfrac{u_{n-1} + u_{n+1}}{2}}$

\begin{bewijs}
\begin{stappen}
u_{n-1} = u_n - v \quad\text{en}\quad u_{n+1} = u_n + v & Een sprong terug en een sprong vooruit. \\
u_{n-1} + u_{n+1} = 2u_n & Optellen: $-v$ en $+v$ vallen weg. \\[1.5ex]
\dfrac{u_{n-1} + u_{n+1}}{2} = u_n & Delen door 2.
\end{stappen}
\end{bewijs}
\end{theorie}

\begin{oefeningen}
\begin{oefening}
Bereken de ontbrekende term.
\begin{enumerate}
  \item $u_7$ als $u_6 = 12$ en $u_8 = 20$
  \item $u_{12}$ als $u_{11} = -3$ en $u_{13} = 5$
\end{enumerate}
\begin{solution}
\textbf{a.}~16 \quad \textbf{b.}~1
\end{solution}
\begin{fullsolution}
\textbf{a.}~$\dfrac{12+20}{2}=16$.\\
\textbf{b.}~$\dfrac{-3+5}{2}=1$.
\end{fullsolution}
\end{oefening}
\end{oefeningen}

\subsection{Grafiek en lineaire groei}

\begin{theorie}
\begin{minipage}[t]{.64\linewidth}
\textbf{Grafiek.} De punten $(n, u_n)$ liggen op een rechte met rico $v$. Enkel de punten horen bij de rij: $n = 2{,}5$ bestaat niet.

\vspace{1.5mm}
\textbf{Lineaire groei.} Stijgt of daalt een grootheid per tijdseenheid met hetzelfde bedrag en is $u_0$ de beginwaarde, dan is na $n$ tijdseenheden: \quad $\formule{u_n = u_0 + n\,v}$
\end{minipage}\hfill
\begin{minipage}[t]{.32\linewidth}
\vspace{-2mm}
\begin{tikzpicture}[x=.74cm,y=.21cm,>={Stealth[length=1.6mm]},font=\scriptsize]
  \draw[->,lijngrijs] (-.1,0) -- (5.7,0) node[right,text=black] {$n$};
  \draw[->,lijngrijs] (0,-.4) -- (0,10.8) node[right,text=black] {$u_n$};
  \foreach \x in {1,...,5} \draw[lijngrijs] (\x,.2) -- (\x,-.2) node[below=-1pt,text=grijs] {\x};
  \draw[lijngrijs,dashed] (.5,0) -- (5.5,10);
  \draw[accent,thick] (3,5) -- node[below=-1pt,text=black] {$+1$} (4,5) -- node[right=-1pt,text=black] {$+v$} (4,7);
  \foreach \x/\y in {1/1,2/3,3/5,4/7,5/9} \fill[accent] (\x,\y) circle (1.9pt);
\end{tikzpicture}
\end{minipage}
\end{theorie}

\voorbeeld{}{Er staat \eur 350 op een rekening. Elke maand komt er \eur 25 bij. Na hoeveel maanden staat er minstens \eur 1000 op?\label{vb:groei}}{%
\begin{stappen}
u_n = 350 + 25n & Beginwaarde $u_0$, want $n$ telt de maanden \emph{na} de start. \\
350 + 25n \geqslant 1000 \ \Rightarrow \ n \geqslant 26 & Na \textbf{26 maanden}.
\end{stappen}}

\lkr{Vergelijk met de stoelen: daar start de rij bij $u_1$, hier bij $u_0$. Laat leerlingen uitleggen waarom het voorschrift dan $n$ bevat en niet $n-1$.}

\nujij{Een waterton bevat 1200 liter en verliest 15 liter per uur. Geef een voorschrift voor de inhoud na $n$ uur. Na hoeveel uur is de ton leeg?}{$u_n = 1200 - 15n$; leeg na 80 uur.}

\begin{oefeningen}
\begin{oefening}
Een zwembad bevat 120 m\textsuperscript{3} water en wordt bijgevuld met 1,5 m\textsuperscript{3} per uur. Geef een voorschrift. Wanneer zit er 180 m\textsuperscript{3} in?
\begin{solution}
$u_n=120+1{,}5\,n$; na 40 uur
\end{solution}
\begin{fullsolution}
$u_n=120+1{,}5\,n=180$ geeft $n=40$.
\end{fullsolution}
\end{oefening}
\end{oefeningen}

%========================================================================================
\section{Meetkundige rijen}

\begin{oefeningen}
\begin{oefening}\label{oef:mverken}
\blokkop{Verken}\quad Op dag 1 kent één persoon een gerucht. Elke dag wordt het aantal mensen dat het kent drie keer zo groot.
\begin{enumerate}
  \item Hoeveel mensen kennen het op dag 2, 3 en 4?
  \item Hoeveel mensen kennen het op dag 10?
  \item Op welke dag kennen voor het eerst meer dan een miljoen mensen het gerucht?
\end{enumerate}
\begin{solution}
\textbf{a.}~3, 9, 27 \quad \textbf{b.}~19\,683 \quad \textbf{c.}~dag 14
\end{solution}
\begin{fullsolution}
\textbf{a.}~Telkens maal 3: 3, 9 en 27.\\
\textbf{b.}~Van dag 1 naar dag 10 zijn het 9 sprongen: $3^9=19\,683$.\\
\textbf{c.}~$3^{n-1}>1\,000\,000$ geeft $n-1>12{,}58$, dus $n-1=13$ en $n=14$.
\end{fullsolution}
\end{oefening}
\end{oefeningen}

\lkr{Laat de leerlingen de vergelijking maken met de stoelen: wat blijft hetzelfde (het aantal sprongen) en wat verandert (optellen wordt vermenigvuldigen)?}

\subsection{Definitie en voorschriften}

\begin{theorie}
\begin{minipage}[t]{.5\linewidth}
Elke term is de vorige term maal hetzelfde getal: het \emph{quotiënt} $q$.

\vspace{1mm}
Recursief: \ $\formule{u_n = u_{n-1} \cdot q}$, met $u_1$ gegeven\\[1mm]
Expliciet: \ \,$\formule{u_n = u_1 \cdot q^{\,n-1}}$
\end{minipage}\hfill
\begin{minipage}[t]{.47\linewidth}
\vspace{-3mm}
\sprongfiguur{\cdot\, q}
\end{minipage}
\end{theorie}

\voorbeeld{}{Bepaal het expliciet voorschrift van $3,\ 6,\ 12,\ \ldots$ en bereken $u_{10}$.\label{vb:mterm}}{%
\begin{stappen}
q = \dfrac{6}{3} = 2 & Quotiënt bepalen. Controle: $\dfrac{12}{6} = 2$. \\[1.5ex]
u_n = 3 \cdot 2^{\,n-1} & Invullen in het voorschrift. Dit mag je niet verder samenvoegen tot $6^{\,n-1}$. \\
u_{10} = 3 \cdot 2^{9} = 1536 & $n = 10$ invullen: 9 sprongen.
\end{stappen}}

\nujij{Bepaal het expliciet voorschrift van $4,\ {-8},\ 16,\ \ldots$ en bereken $u_7$.}{$q = -2$, $u_n = 4 \cdot (-2)^{n-1}$, $u_7 = 256$.}

\begin{oefeningen}
\begin{oefening}
Bepaal het expliciet voorschrift van $5,\ 15,\ 45,\ \ldots$ en bereken $u_6$.
\begin{solution}
$u_n=5\cdot 3^{n-1}$; $u_6=1215$
\end{solution}
\begin{fullsolution}
$q=3$, dus $u_n=5\cdot 3^{n-1}$ en $u_6=5\cdot 3^5=1215$.
\end{fullsolution}
\end{oefening}
\end{oefeningen}

\subsection{Term of geen term}

\begin{werkwijze}
Los $u_n = x$ op, zo nodig met een logaritme. Is $n$ een natuurlijk getal (minstens 1), dan is $x$ de $n$-de term.
\end{werkwijze}

\voorbeeld{}{Is 768 een term van $3,\ 6,\ 12,\ \ldots$? En 1000?\label{vb:misterm}}{%
\begin{stappen}
3 \cdot 2^{\,n-1} = 768 \ \Rightarrow \ 2^{\,n-1} = 256 & Delen door 3. \\
n - 1 = 8 \ \Rightarrow \ n = 9 & Want $2^8 = 256$. \textbf{Ja}, 768 is de 9de term. \\
3 \cdot 2^{\,n-1} = 1000 \ \Rightarrow \ 2^{\,n-1} = 333{,}33\ldots & Hetzelfde voor 1000. \\
n - 1 = \log_2 333{,}33\ldots \approx 8{,}38 & Geen natuurlijk getal: \textbf{nee}.
\end{stappen}}

\nujij{Is 405 een term van $5,\ 15,\ 45,\ \ldots$? En 500?}{405 is de 5de term; bij 500 is $3^{n-1} = 100$ en $n - 1 \approx 4{,}19$, dus geen term.}

\begin{oefeningen}
\begin{oefening}
Is 486 een term van $2,\ 6,\ 18,\ \ldots$? En 500?
\begin{solution}
486 is de 6de term; 500 is geen term.
\end{solution}
\begin{fullsolution}
$2\cdot 3^{n-1}=486$ geeft $3^{n-1}=243=3^5$, dus $n=6$. $2\cdot 3^{n-1}=500$ geeft $3^{n-1}=250$ en $n-1\approx 5{,}03$: geen natuurlijk getal.
\end{fullsolution}
\end{oefening}
\end{oefeningen}

\subsection{De rij vinden uit twee termen}

\begin{theorie}
Tel de sprongen: \quad $\formule{u_n = u_m \cdot q^{\,n-m}}$ \quad Let op: bij een even aantal sprongen zijn er twee mogelijke waarden voor $q$.
\end{theorie}

\voorbeeld{}{$u_3 = 12$ en $u_5 = 48$. Bepaal $u_1$ en $q$.\label{vb:mtwee}}{%
\begin{stappen}
12 \cdot q^2 = 48 \ \Rightarrow \ q^2 = 4 & Van $u_3$ naar $u_5$: 2 sprongen. \\
q = 2 \quad\text{of}\quad q = -2 & Even aantal sprongen: \textbf{twee oplossingen}. \\[1.5ex]
u_1 = \dfrac{12}{q^2} = 3 & Twee sprongen terug: in beide gevallen $u_1 = 3$. \\[1.5ex]
3,\ 6,\ 12,\ 24,\ 48 \quad\text{of}\quad 3,\ {-6},\ 12,\ {-24},\ 48 & Controle: beide rijen voldoen.
\end{stappen}}

\nujij{$u_2 = 4$ en $u_4 = 36$. Bepaal $u_1$ en $q$.}{$q = 3$ en $u_1 = \tfrac43$, of $q = -3$ en $u_1 = -\tfrac43$.}

\begin{oefeningen}
\begin{oefening}
$u_2 = 10$ en $u_5 = 80$. Bepaal $u_1$ en $q$.
\begin{solution}
$q=2$, $u_1=5$
\end{solution}
\begin{fullsolution}
3 sprongen: $10\cdot q^3=80$, dus $q^3=8$ en $q=2$. Een sprong terug: $u_1=\tfrac{10}{2}=5$.
\end{fullsolution}
\end{oefening}
\end{oefeningen}

\subsection{Meetkundig midden}

\begin{theorie}
Het kwadraat van elke term (behalve de eerste) is het product van de twee termen ernaast: \quad $\formule{u_n^{\,2} = u_{n-1} \cdot u_{n+1}}$

\begin{bewijs}
\begin{stappen}
u_{n-1} = \dfrac{u_n}{q} \quad\text{en}\quad u_{n+1} = u_n \cdot q & Een sprong terug en een sprong vooruit. \\[1.5ex]
u_{n-1} \cdot u_{n+1} = \dfrac{u_n}{q} \cdot u_n \cdot q = u_n^{\,2} & Vermenigvuldigen: $q$ valt weg.
\end{stappen}
\end{bewijs}
\end{theorie}

\begin{oefeningen}
\begin{oefening}
$u_4 = 3$ en $u_6 = 12$. Welke waarden kan $u_5$ hebben?
\begin{solution}
6 of $-6$
\end{solution}
\begin{fullsolution}
$u_5^{\,2}=3\cdot 12=36$, dus $u_5=6$ of $u_5=-6$.
\end{fullsolution}
\end{oefening}
\end{oefeningen}

\subsection{Grafiek en exponentiële groei}

\begin{theorie}
\begin{minipage}[t]{.5\linewidth}
\textbf{Grafiek.} Bij $q > 0$ (en $q \neq 1$) liggen de punten $(n, u_n)$ op een exponentiële kromme. Bij $q < 0$ wisselen de termen telkens van teken.

\vspace{1.5mm}
\textbf{Exponentiële groei.} Bij een groei van $p\%$ per tijdseenheid is $\formule{q = 1 + \tfrac{p}{100}}$, bij een daling $\formule{q = 1 - \tfrac{p}{100}}$. Met beginwaarde $u_0$ is dan: \quad $\formule{u_n = u_0 \cdot q^{\,n}}$
\end{minipage}\hfill
\begin{minipage}[t]{.47\linewidth}
\vspace{-2mm}
\begin{tikzpicture}[x=.36cm,y=.17cm,>={Stealth[length=1.4mm]},font=\scriptsize]
  \draw[->,lijngrijs] (-.2,0) -- (5.9,0) node[right,text=black] {$n$};
  \draw[->,lijngrijs] (0,-.5) -- (0,17.5) node[above,text=black] {$u_n$};
  \draw[lijngrijs,dashed,domain=.3:5.2,samples=30] plot (\x,{2^(\x-1)});
  \foreach \x/\y in {1/1,2/2,3/4,4/8,5/16} \fill[accent] (\x,\y) circle (1.7pt);
  \node[text=grijs,anchor=north] at (2.9,-.8) {$q = 2$};
\end{tikzpicture}\hspace{3mm}%
\begin{tikzpicture}[x=.36cm,y=.17cm,>={Stealth[length=1.4mm]},font=\scriptsize]
  \draw[->,lijngrijs] (-.2,0) -- (5.9,0) node[right,text=black] {$n$};
  \draw[->,lijngrijs] (0,-5.5) -- (0,11.5) node[above,text=black] {$u_n$};
  \foreach \x/\y in {1/8,2/-4,3/2,4/-1,5/0.5} \fill[accent] (\x,\y) circle (1.7pt);
  \node[text=grijs,anchor=north] at (2.9,-6.3) {$q = -\tfrac12$};
\end{tikzpicture}
\end{minipage}
\end{theorie}

\voorbeeld{}{Een stad telt 8000 inwoners en groeit met 2\% per jaar. Hoeveel inwoners zijn er na 10 jaar? Wanneer zijn er 12\,000 inwoners?\label{vb:mgroei}}{%
\begin{stappen}
q = 1 + \dfrac{2}{100} = 1{,}02 & Groei van 2\%. \\[1.5ex]
u_n = 8000 \cdot 1{,}02^{\,n} & Beginwaarde $u_0$, want $n$ telt de jaren \emph{na} de start. \\
u_{10} = 8000 \cdot 1{,}02^{10} \approx 9752 & Na 10 jaar: ongeveer \textbf{9752 inwoners}. \\
1{,}02^{\,n} = 1{,}5 \ \Rightarrow \ n = \log_{1{,}02} 1{,}5 \approx 20{,}5 & Delen door 8000. In het \textbf{21ste jaar}.
\end{stappen}}

\nujij{Een stad telt 50\,000 inwoners en krimpt met 1,5\% per jaar. Hoeveel inwoners zijn er na 10 jaar?}{$u_{10} = 50\,000 \cdot 0{,}985^{10} \approx 42\,987$.}

\begin{oefeningen}
\begin{oefening}
Een auto kost \eur 20\,000 en verliest elk jaar 15\% van zijn waarde. Geef een voorschrift. Wat is de waarde na 5 jaar?
\begin{solution}
$u_n=20\,000\cdot 0{,}85^{\,n}$; ongeveer \eur 8874
\end{solution}
\begin{fullsolution}
$q=1-\tfrac{15}{100}=0{,}85$, dus $u_n=20\,000\cdot 0{,}85^{\,n}$ en $u_5\approx 8874$.
\end{fullsolution}
\end{oefening}
\end{oefeningen}

%========================================================================================
\section{Enkelvoudige en samengestelde intrest}

\begin{theorie}
Je belegt een kapitaal $K_0$ aan $p\%$ per jaar. $K_n$ is het kapitaal na $n$ jaar.

\vspace{1mm}
\begin{tabular}{@{}>{\RaggedRight\arraybackslash}p{7cm}>{\RaggedRight\arraybackslash}p{7.6cm}@{}}
\textbf{Enkelvoudig:} intrest enkel op het beginkapitaal & \textbf{Samengesteld:} ook intrest op de intrest \\[1mm]
$\formule{K_n = K_0 \left(1 + n \cdot \dfrac{p}{100}\right)}$ & $\formule{K_n = K_0 \left(1 + \dfrac{p}{100}\right)^{n}}$ \\[3mm]
Rekenkundige rij: $\formule{v = K_0 \cdot \dfrac{p}{100}}$ & Meetkundige rij: $\formule{q = 1 + \dfrac{p}{100}}$
\end{tabular}
\end{theorie}

\voorbeeld{}{Je belegt \eur 2000 aan 3\% per jaar. Hoeveel heb je na 10 jaar, met enkelvoudige en met samengestelde intrest?\label{vb:intrest}}{%
\begin{stappen}
K_{10} = 2000 \,(1 + 10 \cdot 0{,}03) = 2600 & Enkelvoudig: elk jaar \eur 60 erbij. \\
K_{10} = 2000 \cdot 1{,}03^{10} \approx 2687{,}83 & Samengesteld: \eur 87,83 meer, door intrest op intrest.
\end{stappen}}

\nujij{Je belegt \eur 1500 aan 6\% per jaar. Hoeveel heb je na 4 jaar, met enkelvoudige en met samengestelde intrest?}{Enkelvoudig \eur 1860; samengesteld ongeveer \eur 1893,72.}

\begin{oefeningen}
\begin{oefening}
Je belegt \eur 5000 aan 4\% per jaar. Hoeveel heb je na 6 jaar, met enkelvoudige en met samengestelde intrest?
\begin{solution}
Enkelvoudig \eur 6200; samengesteld ongeveer \eur 6326,60
\end{solution}
\begin{fullsolution}
Enkelvoudig: $5000\,(1+6\cdot 0{,}04)=6200$. Samengesteld: $5000\cdot 1{,}04^6\approx 6326{,}60$.
\end{fullsolution}
\end{oefening}
\end{oefeningen}

%----------------------------------------------------------------------------------------
% OPDRACHTEN: het grootste deel van de oefeningen, op een nieuwe pagina
%----------------------------------------------------------------------------------------
\opdrachtenkop

\lkr{De opdrachten staan bewust door elkaar, met stijgende moeilijkheid. In deze versie staat bij elke opdracht het onderwerp, zodat je kunt kiezen. Geef ze pas volledig op als de hele leerstof gezien is, of selecteer per onderwerp.}

\begin{oefeningen}
\begin{oefening}[1]\label{oef:soort}
\thema{gemengd}Is de rij rekenkundig, meetkundig, beide of geen van beide?
\begin{reeks}
  \item $7,\ 14,\ 21,\ \ldots$
  \item $7,\ 14,\ 28,\ \ldots$
  \item $1,\ 1,\ 1,\ \ldots$
  \item $1,\ 2,\ 3,\ 5,\ \ldots$
  \item $9,\ 3,\ 1,\ \ldots$
  \item $9,\ 3,\ {-3},\ \ldots$
\end{reeks}
\begin{solution}
\textbf{a.}~rekenkundig \quad \textbf{b.}~meetkundig \quad \textbf{c.}~beide \quad \textbf{d.}~geen van beide \quad \textbf{e.}~meetkundig \quad \textbf{f.}~rekenkundig
\end{solution}
\begin{fullsolution}
\textbf{a.}~$v=7$.\\
\textbf{b.}~$q=2$.\\
\textbf{c.}~$v=0$ en $q=1$.\\
\textbf{d.}~Verschillen 1, 1, 2 en quotiënten 2, $\tfrac32$, $\tfrac53$: geen van beide.\\
\textbf{e.}~$q=\tfrac13$.\\
\textbf{f.}~$v=-6$.
\end{fullsolution}
\end{oefening}

\begin{oefening}[1]
\thema{rekenkundig: voorschriften}Bereken de gevraagde term.
\begin{enumerate}
  \item $u_{15}$ als $u_1 = -8$ en $v = 6$
  \item $u_{21}$ als $u_1 = 40$ en $v = -2{,}5$
  \item $u_{10}$ als $u_1 = \tfrac13$ en $v = \tfrac23$
\end{enumerate}
\begin{solution}
\textbf{a.}~76 \quad \textbf{b.}~$-10$ \quad \textbf{c.}~$\tfrac{19}{3}$
\end{solution}
\begin{fullsolution}
\textbf{a.}~$u_{15}=-8+14\cdot 6=76$.\\
\textbf{b.}~$u_{21}=40+20\cdot(-2{,}5)=-10$.\\
\textbf{c.}~$u_{10}=\tfrac13+9\cdot\tfrac23=\tfrac{19}{3}$.
\end{fullsolution}
\end{oefening}

\begin{oefening}[1]\label{oef:isterm}
\thema{rekenkundig: term of geen term}Is 370 een term van $6,\ 13,\ 20,\ \ldots$? En 500?
\begin{solution}
370 is de 53ste term; 500 is geen term.
\end{solution}
\begin{fullsolution}
$6+(n-1)\cdot 7=370$ geeft $n-1=52$, dus $n=53$. $6+(n-1)\cdot 7=500$ geeft $n-1=\tfrac{494}{7}$: geen natuurlijk getal.
\end{fullsolution}
\end{oefening}

\begin{oefening}[1]\label{oef:mvoorschrift}
\thema{meetkundig: voorschriften}Geef een recursief en een expliciet voorschrift.
\begin{reeks}
  \item $4,\ 12,\ 36,\ \ldots$
  \item $80,\ 40,\ 20,\ \ldots$
  \item $-1,\ 2,\ {-4},\ \ldots$
\end{reeks}
\begin{solution}
\textbf{a.}~$u_n=3u_{n-1}$, $u_1=4$; $u_n=4\cdot 3^{n-1}$ \quad \textbf{b.}~$u_n=\tfrac12 u_{n-1}$, $u_1=80$; $u_n=80\cdot(\tfrac12)^{n-1}$ \quad \textbf{c.}~$u_n=-2u_{n-1}$, $u_1=-1$; $u_n=-(-2)^{n-1}$
\end{solution}
\begin{fullsolution}
\textbf{a.}~$q=3$: $u_n=4\cdot 3^{n-1}$.\\
\textbf{b.}~$q=\tfrac12$: $u_n=80\cdot(\tfrac12)^{n-1}$.\\
\textbf{c.}~$q=-2$: $u_n=-1\cdot(-2)^{n-1}=-(-2)^{n-1}$.\\
Het recursief voorschrift is telkens $u_n=u_{n-1}\cdot q$ met de eerste term erbij.
\end{fullsolution}
\end{oefening}

\begin{oefening}[1]
\thema{gemengd}De eerste twee termen van een rij zijn 2 en 6. Wat is de derde term als de rij rekenkundig is? En als ze meetkundig is?
\begin{solution}
10; 18
\end{solution}
\begin{fullsolution}
Rekenkundig: $v=4$, dus $6+4=10$. Meetkundig: $q=3$, dus $6\cdot 3=18$.
\end{fullsolution}
\end{oefening}

\begin{oefening}[1]\label{oef:tweetermen}
\thema{rekenkundig: twee termen}Bepaal $u_1$ en $v$.
\begin{enumerate}
  \item $u_2 = 9$ en $u_7 = 34$
  \item $u_5 = -1$ en $u_{12} = -22$
  \item $u_6 = 4$ en $u_{16} = 9$
\end{enumerate}
\begin{solution}
\textbf{a.}~$v=5$, $u_1=4$ \quad \textbf{b.}~$v=-3$, $u_1=11$ \quad \textbf{c.}~$v=\tfrac12$, $u_1=\tfrac32$
\end{solution}
\begin{fullsolution}
\textbf{a.}~5 sprongen: $5v=25$, dus $v=5$ en $u_1=9-5=4$.\\
\textbf{b.}~7 sprongen: $7v=-21$, dus $v=-3$ en $u_1=-1-4\cdot(-3)=11$.\\
\textbf{c.}~10 sprongen: $10v=5$, dus $v=\tfrac12$ en $u_1=4-5\cdot\tfrac12=\tfrac32$.
\end{fullsolution}
\end{oefening}

\begin{oefening}[1]
\thema{meetkundig: voorschriften}Bereken de gevraagde term.
\begin{enumerate}
  \item $u_8$ als $u_1 = 2$ en $q = 3$
  \item $u_7$ als $u_1 = 64$ en $q = \tfrac12$
  \item $u_6$ als $u_1 = -1$ en $q = -2$
\end{enumerate}
\begin{solution}
\textbf{a.}~4374 \quad \textbf{b.}~1 \quad \textbf{c.}~32
\end{solution}
\begin{fullsolution}
\textbf{a.}~$u_8=2\cdot 3^7=4374$.\\
\textbf{b.}~$u_7=64\cdot(\tfrac12)^6=1$.\\
\textbf{c.}~$u_6=-1\cdot(-2)^5=32$.
\end{fullsolution}
\end{oefening}

\begin{oefening}[1]
\thema{gemengd}Gegeven: $u_n = n^2 - 1$. Geef de eerste vier termen. Is de rij rekenkundig of meetkundig?
\begin{solution}
0, 3, 8, 15; geen van beide
\end{solution}
\begin{fullsolution}
$u_1=0$, $u_2=3$, $u_3=8$, $u_4=15$. De verschillen 3, 5, 7 zijn niet constant, en de quotiënten evenmin.
\end{fullsolution}
\end{oefening}

\begin{oefening}[1]\label{oef:misterm}
\thema{meetkundig: term of geen term}Is 1458 een term van $2,\ 6,\ 18,\ \ldots$? En 2000?
\begin{solution}
1458 is de 7de term; 2000 is geen term.
\end{solution}
\begin{fullsolution}
$2\cdot 3^{n-1}=1458$ geeft $3^{n-1}=729=3^6$, dus $n=7$. $2\cdot 3^{n-1}=2000$ geeft $3^{n-1}=1000$ en $n-1\approx 6{,}29$: geen natuurlijk getal.
\end{fullsolution}
\end{oefening}

\begin{oefening}[1]\label{oef:voorschrift}
\thema{rekenkundig: voorschriften}Geef een recursief en een expliciet voorschrift.
\begin{reeks}
  \item $5,\ 12,\ 19,\ \ldots$
  \item $20,\ 17,\ 14,\ \ldots$
  \item $\tfrac14,\ \tfrac34,\ \tfrac54,\ \ldots$
\end{reeks}
\begin{solution}
\textbf{a.}~$u_n=u_{n-1}+7$, $u_1=5$; $u_n=7n-2$ \quad \textbf{b.}~$u_n=u_{n-1}-3$, $u_1=20$; $u_n=23-3n$ \quad \textbf{c.}~$u_n=u_{n-1}+\tfrac12$, $u_1=\tfrac14$; $u_n=\tfrac12 n-\tfrac14$
\end{solution}
\begin{fullsolution}
\textbf{a.}~$v=7$: $u_n=5+(n-1)\cdot 7=7n-2$.\\
\textbf{b.}~$v=-3$: $u_n=20+(n-1)\cdot(-3)=23-3n$.\\
\textbf{c.}~$v=\tfrac12$: $u_n=\tfrac14+(n-1)\cdot\tfrac12=\tfrac12 n-\tfrac14$.\\
Het recursief voorschrift is telkens $u_n=u_{n-1}+v$ met de eerste term erbij.
\end{fullsolution}
\end{oefening}

\begin{oefening}[1]\label{oef:mtwee}
\thema{meetkundig: twee termen}Bepaal $u_1$ en $q$.
\begin{enumerate}
  \item $u_2 = 6$ en $u_5 = 162$
  \item $u_3 = 20$ en $u_6 = 2{,}5$
  \item $u_2 = -4$ en $u_5 = 32$
\end{enumerate}
\begin{solution}
\textbf{a.}~$q=3$, $u_1=2$ \quad \textbf{b.}~$q=\tfrac12$, $u_1=80$ \quad \textbf{c.}~$q=-2$, $u_1=2$
\end{solution}
\begin{fullsolution}
\textbf{a.}~3 sprongen: $q^3=27$, dus $q=3$ en $u_1=\tfrac63=2$.\\
\textbf{b.}~3 sprongen: $q^3=\tfrac18$, dus $q=\tfrac12$ en $u_1=20\cdot 2^2=80$.\\
\textbf{c.}~3 sprongen: $q^3=-8$, dus $q=-2$ en $u_1=\tfrac{-4}{-2}=2$.
\end{fullsolution}
\end{oefening}

\begin{oefening}[1]
\thema{rekenkundig midden}Bepaal $x$ zodat de drie getallen opeenvolgende termen van een rekenkundige rij zijn.
\begin{reeks}
  \item $3,\ x,\ 11$
  \item $-4,\ x,\ 10$
  \item $x,\ 6,\ 14$
\end{reeks}
\begin{solution}
\textbf{a.}~7 \quad \textbf{b.}~3 \quad \textbf{c.}~$-2$
\end{solution}
\begin{fullsolution}
\textbf{a.}~$x=\dfrac{3+11}{2}=7$.\\
\textbf{b.}~$x=\dfrac{-4+10}{2}=3$.\\
\textbf{c.}~$6=\dfrac{x+14}{2}$ geeft $x=-2$.
\end{fullsolution}
\end{oefening}

\begin{oefening}[1]
\thema{rekenkundig: voorschriften}Iemand schrijft: `$u_1 = 5$ en $v = 4$, dus $u_{10} = 5 + 10 \cdot 4 = 45$.' Wat is er fout?
\begin{solution}
9 sprongen: $u_{10}=41$.
\end{solution}
\begin{fullsolution}
Van $u_1$ naar $u_{10}$ zijn het 9 sprongen, geen 10. Dus $u_{10}=5+9\cdot 4=41$.
\end{fullsolution}
\end{oefening}

\begin{oefening}[1]\label{oef:intrest}
\thema{intrest}Je belegt \eur 3000 aan 2,5\% per jaar. Hoeveel heb je na 8 jaar, met enkelvoudige en met samengestelde intrest?
\begin{solution}
Enkelvoudig \eur 3600; samengesteld ongeveer \eur 3655,21
\end{solution}
\begin{fullsolution}
Enkelvoudig: $3000\,(1+8\cdot 0{,}025)=3600$. Samengesteld: $3000\cdot 1{,}025^8\approx 3655{,}21$.
\end{fullsolution}
\end{oefening}

\begin{oefening}[1]
\thema{rekenkundig: voorschriften}Bepaal $u_1$ en $v$ uit het voorschrift.
\begin{reeks}
  \item $u_n = 5n - 2$
  \item $u_n = 7 - 2n$
  \item $u_n = \tfrac13 n + 1$
\end{reeks}
\begin{solution}
\textbf{a.}~$u_1=3$, $v=5$ \quad \textbf{b.}~$u_1=5$, $v=-2$ \quad \textbf{c.}~$u_1=\tfrac43$, $v=\tfrac13$
\end{solution}
\begin{fullsolution}
Vul $n=1$ in voor $u_1$; het verschil is de coëfficiënt van $n$.\\
\textbf{a.}~$u_1=3$, $v=5$.\\
\textbf{b.}~$u_1=5$, $v=-2$.\\
\textbf{c.}~$u_1=\tfrac43$, $v=\tfrac13$.
\end{fullsolution}
\end{oefening}

\begin{oefening}[1]
\thema{meetkundig: voorschriften}Is de rij meetkundig? Zo ja, geef $q$.
\begin{reeks}
  \item $2,\ 6,\ 18,\ 54,\ \ldots$
  \item $5,\ 10,\ 15,\ 20,\ \ldots$
  \item $16,\ {-8},\ 4,\ {-2},\ \ldots$
  \item $7,\ 7,\ 7,\ 7,\ \ldots$
  \item $1,\ 4,\ 9,\ 16,\ \ldots$
  \item $0{,}3;\ 0{,}03;\ 0{,}003;\ \ldots$
\end{reeks}
\begin{solution}
\textbf{a.}~ja, $q=3$ \quad \textbf{b.}~nee \quad \textbf{c.}~ja, $q=-\tfrac12$ \quad \textbf{d.}~ja, $q=1$ \quad \textbf{e.}~nee \quad \textbf{f.}~ja, $q=0{,}1$
\end{solution}
\begin{fullsolution}
Kijk naar het quotiënt van opeenvolgende termen.\\
\textbf{a.}~Telkens 3.\\
\textbf{b.}~Quotiënten $2$ en $\tfrac32$: niet meetkundig (wel rekenkundig).\\
\textbf{c.}~Telkens $-\tfrac12$.\\
\textbf{d.}~Telkens 1.\\
\textbf{e.}~Quotiënten $4$ en $\tfrac94$: niet meetkundig.\\
\textbf{f.}~Telkens $0{,}1$.
\end{fullsolution}
\end{oefening}

\begin{oefening}[1]
\thema{meetkundig midden}Bepaal $x$ zodat de drie getallen opeenvolgende termen van een meetkundige rij zijn.
\begin{reeks}
  \item $4,\ x,\ 9$
  \item $2,\ x,\ 32$
  \item $x,\ 6,\ 12$
\end{reeks}
\begin{solution}
\textbf{a.}~6 of $-6$ \quad \textbf{b.}~8 of $-8$ \quad \textbf{c.}~3
\end{solution}
\begin{fullsolution}
\textbf{a.}~$x^2=36$.\\
\textbf{b.}~$x^2=64$.\\
\textbf{c.}~$36=12x$, dus $x=3$.
\end{fullsolution}
\end{oefening}

\begin{oefening}[1]\label{oef:groei}
\thema{rekenkundig: lineaire groei}Een bibliotheek rekent \eur 2 boete plus \eur 0,50 per dag dat je te laat bent.
\begin{enumerate}
  \item Geef een voorschrift voor de boete na $n$ dagen.
  \item Hoeveel betaal je na 14 dagen?
  \item Na hoeveel dagen is de boete \eur 12?
\end{enumerate}
\begin{solution}
\textbf{a.}~$u_n=2+0{,}5\,n$ \quad \textbf{b.}~\eur 9 \quad \textbf{c.}~20 dagen
\end{solution}
\begin{fullsolution}
\textbf{a.}~$u_0=2$ en $v=0{,}5$, dus $u_n=2+0{,}5\,n$.\\
\textbf{b.}~$u_{14}=2+7=9$.\\
\textbf{c.}~$2+0{,}5\,n=12$ geeft $n=20$.
\end{fullsolution}
\end{oefening}

\begin{oefening}[1]\label{oef:herken}
\thema{rekenkundig: voorschriften}Is de rij rekenkundig? Zo ja, geef $v$.
\begin{reeks}
  \item $4,\ 9,\ 14,\ 19,\ \ldots$
  \item $2,\ 4,\ 8,\ 16,\ \ldots$
  \item $-3,\ {-3},\ {-3},\ \ldots$
  \item $\tfrac12,\ \tfrac13,\ \tfrac16,\ 0,\ \ldots$
  \item $1,\ 4,\ 9,\ 16,\ \ldots$
  \item $10;\ 7{,}5;\ 5;\ 2{,}5;\ \ldots$
\end{reeks}
\begin{solution}
\textbf{a.}~ja, $v=5$ \quad \textbf{b.}~nee \quad \textbf{c.}~ja, $v=0$ \quad \textbf{d.}~ja, $v=-\tfrac16$ \quad \textbf{e.}~nee \quad \textbf{f.}~ja, $v=-2{,}5$
\end{solution}
\begin{fullsolution}
Kijk naar het verschil van opeenvolgende termen.\\
\textbf{a.}~Telkens 5.\\
\textbf{b.}~Verschillen 2, 4, 8: niet rekenkundig.\\
\textbf{c.}~Telkens 0.\\
\textbf{d.}~Telkens $-\tfrac16$.\\
\textbf{e.}~Verschillen 3, 5, 7: niet rekenkundig.\\
\textbf{f.}~Telkens $-2{,}5$.
\end{fullsolution}
\end{oefening}

\begin{oefening}[1]
\thema{meetkundig: voorschriften}Iemand schrijft: `$u_1 = 3$ en $q = 2$, dus $u_5 = 6^4 = 1296$.' Wat is er fout?
\begin{solution}
Alleen $q$ staat in de macht: $u_5=48$.
\end{solution}
\begin{fullsolution}
De macht hoort alleen bij $q$: $u_5=3\cdot 2^4=48$. Je mag $3\cdot 2^4$ niet schrijven als $6^4$.
\end{fullsolution}
\end{oefening}

\begin{oefening}[1]
\thema{rekenkundig: twee termen}$u_1 = 7$ en $u_{12} = -37$. Bepaal $v$ en $u_{20}$.
\begin{solution}
$v=-4$; $u_{20}=-69$
\end{solution}
\begin{fullsolution}
11 sprongen: $11v=-44$, dus $v=-4$. Dan is $u_{20}=7+19\cdot(-4)=-69$.
\end{fullsolution}
\end{oefening}

\begin{oefening}[1]
\thema{gemengd}$u_1 = 1$ en $u_n = 2u_{n-1} + 1$. Bereken $u_2$, $u_3$ en $u_4$. Is de rij rekenkundig of meetkundig?
\begin{solution}
3, 7, 15; geen van beide
\end{solution}
\begin{fullsolution}
$u_2=2\cdot 1+1=3$, $u_3=2\cdot 3+1=7$, $u_4=2\cdot 7+1=15$. Verschillen 2, 4, 8 en quotiënten 3, $\tfrac73$, $\tfrac{15}{7}$ zijn niet constant.
\end{fullsolution}
\end{oefening}

\begin{oefening}[1]
\thema{rekenkundig: voorschriften}Gegeven: $u_1 = -4$ en $u_n = u_{n-1} + 3$. Geef het expliciet voorschrift en bereken $u_{50}$.
\begin{solution}
$u_n=3n-7$; $u_{50}=143$
\end{solution}
\begin{fullsolution}
$u_n=-4+(n-1)\cdot 3=3n-7$, dus $u_{50}=3\cdot 50-7=143$.
\end{fullsolution}
\end{oefening}

\begin{oefening}[1]
\thema{gemengd}Geef een expliciet voorschrift van $\tfrac12,\ \tfrac23,\ \tfrac34,\ \tfrac45,\ \ldots$
\begin{solution}
$u_n=\dfrac{n}{n+1}$
\end{solution}
\begin{fullsolution}
De teller is het rangnummer, de noemer is één meer: $u_n=\dfrac{n}{n+1}$.
\end{fullsolution}
\end{oefening}

\begin{oefening}[1]
\thema{rekenkundig: lineaire groei}Een taxirit kost \eur 4,50 instapgeld plus \eur 1,80 per kilometer.
\begin{enumerate}
  \item Wat kost een rit van 12 km?
  \item Hoe lang was een rit van \eur 31,50?
\end{enumerate}
\begin{solution}
\textbf{a.}~\eur 26,10 \quad \textbf{b.}~15 km
\end{solution}
\begin{fullsolution}
$u_n=4{,}50+1{,}80\,n$.\\
\textbf{a.}~$u_{12}=4{,}50+21{,}60=26{,}10$.\\
\textbf{b.}~$4{,}50+1{,}80\,n=31{,}50$ geeft $n=15$.
\end{fullsolution}
\end{oefening}

\begin{oefening}[1]
\thema{gemengd}Wat is het kleinste natuurlijke getal $n$ waarvoor $5n - 3 > 100$?
\begin{solution}
$n=21$
\end{solution}
\begin{fullsolution}
$5n>103$, dus $n>20{,}6$ en $n=21$.
\end{fullsolution}
\end{oefening}

\begin{oefening}[2]
\thema{meetkundig: voorschriften}Bepaal $u_1$ en $q$ uit het voorschrift.
\begin{reeks}
  \item $u_n = 5 \cdot 4^{\,n-1}$
  \item $u_n = 3 \cdot 2^{\,n}$
  \item $u_n = 10 \cdot (\tfrac12)^{n+1}$
\end{reeks}
\begin{solution}
\textbf{a.}~$u_1=5$, $q=4$ \quad \textbf{b.}~$u_1=6$, $q=2$ \quad \textbf{c.}~$u_1=\tfrac52$, $q=\tfrac12$
\end{solution}
\begin{fullsolution}
Vul $n=1$ in voor $u_1$; het quotiënt is het grondtal.\\
\textbf{a.}~$u_1=5$, $q=4$.\\
\textbf{b.}~$u_1=3\cdot 2=6$, $q=2$.\\
\textbf{c.}~$u_1=10\cdot\tfrac14=\tfrac52$, $q=\tfrac12$.
\end{fullsolution}
\end{oefening}

\begin{oefening}[2]
\thema{rekenkundig: voorschriften}Een zaal heeft 14 stoelen op de eerste rij. Elke volgende rij telt 2 stoelen meer.
\begin{enumerate}
  \item Hoeveel stoelen staan er op rij 18?
  \item Op welke rij staan 60 stoelen?
\end{enumerate}
\begin{solution}
\textbf{a.}~48 \quad \textbf{b.}~rij 24
\end{solution}
\begin{fullsolution}
\textbf{a.}~$u_{18}=14+17\cdot 2=48$.\\
\textbf{b.}~$14+(n-1)\cdot 2=60$ geeft $n-1=23$, dus $n=24$.
\end{fullsolution}
\end{oefening}

\begin{oefening}[2]
\thema{meetkundig: exponentiële groei}Een algenlaag van 40 m\textsuperscript{2} groeit met 25\% per week.
\begin{enumerate}
  \item Geef een voorschrift voor de oppervlakte na $n$ weken.
  \item Hoe groot is de laag na 6 weken?
  \item Wanneer is de laag voor het eerst groter dan 500 m\textsuperscript{2}?
\end{enumerate}
\begin{solution}
\textbf{a.}~$u_n=40\cdot 1{,}25^{\,n}$ \quad \textbf{b.}~ongeveer 153 m\textsuperscript{2} \quad \textbf{c.}~na 12 weken
\end{solution}
\begin{fullsolution}
\textbf{a.}~$q=1{,}25$, dus $u_n=40\cdot 1{,}25^{\,n}$.\\
\textbf{b.}~$u_6=40\cdot 1{,}25^6\approx 152{,}6$.\\
\textbf{c.}~$1{,}25^{\,n}>12{,}5$ geeft $n>11{,}32$, dus na 12 weken.
\end{fullsolution}
\end{oefening}

\begin{oefening}[2]
\thema{rekenkundig: twee termen}Geef het expliciet voorschrift als $u_{10} = -4$ en $u_{20} = -34$.
\begin{solution}
$u_n=26-3n$
\end{solution}
\begin{fullsolution}
10 sprongen: $10v=-30$, dus $v=-3$. Dan is $u_n=-4+(n-10)\cdot(-3)=26-3n$.
\end{fullsolution}
\end{oefening}

\begin{oefening}[2]
\thema{meetkundig: twee termen}$u_2 = 5$ en $u_4 = 45$. Bepaal $u_1$ en $q$. Let op: er zijn twee oplossingen.
\begin{solution}
$q=3$, $u_1=\tfrac53$ of $q=-3$, $u_1=-\tfrac53$
\end{solution}
\begin{fullsolution}
2 sprongen: $q^2=9$, dus $q=3$ of $q=-3$. Een sprong terug: $u_1=\tfrac{5}{q}$, dus $\tfrac53$ of $-\tfrac53$.
\end{fullsolution}
\end{oefening}

\begin{oefening}[2]
\thema{rekenkundig: term of geen term}Welke term van $3,\ 7,\ 11,\ \ldots$ is de eerste die groter is dan 200?
\begin{solution}
$u_{51}=203$
\end{solution}
\begin{fullsolution}
$u_n=4n-1>200$ geeft $n>50{,}25$, dus $n=51$ en $u_{51}=203$.
\end{fullsolution}
\end{oefening}

\begin{oefening}[2]
\thema{meetkundig: term of geen term}Welke term van $3,\ 6,\ 12,\ \ldots$ is de eerste die groter is dan 1000?
\begin{solution}
$u_{10}=1536$
\end{solution}
\begin{fullsolution}
$3\cdot 2^{n-1}>1000$ geeft $2^{n-1}>333{,}33\ldots$ en $n-1>8{,}38$, dus $n=10$ en $u_{10}=1536$.
\end{fullsolution}
\end{oefening}

\begin{oefening}[2]
\thema{rekenkundig midden}$x + 2$, $2x + 1$ en $4x - 6$ zijn opeenvolgende termen van een rekenkundige rij. Bepaal $x$ en de drie termen.
\begin{solution}
$x=6$; 8, 13, 18
\end{solution}
\begin{fullsolution}
$2(2x+1)=(x+2)+(4x-6)$ geeft $4x+2=5x-4$, dus $x=6$. De termen zijn 8, 13 en 18.
\end{fullsolution}
\end{oefening}

\begin{oefening}[2]
\thema{intrest}Hoeveel jaar duurt het voor een kapitaal verdubbelt aan 4\% samengestelde intrest?
\begin{solution}
18 jaar
\end{solution}
\begin{fullsolution}
$1{,}04^{\,n}=2$ geeft $n\approx 17{,}67$, dus na 18 jaar is het kapitaal zeker verdubbeld.
\end{fullsolution}
\end{oefening}

\begin{oefening}[2]
\thema{rekenkundig: lineaire groei}Je legt vierkanten met lucifers, op één rij.

\begin{tikzpicture}[x=.42cm,y=.42cm,line width=1.1pt,line cap=round]
  \foreach \s/\k in {0/1,2/2,5/3}{
    \draw (\s,0) rectangle +(\k,1);
    \foreach \i in {1,...,\k}{\ifnum\i<\k\relax\draw (\s+\i,0) -- (\s+\i,1);\fi}}
\end{tikzpicture}
\begin{enumerate}
  \item Hoeveel lucifers zijn nodig voor 40 vierkanten?
  \item Hoeveel vierkanten kun je leggen met 200 lucifers?
\end{enumerate}
\begin{solution}
\textbf{a.}~121 \quad \textbf{b.}~66
\end{solution}
\begin{fullsolution}
De rij is $4,\ 7,\ 10,\ \ldots$ met $v=3$, dus $u_n=3n+1$.\\
\textbf{a.}~$u_{40}=121$.\\
\textbf{b.}~$3n+1\leqslant 200$ geeft $n\leqslant 66{,}33\ldots$, dus 66 vierkanten (199 lucifers).
\end{fullsolution}
\end{oefening}

\begin{oefening}[2]
\thema{meetkundig: voorschriften}Een bal valt van 2 m hoogte en stuit telkens terug tot 60\% van de vorige hoogte.
\begin{enumerate}
  \item Hoe hoog komt de bal na de 4de keer stuiten?
  \item Na hoeveel keer stuiten komt de bal voor het eerst niet meer boven 5 cm?
\end{enumerate}
\begin{solution}
\textbf{a.}~ongeveer 26 cm \quad \textbf{b.}~8 keer
\end{solution}
\begin{fullsolution}
$u_n=2\cdot 0{,}6^{\,n}$ in m.\\
\textbf{a.}~$u_4=2\cdot 0{,}1296\approx 0{,}26$.\\
\textbf{b.}~$2\cdot 0{,}6^{\,n}\leqslant 0{,}05$ geeft $0{,}6^{\,n}\leqslant 0{,}025$ en $n\geqslant 7{,}22$, dus 8 keer.
\end{fullsolution}
\end{oefening}

\begin{oefening}[2]
\thema{rekenkundig: lineaire groei}De punten $(1, 8)$, $(2, 5)$ en $(3, 2)$ liggen op de grafiek van een rij. Geef het voorschrift. Bij welke $n$ komt het eerste punt onder de $n$-as?
\begin{solution}
$u_n=11-3n$; bij $n=4$
\end{solution}
\begin{fullsolution}
$v=-3$ en $u_1=8$, dus $u_n=11-3n$. $11-3n<0$ geeft $n>3{,}67$, dus $n=4$.
\end{fullsolution}
\end{oefening}

\begin{oefening}[2]
\thema{meetkundig: exponentiële groei}De punten $(1, 3)$, $(2, 6)$ en $(3, 12)$ liggen op de grafiek van een meetkundige rij. Geef het voorschrift. Bij welke $n$ komt de grafiek voor het eerst boven 100?
\begin{solution}
$u_n=3\cdot 2^{n-1}$; bij $n=7$
\end{solution}
\begin{fullsolution}
$q=2$ en $u_1=3$. $3\cdot 2^{n-1}>100$ geeft $n-1>5{,}06$, dus $n=7$ ($u_6=96$ en $u_7=192$).
\end{fullsolution}
\end{oefening}

\begin{oefening}[2]
\thema{rekenkundig: voorschriften}De onderste sport van een ladder zit 30 cm boven de grond. Elke volgende sport zit 28 cm hoger.
\begin{enumerate}
  \item Hoe hoog zit sport 12?
  \item Welke sport zit op 198 cm?
\end{enumerate}
\begin{solution}
\textbf{a.}~338 cm \quad \textbf{b.}~sport 7
\end{solution}
\begin{fullsolution}
\textbf{a.}~$u_{12}=30+11\cdot 28=338$.\\
\textbf{b.}~$30+(n-1)\cdot 28=198$ geeft $n-1=6$, dus $n=7$.
\end{fullsolution}
\end{oefening}

\begin{oefening}[2]
\thema{meetkundig: exponentiële groei}Na inname zit er 200 mg van een medicijn in het bloed. Elk uur wordt 20\% afgebroken.
\begin{enumerate}
  \item Hoeveel zit er nog in het bloed na 5 uur?
  \item Na hoeveel uur is dat voor het eerst minder dan 10 mg?
\end{enumerate}
\begin{solution}
\textbf{a.}~ongeveer 65,5 mg \quad \textbf{b.}~na 14 uur
\end{solution}
\begin{fullsolution}
$u_n=200\cdot 0{,}8^{\,n}$.\\
\textbf{a.}~$u_5=200\cdot 0{,}32768\approx 65{,}5$.\\
\textbf{b.}~$0{,}8^{\,n}<0{,}05$ geeft $n>13{,}43$, dus na 14 uur.
\end{fullsolution}
\end{oefening}

\begin{oefening}[2]
\thema{rekenkundig: term of geen term}Hoeveel positieve termen telt de rij $50,\ 47,\ 44,\ \ldots$?
\begin{solution}
17
\end{solution}
\begin{fullsolution}
$u_n=53-3n>0$ geeft $n<17{,}67$, dus 17 termen. Controle: $u_{17}=2$ en $u_{18}=-1$.
\end{fullsolution}
\end{oefening}

\begin{oefening}[2]
\thema{meetkundig: term of geen term}Welke term van $1000,\ 500,\ 250,\ \ldots$ is de eerste die kleiner is dan 1?
\begin{solution}
$u_{11}\approx 0{,}98$
\end{solution}
\begin{fullsolution}
$1000\cdot(\tfrac12)^{n-1}<1$ geeft $2^{n-1}>1000$ en $n-1>9{,}97$, dus $n=11$ en $u_{11}=\tfrac{1000}{1024}\approx 0{,}98$.
\end{fullsolution}
\end{oefening}

\begin{oefening}[2]
\thema{meetkundig midden}$x - 2$, $x + 1$ en $x + 7$ zijn opeenvolgende termen van een meetkundige rij. Bepaal $x$ en de drie termen.
\begin{solution}
$x=5$; 3, 6, 12
\end{solution}
\begin{fullsolution}
$(x+1)^2=(x-2)(x+7)$ geeft $x^2+2x+1=x^2+5x-14$, dus $x=5$. De termen zijn 3, 6 en 12.
\end{fullsolution}
\end{oefening}

\begin{oefening}[2]
\thema{lineaire en exponentiële groei}Een loon van \eur 2000 stijgt elke maand met \eur 50. Een ander loon van \eur 2000 stijgt elke maand met 2\%. Welk loon is het hoogst na 12 maanden? En na 36 maanden?
\begin{solution}
Na 12 maanden het eerste; na 36 maanden het tweede.
\end{solution}
\begin{fullsolution}
Na 12 maanden: $2000+12\cdot 50=2600$ en $2000\cdot 1{,}02^{12}\approx 2536{,}48$. Na 36 maanden: $2000+36\cdot 50=3800$ en $2000\cdot 1{,}02^{36}\approx 4079{,}77$.
\end{fullsolution}
\end{oefening}

\begin{oefening}[2]
\thema{rekenkundig: lineaire groei}Een tank bevat 45 liter. Per 100 km verbruikt de auto 6 liter. Na hoeveel kilometer zit er nog 9 liter in?
\begin{solution}
600 km
\end{solution}
\begin{fullsolution}
Met $n$ het aantal keer 100 km is $u_n=45-6n$. $45-6n=9$ geeft $n=6$, dus 600 km.
\end{fullsolution}
\end{oefening}

\begin{oefening}[2]
\thema{meetkundig: twee termen}Voeg tussen 3 en 96 vier getallen in, zodat je een meetkundige rij krijgt.
\begin{solution}
6, 12, 24, 48
\end{solution}
\begin{fullsolution}
Met vier getallen ertussen zijn het 5 sprongen: $q^5=\tfrac{96}{3}=32$, dus $q=2$.
\end{fullsolution}
\end{oefening}

\begin{oefening}[2]
\thema{rekenkundig: lineaire groei}Op rekening A staat \eur 200, met elke maand \eur 30 erbij. Op B staat \eur 500, met elke maand \eur 10 erbij. Wanneer staat op beide evenveel? Hoeveel is dat?
\begin{solution}
Na 15 maanden; \eur 650
\end{solution}
\begin{fullsolution}
$200+30n=500+10n$ geeft $n=15$. Dan staat er $200+30\cdot 15=650$ euro op beide.
\end{fullsolution}
\end{oefening}

\begin{oefening}[2]
\thema{rekenkundig: term of geen term}Hoeveel termen telt de rij $15,\ 18,\ 21,\ \ldots,\ 303$?
\begin{solution}
97
\end{solution}
\begin{fullsolution}
$15+(n-1)\cdot 3=303$ geeft $n-1=96$, dus $n=97$. Let op: 96 is het aantal sprongen, niet het aantal termen.
\end{fullsolution}
\end{oefening}

\begin{oefening}[2]
\thema{intrest}\eur 1000 groeit in 2 jaar met samengestelde intrest aan tot \eur 1210. Wat is de intrestvoet?
\begin{solution}
10\%
\end{solution}
\begin{fullsolution}
$(1+\tfrac{p}{100})^2=1{,}21$ geeft $1+\tfrac{p}{100}=1{,}1$, dus $p=10$.
\end{fullsolution}
\end{oefening}

\begin{oefening}[2]
\thema{rekenkundig: twee termen}Voeg tussen 4 en 34 vijf getallen in, zodat je een rekenkundige rij krijgt.
\begin{solution}
9, 14, 19, 24, 29
\end{solution}
\begin{fullsolution}
Met vijf getallen ertussen heeft de rij 7 termen, dus 6 sprongen: $6v=30$ en $v=5$.
\end{fullsolution}
\end{oefening}

\begin{oefening}[2]
\thema{meetkundig: voorschriften}Een blad papier is 0,1 mm dik. Bij elke plooi verdubbelt de dikte.
\begin{enumerate}
  \item Hoe dik is het na 7 plooien?
  \item Hoeveel plooien heb je in theorie nodig om dikker te zijn dan 1 meter?
\end{enumerate}
\begin{solution}
\textbf{a.}~12,8 mm \quad \textbf{b.}~14
\end{solution}
\begin{fullsolution}
$u_n=0{,}1\cdot 2^{\,n}$ in mm.\\
\textbf{a.}~$u_7=0{,}1\cdot 128=12{,}8$.\\
\textbf{b.}~$0{,}1\cdot 2^{\,n}>1000$ geeft $2^{\,n}>10\,000$ en $n>13{,}29$, dus 14 plooien.
\end{fullsolution}
\end{oefening}

\begin{oefening}[2]
\thema{meetkundig: term of geen term}Hoeveel termen telt de rij $5,\ 10,\ 20,\ \ldots,\ 5120$?
\begin{solution}
11
\end{solution}
\begin{fullsolution}
$5\cdot 2^{n-1}=5120$ geeft $2^{n-1}=1024=2^{10}$, dus $n-1=10$ en $n=11$.
\end{fullsolution}
\end{oefening}

\begin{oefening}[2]
\thema{rekenkundig: term of geen term}Welke term van $2,\ 5,\ 8,\ \ldots$ ligt het dichtst bij 1000?
\begin{solution}
$u_{334}=1001$
\end{solution}
\begin{fullsolution}
$u_n=3n-1$. Uit $3n-1=1000$ volgt $n=333{,}67$. Vergelijk: $u_{333}=998$ en $u_{334}=1001$, dus $u_{334}$ ligt het dichtst bij.
\end{fullsolution}
\end{oefening}

\begin{oefening}[2]
\thema{rekenkundig: voorschriften}$u_n = -2n + 9$. Bepaal $u_1$ en $v$. Vanaf welke $n$ zijn de termen negatief?
\begin{solution}
$u_1=7$, $v=-2$; vanaf $n=5$
\end{solution}
\begin{fullsolution}
$u_1=7$ en $v=-2$. $-2n+9<0$ geeft $n>4{,}5$, dus vanaf $n=5$.
\end{fullsolution}
\end{oefening}

\begin{oefening}[2]
\thema{rekenkundig: twee termen}De punten $(2, 5)$ en $(6, 17)$ liggen op de grafiek van een rekenkundige rij. Ligt $(30, 90)$ er ook op?
\begin{solution}
Nee, $u_{30}=89$.
\end{solution}
\begin{fullsolution}
4 sprongen: $4v=12$, dus $v=3$ en $u_1=2$. Dan is $u_{30}=2+29\cdot 3=89$.
\end{fullsolution}
\end{oefening}

\begin{oefening}[2]
\thema{intrest}Hoeveel moet je nu beleggen aan 3\% samengestelde intrest om over 5 jaar \eur 10\,000 te hebben?
\begin{solution}
Ongeveer \eur 8626,09
\end{solution}
\begin{fullsolution}
$K_0\cdot 1{,}03^5=10\,000$ geeft $K_0=\dfrac{10\,000}{1{,}03^5}\approx 8626{,}09$.
\end{fullsolution}
\end{oefening}

\begin{oefening}[3]
\thema{meetkundig: exponentiële groei}Een hoeveelheid daalt met 10\% per jaar. Na hoeveel jaar is ze gehalveerd?
\begin{solution}
Na ongeveer 6,6 jaar
\end{solution}
\begin{fullsolution}
$0{,}9^{\,n}=0{,}5$ geeft $n=\log_{0{,}9} 0{,}5\approx 6{,}6$. Dit hangt niet af van de beginhoeveelheid.
\end{fullsolution}
\end{oefening}

\begin{oefening}[3]
\thema{rekenkundig: term of geen term}Hoeveel getallen van drie cijfers zijn een veelvoud van 7?
\begin{solution}
128
\end{solution}
\begin{fullsolution}
Het kleinste is $105=7\cdot 15$ en het grootste is $994=7\cdot 142$. De rij $105,\ 112,\ \ldots,\ 994$ telt $142-15+1=128$ termen.
\end{fullsolution}
\end{oefening}

\begin{oefening}[3]
\thema{rekenkundig: twee termen}$u_2 + u_4 = 22$ en $u_6 = 23$. Bepaal $u_1$ en $v$.
\begin{solution}
$v=4$, $u_1=3$
\end{solution}
\begin{fullsolution}
$2u_1+4v=22$ en $u_1+5v=23$. Uit de eerste volgt $u_1=11-2v$; invullen geeft $11+3v=23$, dus $v=4$ en $u_1=3$.
\end{fullsolution}
\end{oefening}

\begin{oefening}[3]
\thema{meetkundig midden}Drie opeenvolgende termen van een meetkundige rij hebben als product 216 en als som 19. Welke termen zijn het?
\begin{solution}
4, 6, 9 (of 9, 6, 4)
\end{solution}
\begin{fullsolution}
Noem de termen $\tfrac{a}{q}$, $a$ en $aq$. Het product is $a^3=216$, dus $a=6$. De som geeft $\tfrac6q+6+6q=19$, dus $6q^2-13q+6=0$ en $q=\tfrac32$ of $q=\tfrac23$.
\end{fullsolution}
\end{oefening}

\begin{oefening}[3]
\thema{rekenkundig: voorschriften}Bewijs: als $u_n = pn + q$ met $p$ en $q$ vaste getallen, dan is de rij rekenkundig. Wat is het verschil?
\begin{solution}
Het verschil is $p$.
\end{solution}
\begin{fullsolution}
$u_{n+1}-u_n=p(n+1)+q-(pn+q)=p$. Dit hangt niet af van $n$, dus de rij is rekenkundig met $v=p$.
\end{fullsolution}
\end{oefening}

\begin{oefening}[3]
\thema{rekenkundig midden}Drie opeenvolgende termen van een rekenkundige rij hebben als som 33 en als product 1155. Welke termen zijn het?
\begin{solution}
7, 11, 15 (of 15, 11, 7)
\end{solution}
\begin{fullsolution}
Noem de termen $a-v$, $a$ en $a+v$. De som is $3a=33$, dus $a=11$. Het product is $11(121-v^2)=1155$, dus $v^2=16$ en $v=\pm 4$.
\end{fullsolution}
\end{oefening}

\begin{oefening}[3]
\thema{meetkundig: voorschriften}Bewijs: als $u_n = a \cdot b^{\,n}$ met $a$ en $b$ vaste getallen verschillend van 0, dan is de rij meetkundig. Wat is het quotiënt?
\begin{solution}
Het quotiënt is $b$.
\end{solution}
\begin{fullsolution}
$\dfrac{u_{n+1}}{u_n}=\dfrac{a\cdot b^{\,n+1}}{a\cdot b^{\,n}}=b$. Dit hangt niet af van $n$, dus de rij is meetkundig met $q=b$.
\end{fullsolution}
\end{oefening}

\end{oefeningen}

%----------------------------------------------------------------------------------------
% KOPPEL: verbanden met andere leerstof. Enkel waar het zinvol is; geen koppeling tussen
% leerstof die vanzelf uit elkaar volgt (zoals limieten, afgeleiden en integralen).
%----------------------------------------------------------------------------------------
\opdrachtsectie{Koppel}

\lkr{Kijk na of deze leerstof al gezien is in jouw jaarplanning. Kansrekenen staat bv. ook in het 6de jaar.}

\begin{oefeningen}
\begin{oefening}[2]
\koppel{analyse}{Exponentiële functies} Gegeven: $f(x) = 5 \cdot 2^{\,x}$.
\begin{enumerate}
  \item Bereken $f(1)$, $f(2)$, $f(3)$ en $f(4)$. Welke soort rij vormen die getallen? Wat is het quotiënt?
  \item Wat is het quotiënt als je alleen $f(1)$, $f(3)$, $f(5)$, \ldots\ neemt?
\end{enumerate}
\begin{solution}
\textbf{a.}~10, 20, 40, 80; meetkundig, $q=2$ \quad \textbf{b.}~$q=4$
\end{solution}
\begin{fullsolution}
\textbf{a.}~Elke stap van 1 in $x$ vermenigvuldigt met 2, dus $q=2$.\\
\textbf{b.}~$\dfrac{f(3)}{f(1)}=2^2=4$.
\end{fullsolution}
\end{oefening}

\begin{oefening}[3]
\koppel{analyse}{Logaritmen} Gegeven is de meetkundige rij $2,\ 8,\ 32,\ 128,\ \ldots$
\begin{enumerate}
  \item Bereken $\log_2$ van elke term. Welke soort rij krijg je?
  \item Toon aan: de logaritmen van de termen van een meetkundige rij met positieve termen vormen een rekenkundige rij. Wat is het verschil?
\end{enumerate}
\begin{solution}
\textbf{a.}~1, 3, 5, 7; rekenkundig, $v=2$ \quad \textbf{b.}~$v=\log q$
\end{solution}
\begin{fullsolution}
\textbf{a.}~$\log_2 2=1$, $\log_2 8=3$, $\log_2 32=5$, $\log_2 128=7$.\\
\textbf{b.}~$\log u_n=\log\left(u_1 q^{\,n-1}\right)=\log u_1+(n-1)\log q$. Dit is een rekenkundige rij met eerste term $\log u_1$ en verschil $\log q$.
\end{fullsolution}
\end{oefening}

\begin{oefening}[2]
\koppel{algebra}{Matrices} Gegeven: $A = \begin{pmatrix} 1 & 2 \\ 0 & 1 \end{pmatrix}$.
\begin{enumerate}
  \item Bereken $A^2$ en $A^3$. Welke rij vormen de elementen rechtsboven in $A$, $A^2$, $A^3$?
  \item Bepaal $A^{10}$ zonder verder te vermenigvuldigen.
\end{enumerate}
\begin{solution}
\textbf{a.}~Rechtsboven: 2, 4, 6; rekenkundig, $v=2$ \quad \textbf{b.}~$A^{10}=\begin{pmatrix} 1 & 20 \\ 0 & 1 \end{pmatrix}$
\end{solution}
\begin{fullsolution}
\textbf{a.}~$A^2=\begin{pmatrix} 1 & 4 \\ 0 & 1 \end{pmatrix}$ en $A^3=\begin{pmatrix} 1 & 6 \\ 0 & 1 \end{pmatrix}$.\\
\textbf{b.}~Het element rechtsboven is $u_n=2n$, dus $u_{10}=20$.
\end{fullsolution}
\end{oefening}

\begin{oefening}[2]
\koppel{algebra}{Complexe getallen} Gegeven: $u_n = (1 + \mathrm{i})^n$.
\begin{enumerate}
  \item Bereken $u_1$, $u_2$, $u_3$ en $u_4$. Wat is het quotiënt van deze rij?
  \item Welke rij vormen de moduli van deze termen?
\end{enumerate}
\begin{solution}
\textbf{a.}~$1+\mathrm{i}$, $2\mathrm{i}$, $-2+2\mathrm{i}$, $-4$; $q=1+\mathrm{i}$ \quad \textbf{b.}~meetkundig, $q=\sqrt2$
\end{solution}
\begin{fullsolution}
\textbf{a.}~Telkens vermenigvuldigen met $1+\mathrm{i}$: $(1+\mathrm{i})^2=2\mathrm{i}$, $2\mathrm{i}(1+\mathrm{i})=-2+2\mathrm{i}$, $(2\mathrm{i})^2=-4$.\\
\textbf{b.}~De moduli zijn $\sqrt2$, $2$, $2\sqrt2$, $4$: telkens maal $\sqrt2$.
\end{fullsolution}
\end{oefening}

\needspace{14\baselineskip}% figuur en deelvragen bij de opgave houden
\begin{oefening}[2]
\koppel{meetkunde}{Vlakke meetkunde} Een vierkant heeft een oppervlakte van 16 cm\textsuperscript{2}. Door de middens van de zijden te verbinden, krijg je een nieuw vierkant, enzovoort.

\begin{tikzpicture}[scale=.34,line width=.7pt]
  \draw (0,0) -- (4,0) -- (4,4) -- (0,4) -- cycle;
  \draw (2,0) -- (4,2) -- (2,4) -- (0,2) -- cycle;
  \draw (1,1) -- (3,1) -- (3,3) -- (1,3) -- cycle;
\end{tikzpicture}
\begin{enumerate}
  \item Toon aan dat de oppervlakte telkens halveert.
  \item Wat is de oppervlakte van het 6de vierkant?
\end{enumerate}
\begin{solution}
\textbf{b.}~$\tfrac12$ cm\textsuperscript{2}
\end{solution}
\begin{fullsolution}
\textbf{a.}~De vier hoekdriehoeken die wegvallen, zijn samen even groot als het nieuwe vierkant. De oppervlakte halveert dus: $q=\tfrac12$.\\
\textbf{b.}~$u_6=16\cdot(\tfrac12)^5=\tfrac12$.
\end{fullsolution}
\end{oefening}

\begin{oefening}[2]
\begin{metprent}{dobbelstenen}
\koppel{data}{Kansrekenen} Je gooit met een dobbelsteen tot je een zes gooit.
\begin{enumerate}
  \item Bereken de kans dat de eerste zes valt bij worp 1, worp 2 en worp 3. Welke soort rij is dat?
  \item Vanaf welke worp is die kans kleiner dan 1\%?
\end{enumerate}
\end{metprent}
\begin{solution}
\textbf{a.}~$\tfrac16$, $\tfrac{5}{36}$, $\tfrac{25}{216}$; meetkundig, $q=\tfrac56$ \quad \textbf{b.}~vanaf worp 17
\end{solution}
\begin{fullsolution}
\textbf{a.}~Eerst $n-1$ keer geen zes, dan wel een zes: $u_n=\tfrac16\cdot(\tfrac56)^{n-1}$.\\
\textbf{b.}~$\tfrac16\cdot(\tfrac56)^{n-1}<0{,}01$ geeft $(\tfrac56)^{n-1}<0{,}06$ en $n-1>15{,}43$, dus $n=17$.
\end{fullsolution}
\end{oefening}
\end{oefeningen}

%----------------------------------------------------------------------------------------
\opdrachtsectie{Test jezelf}

\begin{oefeningen}
\begin{oefening}\label{oef:t1}
Is de rij rekenkundig, meetkundig of geen van beide? Geef $v$ of $q$.
\begin{reeks}
  \item $11,\ 8,\ 5,\ 2,\ \ldots$
  \item $3,\ {-6},\ 12,\ {-24},\ \ldots$
  \item $1,\ 2,\ 4,\ 7,\ \ldots$
\end{reeks}
\begin{solution}
\textbf{a.}~rekenkundig, $v=-3$ \quad \textbf{b.}~meetkundig, $q=-2$ \quad \textbf{c.}~geen van beide
\end{solution}
\begin{fullsolution}
\textbf{a.}~Telkens $-3$.\\
\textbf{b.}~Telkens maal $-2$.\\
\textbf{c.}~Verschillen 1, 2, 3 en quotiënten 2, 2, $\tfrac74$: niet constant.
\end{fullsolution}
\end{oefening}

\begin{oefening}\label{oef:t2}
Geef een recursief en een expliciet voorschrift van $4,\ 10,\ 16,\ \ldots$
\begin{solution}
$u_n=u_{n-1}+6$, $u_1=4$; $u_n=6n-2$
\end{solution}
\begin{fullsolution}
$v=6$, dus $u_n=u_{n-1}+6$ met $u_1=4$, en $u_n=4+(n-1)\cdot 6=6n-2$.
\end{fullsolution}
\end{oefening}

\begin{oefening}\label{oef:t3}
Geef een recursief en een expliciet voorschrift van $2,\ 10,\ 50,\ \ldots$
\begin{solution}
$u_n=5u_{n-1}$, $u_1=2$; $u_n=2\cdot 5^{n-1}$
\end{solution}
\begin{fullsolution}
$q=5$, dus $u_n=5u_{n-1}$ met $u_1=2$, en $u_n=2\cdot 5^{n-1}$.
\end{fullsolution}
\end{oefening}

\begin{oefening}\label{oef:t4}
Een rekenkundige rij heeft $u_4 = 13$ en $u_{11} = 41$. Bepaal $u_1$ en $v$.
\begin{solution}
$v=4$, $u_1=1$
\end{solution}
\begin{fullsolution}
7 sprongen: $7v=28$, dus $v=4$. Drie sprongen terug: $u_1=13-3\cdot 4=1$.
\end{fullsolution}
\end{oefening}

\begin{oefening}\label{oef:t5}
Een meetkundige rij heeft $u_2 = 6$ en $u_5 = 48$. Bepaal $u_1$ en $q$.
\begin{solution}
$q=2$, $u_1=3$
\end{solution}
\begin{fullsolution}
3 sprongen: $q^3=8$, dus $q=2$. Een sprong terug: $u_1=\tfrac62=3$.
\end{fullsolution}
\end{oefening}

\begin{oefening}\label{oef:t6}
Is 100 een term van $4,\ 10,\ 16,\ \ldots$? En 101?
\begin{solution}
100 is de 17de term; 101 is geen term.
\end{solution}
\begin{fullsolution}
$6n-2=100$ geeft $n=17$. $6n-2=101$ geeft $n=\tfrac{103}{6}$: geen natuurlijk getal.
\end{fullsolution}
\end{oefening}

\begin{oefening}\label{oef:t7}
Is 1024 een term van $2,\ 4,\ 8,\ \ldots$? En 1000?
\begin{solution}
1024 is de 10de term; 1000 is geen term.
\end{solution}
\begin{fullsolution}
$u_n=2^{\,n}$. $2^{\,n}=1024$ geeft $n=10$. $2^{\,n}=1000$ geeft $n\approx 9{,}97$: geen natuurlijk getal.
\end{fullsolution}
\end{oefening}

\begin{oefening}\label{oef:t8}
Een kaars is 30 cm lang en brandt 1,5 cm per uur op. Geef een voorschrift voor de lengte na $n$ uur. Wanneer is de kaars op?
\begin{solution}
$u_n=30-1{,}5\,n$; na 20 uur
\end{solution}
\begin{fullsolution}
$u_0=30$ en $v=-1{,}5$, dus $u_n=30-1{,}5\,n$. Uit $u_n=0$ volgt $n=20$.
\end{fullsolution}
\end{oefening}

\begin{oefening}\label{oef:t9}
Je belegt \eur 4000 aan 2\% samengestelde intrest. Hoeveel heb je na 5 jaar?
\begin{solution}
Ongeveer \eur 4416,32
\end{solution}
\begin{fullsolution}
$K_5=4000\cdot 1{,}02^5\approx 4416{,}32$.
\end{fullsolution}
\end{oefening}
\end{oefeningen}

\herhaalschema{%
  \fzv{oef:t1}{oef:soort}%
  \fv{oef:t2}{vb:term}{oef:voorschrift}%
  \fv{oef:t3}{vb:mterm}{oef:mvoorschrift}%
  \fv{oef:t4}{vb:tweetermen}{oef:tweetermen}%
  \fv{oef:t5}{vb:mtwee}{oef:mtwee}%
  \fv{oef:t6}{vb:isterm}{oef:isterm}%
  \fv{oef:t7}{vb:misterm}{oef:misterm}%
  \fv{oef:t8}{vb:groei}{oef:groei}%
  \fv{oef:t9}{vb:intrest}{oef:intrest}}

\ifzelfstudie
\section*{Maak je eigen samenvatting}
Sluit dit boek en schrijf in je schrift: wat is een rekenkundige en wat is een meetkundige rij, welke voorschriften horen erbij, hoe vind je de rij uit twee termen, hoe test je of een getal een term is, wat is het verschil tussen enkelvoudige en samengestelde intrest, en welke fouten wil je vermijden? Vergelijk daarna met de kaders.
\fi

%----------------------------------------------------------------------------------------
% SOMFORMULES. Alle opgaven zijn eigen werk.
%
% Opbouw: Start, leerstof (kern, uitgewerkt voorbeeld, korte oefening), en daarna op een
% nieuwe pagina de Opdrachten (door elkaar, stijgende moeilijkheid), Koppel en Test jezelf.
% Labels van dit hoofdstuk beginnen met h3-.
%----------------------------------------------------------------------------------------

\mychapter{3}{Somformules}

\lkr{Leerplan B+S'': LPD 44, formule voor de som van de eerste $n$ termen, met samengestelde intrest als context. Twee bewijzen die leerlingen moeten kunnen volgen en navertellen: de som tweemaal schrijven, en $q \cdot s_n - s_n$. De rode draad in de opdrachten: eerst herkennen welke soort rij het is, dan het aantal termen bepalen, dan pas de formule gebruiken. Oneindige sommen komen niet in dit hoofdstuk: die vragen de limiet uit hoofdstuk 4.}

\begin{oefeningen}[Start]
\begin{oefening}
\blokkop{Opwarmer}
\begin{enumerate}
  \item Een rekenkundige rij heeft $u_1 = 4$ en $v = 3$. Bereken $u_{10}$.
  \item Een meetkundige rij heeft $u_1 = 2$ en $q = 3$. Bereken $u_5$.
  \item Hoeveel termen telt de rij $8,\ 11,\ 14,\ \ldots,\ 77$?
  \item Los op: $n^2 + 3n - 40 = 0$.
  \item Werk uit: $(x-1)(x^3 + x^2 + x + 1)$.
  \item Los op: $2^x = 1024$ en $3^x = 200$.
\end{enumerate}
\begin{solution}
\textbf{a.}~31 \quad \textbf{b.}~162 \quad \textbf{c.}~24 \quad \textbf{d.}~$n=5$ of $n=-8$ \quad \textbf{e.}~$x^4-1$ \quad \textbf{f.}~$x=10$; $x\approx 4{,}82$
\end{solution}
\begin{fullsolution}
\textbf{a.}~$u_{10}=4+9\cdot 3=31$.\\
\textbf{b.}~$u_5=2\cdot 3^4=162$.\\
\textbf{c.}~$8+(n-1)\cdot 3=77$ geeft $n-1=23$, dus $n=24$.\\
\textbf{d.}~$(n-5)(n+8)=0$.\\
\textbf{e.}~$x^4+x^3+x^2+x-x^3-x^2-x-1=x^4-1$: bijna alles valt weg.\\
\textbf{f.}~$2^{10}=1024$, dus $x=10$. $x=\log_3 200=\dfrac{\log 200}{\log 3}\approx 4{,}82$.
\end{fullsolution}
\end{oefening}

\begin{oefening}\label{oef:h3-verken}
\blokkop{Verken}\quad Je spaart tien weken lang. In week 1 leg je \eur 3 opzij, en elke week \eur 2 meer dan de week ervoor.
\begin{enumerate}
  \item Schrijf de tien bedragen op een rij. Schrijf daaronder dezelfde bedragen in omgekeerde volgorde.
  \item Tel per kolom de twee getallen op. Wat valt op?
  \item Hoeveel heb je na tien weken gespaard? Gebruik de twee rijen, zonder alle bedragen op te tellen.
\end{enumerate}
\begin{solution}
\textbf{b.}~elke kolom geeft 24 \quad \textbf{c.}~\eur 120
\end{solution}
\begin{fullsolution}
\textbf{a.}~$3,\ 5,\ 7,\ 9,\ 11,\ 13,\ 15,\ 17,\ 19,\ 21$ en daaronder $21,\ 19,\ \ldots,\ 3$.\\
\textbf{b.}~Bovenaan komt er telkens 2 bij, onderaan gaat er telkens 2 af: elke kolom geeft 24.\\
\textbf{c.}~De twee rijen samen zijn $10\cdot 24=240$. Dat is twee keer het gespaarde bedrag, dus \eur 120.
\end{fullsolution}
\end{oefening}
\end{oefeningen}

\lkr{Laat oefening \ref{oef:h3-verken} maken vóór de uitleg. Vraag waarom je deelt door 2: dat is de stap die leerlingen het vaakst vergeten als ze de formule later moeten verklaren.}

%========================================================================================
\section{Som van een rekenkundige rij}

\subsection{De somformule}

\begin{theorie}
\begin{minipage}[t]{.64\linewidth}
$s_n$ is de som van de eerste $n$ termen:

\qquad $\formule{s_n = u_1 + u_2 + \cdots + u_n}$

\vspace{1.5mm}
Bij een rekenkundige rij is dat $n$ keer het gemiddelde van de eerste en de laatste term:

\vspace{1mm}
\qquad $\formule{s_n = n \cdot \dfrac{u_1 + u_n}{2}}$

\vspace{1.5mm}
De figuur toont het idee: twee keer dezelfde rij vormt een rechthoek van $n$ op $u_1 + u_n$.
\end{minipage}\hfill
\begin{minipage}[t]{.32\linewidth}
\vspace{-1mm}
\begin{tikzpicture}[x=5.4mm,y=2.7mm,font=\scriptsize]
  \foreach \i/\h in {0/2,1/3,2/4,3/5,4/6}{
    \fill[accent!45] (\i,0) rectangle (\i+1,\h);
    \draw[accent] (\i,0) rectangle (\i+1,\h);
    \draw[accent] (\i,\h) rectangle (\i+1,8);}
  \draw[decorate,decoration={brace,mirror,amplitude=3pt},grijs] (0,-.3) -- (5,-.3) node[midway,below=2pt,text=black] {$n$ termen};
  \draw[decorate,decoration={brace,mirror,amplitude=3pt},grijs] (5.15,0) -- (5.15,8) node[midway,right=2pt,text=black] {$u_1 + u_n$};
\end{tikzpicture}
\end{minipage}

\begin{bewijs}
\begin{stappen}
s_n = u_1 + u_2 + \cdots + u_n & De som. \\
s_n = u_n + u_{n-1} + \cdots + u_1 & Dezelfde som, omgekeerd. \\
2 s_n = (u_1 + u_n) + \cdots + (u_n + u_1) & Optellen per kolom. Bovenaan komt er telkens $v$ bij, onderaan gaat er telkens $v$ af: elke kolom geeft $u_1 + u_n$. \\
2 s_n = n \,(u_1 + u_n) & Er zijn $n$ kolommen. \\[1ex]
s_n = n \cdot \dfrac{u_1 + u_n}{2} & Delen door 2.
\end{stappen}
\end{bewijs}
\end{theorie}

\voorbeeld{}{Bereken $s_{20}$ van de rij $7,\ 11,\ 15,\ \ldots$\label{vb:h3-rr}}{%
\begin{stappen}
v = 11 - 7 = 4 & Rekenkundig. Controle: $15 - 11 = 4$. \\
u_{20} = 7 + 19 \cdot 4 = 83 & Eerst de laatste term: 19 sprongen. \\[1ex]
s_{20} = 20 \cdot \dfrac{7 + 83}{2} = 900 & Twintig keer het gemiddelde, en dat is 45.
\end{stappen}}

\nujij{Bereken $s_{30}$ van de rij $20,\ 18,\ 16,\ \ldots$}{$v = -2$, $u_{30} = -38$, $s_{30} = 30 \cdot \tfrac{20 - 38}{2} = -270$.}

\thuis{Vul je $u_n = u_1 + (n-1)v$ in, dan krijg je een tweede vorm: $s_n = \tfrac{n}{2}\bigl(2u_1 + (n-1)v\bigr)$. Die is handig als je $u_n$ nog niet kent.}

\begin{oefeningen}
\begin{oefening}
Bereken $s_{25}$ van de rij $9,\ 13,\ 17,\ \ldots$
\begin{solution}
1425
\end{solution}
\begin{fullsolution}
$v=4$, dus $u_{25}=9+24\cdot 4=105$ en $s_{25}=25\cdot\dfrac{9+105}{2}=1425$.
\end{fullsolution}
\end{oefening}
\end{oefeningen}

\subsection{Eerst het aantal termen}

\begin{werkwijze}
Is de laatste term gegeven en niet het aantal termen, bereken dan eerst $n$ met het expliciet voorschrift. Het aantal termen is één meer dan het aantal sprongen.
\end{werkwijze}

\voorbeeld{}{Bereken de som van alle veelvouden van 6 tussen 100 en 400.\label{vb:h3-aantal}}{%
\begin{stappen}
102 + 108 + \cdots + 396 & Eerst uitschrijven: $102 = 6 \cdot 17$ en $396 = 6 \cdot 66$. \\
102 + (n-1) \cdot 6 = 396 \ \Rightarrow \ n = 50 & 49 sprongen, dus 50 termen. \\[1ex]
s_{50} = 50 \cdot \dfrac{102 + 396}{2} = 12\,450 &
\end{stappen}}

\nujij{Bereken $50 + 54 + 58 + \cdots + 250$.}{$n = 51$, en de som is $51 \cdot 150 = 7650$.}

\begin{oefeningen}
\begin{oefening}
Bereken $15 + 22 + 29 + \cdots + 204$.
\begin{solution}
3066
\end{solution}
\begin{fullsolution}
$15+(n-1)\cdot 7=204$ geeft $n-1=27$, dus $n=28$. Dan is de som $28\cdot\dfrac{15+204}{2}=3066$.
\end{fullsolution}
\end{oefening}
\end{oefeningen}

\subsection{Het aantal termen uit de som}

\begin{werkwijze}
Is de som gegeven, schrijf dan $s_n$ met $n$ als onbekende. Dat geeft een vierkantsvergelijking. Enkel een oplossing in $\mathbb{N}_0$ is een aantal termen.
\end{werkwijze}

\voorbeeld{}{Hoeveel termen van $3,\ 7,\ 11,\ \ldots$ moet je optellen om 210 te krijgen?\label{vb:h3-nuit}}{%
\begin{stappen}
u_n = 3 + (n-1) \cdot 4 = 4n - 1 & De laatste term, met $n$ als onbekende. \\[1ex]
s_n = n \cdot \dfrac{3 + 4n - 1}{2} = n\,(2n + 1) & Invullen in de somformule. \\
2n^2 + n - 210 = 0 & $s_n = 210$ stellen. \\[1ex]
n = 10 \quad\text{of}\quad n = -\dfrac{21}{2} & Enkel $n = 10$ kan. Controle: $10 \cdot 21 = 210$.
\end{stappen}}

\nujij{Hoeveel termen van $6,\ 8,\ 10,\ \ldots$ moet je optellen om 150 te krijgen?}{$s_n = n(n+5)$, dus $n = 10$.}

\begin{oefeningen}
\begin{oefening}
Hoeveel termen van $4,\ 7,\ 10,\ \ldots$ moet je optellen om 175 te krijgen?
\begin{solution}
10
\end{solution}
\begin{fullsolution}
$u_n=3n+1$, dus $s_n=n\cdot\dfrac{4+3n+1}{2}=\dfrac{n(3n+5)}{2}$. Uit $3n^2+5n-350=0$ volgt $n=10$ of $n=-\tfrac{35}{3}$. Enkel $n=10$ kan.
\end{fullsolution}
\end{oefening}
\end{oefeningen}

%========================================================================================
\section{Som van een meetkundige rij}

\begin{oefeningen}
\begin{oefening}\label{oef:h3-mverken}
\blokkop{Verken}\quad Bij sneeuw verwittigt de school 1 leerling. Die belt 3 andere leerlingen. Elk van hen belt er weer 3, enzovoort. Noem $S = 1 + 3 + 9 + 27 + 81$ het aantal verwittigde leerlingen na vier belrondes.
\begin{enumerate}
  \item Schrijf ook $3S$ uit, als som van vijf termen.
  \item Bereken $3S - S$. Wat valt weg?
  \item Bereken daarmee $S$ zonder de termen op te tellen.
\end{enumerate}
\begin{solution}
\textbf{a.}~$3+9+27+81+243$ \quad \textbf{b.}~$243-1=242$ \quad \textbf{c.}~121
\end{solution}
\begin{fullsolution}
\textbf{a.}~Elke term maal 3: $3S=3+9+27+81+243$.\\
\textbf{b.}~De termen 3, 9, 27 en 81 staan in beide sommen en vallen weg: $3S-S=243-1=242$.\\
\textbf{c.}~$2S=242$, dus $S=121$.
\end{fullsolution}
\end{oefening}
\end{oefeningen}

\lkr{Laat de leerlingen vergelijken met de spaaroefening: bij een rekenkundige rij kies je voor optellen in omgekeerde volgorde, bij een meetkundige rij voor vermenigvuldigen met $q$. Vraag waarom de truc van de ene niet werkt bij de andere.}

\subsection{De somformule}

\begin{theorie}
Voor een meetkundige rij met $q \neq 1$: \qquad $\formule{s_n = u_1 \cdot \dfrac{q^{\,n} - 1}{q - 1}}$

\vspace{1.5mm}
Voor $q = 1$ zijn alle termen gelijk: \quad $\formule{s_n = n \cdot u_1}$

\begin{bewijs}
\begin{stappen}
s_n = u_1 + u_1 q + \cdots + u_1 q^{\,n-1} & De som. \\
q\, s_n = u_1 q + \cdots + u_1 q^{\,n-1} + u_1 q^{\,n} & Vermenigvuldig met $q$: alles schuift een plaats op. \\
q\, s_n - s_n = u_1 q^{\,n} - u_1 & Aftrekken: de termen ertussen vallen weg. \\
s_n (q - 1) = u_1 (q^{\,n} - 1) & Links en rechts buiten haakjes brengen. \\[1ex]
s_n = u_1 \cdot \dfrac{q^{\,n} - 1}{q - 1} & Delen door $q - 1$, wat mag omdat $q \neq 1$.
\end{stappen}
\end{bewijs}
\end{theorie}

Let op: in de somformule staat de exponent $n$, het aantal termen. In het expliciet voorschrift staat $n - 1$, het aantal sprongen.

\voorbeeld{}{Bereken $s_{10}$ van de rij $3,\ 6,\ 12,\ \ldots$\label{vb:h3-mr}}{%
\begin{stappen}
q = \dfrac{6}{3} = 2 & Meetkundig. Controle: $\dfrac{12}{6} = 2$. \\[1.5ex]
s_{10} = 3 \cdot \dfrac{2^{10} - 1}{2 - 1} = 3069 & Exponent 10, want er zijn 10 termen.
\end{stappen}}

\nujij{Bereken $s_9$ van de rij $1,\ {-2},\ 4,\ \ldots$}{$q = -2$, dus $s_9 = \tfrac{(-2)^9 - 1}{-2 - 1} = \tfrac{-513}{-3} = 171$.}

\begin{oefeningen}
\begin{oefening}
Bereken $s_6$ van de rij $2,\ 10,\ 50,\ \ldots$
\begin{solution}
7812
\end{solution}
\begin{fullsolution}
$q=5$, dus $s_6=2\cdot\dfrac{5^6-1}{5-1}=2\cdot\dfrac{15\,624}{4}=7812$.
\end{fullsolution}
\end{oefening}
\end{oefeningen}

\subsection{Eerst het aantal termen}

\begin{werkwijze}
Is de laatste term gegeven, bereken dan eerst $n$ uit $u_n = u_1 \cdot q^{\,n-1}$, zo nodig met een logaritme. Werk met de exacte machten, niet met afgeronde getallen.
\end{werkwijze}

\voorbeeld{}{Bereken $2 + 6 + 18 + \cdots + 4374$.\label{vb:h3-mrn}}{%
\begin{stappen}
2 \cdot 3^{\,n-1} = 4374 \ \Rightarrow \ 3^{\,n-1} = 2187 & $q = 3$. Delen door 2. \\
n - 1 = 7 \ \Rightarrow \ n = 8 & Want $3^7 = 2187$. Dus 8 termen. \\[1ex]
s_8 = 2 \cdot \dfrac{3^8 - 1}{3 - 1} = 6560 &
\end{stappen}}

\nujij{Bereken $7 + 14 + 28 + \cdots + 448$.}{$2^{n-1} = 64$, dus $n = 7$, en de som is $7 \cdot (2^7 - 1) = 889$.}

\begin{oefeningen}
\begin{oefening}
Bereken $5 + 15 + 45 + \cdots + 3645$.
\begin{solution}
5465
\end{solution}
\begin{fullsolution}
$5\cdot 3^{n-1}=3645$ geeft $3^{n-1}=729=3^6$, dus $n=7$. Dan is $s_7=5\cdot\dfrac{3^7-1}{3-1}=5465$.
\end{fullsolution}
\end{oefening}
\end{oefeningen}

%========================================================================================
\section{Het somteken}

\subsection{Lezen en berekenen}

\begin{theorie}
\begin{minipage}[t]{.56\linewidth}
Een som schrijf je kort met het \emph{somteken} $\Sigma$:

\vspace{1mm}
\qquad $\formule{\displaystyle\sum_{i=1}^{n} u_i = u_1 + u_2 + \cdots + u_n}$

\vspace{1.5mm}
De letter $i$ is de \emph{lopende index}: die neemt elk geheel getal aan van de beginwaarde tot en met de eindwaarde.
\end{minipage}\hfill
\begin{minipage}[t]{.4\linewidth}
\textbf{Aantal termen:} eindwaarde $-$ beginwaarde $+ 1$

\vspace{1.5mm}
\textbf{Rekenregels:}

$\formule{\displaystyle\sum (u_i + w_i) = \sum u_i + \sum w_i}$

$\formule{\displaystyle\sum c \cdot u_i = c \cdot \sum u_i}$
\end{minipage}
\end{theorie}

\thuis{Je leest $\sum_{i=1}^{n} u_i$ als: de som van de $u_i$ voor $i$ van 1 tot $n$. Het teken is de Griekse hoofdletter sigma, de S van \emph{som}. De naam van de index maakt niet uit: $\sum_{i=1}^{5} i$ en $\sum_{k=1}^{5} k$ zijn hetzelfde getal.}

\voorbeeld{}{Bereken $\displaystyle\sum_{i=3}^{12} (2i + 1)$ en $\displaystyle\sum_{j=1}^{5} 3 \cdot 2^{\,j}$.\label{vb:h3-somteken}}{%
\begin{stappen}
7 + 9 + 11 + \cdots + 25 & Vul $i = 3$, 4, 5 en 12 in: rekenkundig. \\
12 - 3 + 1 = 10 \text{ termen} & Niet 12 en niet 9. \\[1ex]
10 \cdot \dfrac{7 + 25}{2} = 160 & De eerste som. \\
6 + 12 + 24 + 48 + 96 & Vul $j = 1$, 2, 3 in: meetkundig, met $u_1 = 6$ en niet 3. \\[1ex]
6 \cdot \dfrac{2^5 - 1}{2 - 1} = 186 & De tweede som.
\end{stappen}}

\lkr{Laat leerlingen altijd de eerste drie termen uitschrijven voor ze een formule kiezen. Dat voorkomt de twee klassieke fouten: het verkeerde aantal termen, en $u_1$ aflezen zonder invullen.}

\nujij{Bereken $\displaystyle\sum_{k=2}^{9} (5k - 3)$.}{Termen $7,\ 12,\ \ldots,\ 42$; 8 termen; de som is $8 \cdot \tfrac{7 + 42}{2} = 196$.}

\begin{oefeningen}
\begin{oefening}
Bereken.
\begin{reeks}[2]
  \item $\displaystyle\sum_{i=2}^{11} (3i + 1)$
  \item $\displaystyle\sum_{j=1}^{6} 2 \cdot 3^{\,j}$
\end{reeks}
\begin{solution}
\textbf{a.}~205 \quad \textbf{b.}~2184
\end{solution}
\begin{fullsolution}
\textbf{a.}~$7+10+\cdots+34$: 10 termen, dus $10\cdot\dfrac{7+34}{2}=205$.\\
\textbf{b.}~$6+18+\cdots$: $u_1=6$, $q=3$ en 6 termen, dus $6\cdot\dfrac{3^6-1}{3-1}=2184$.
\end{fullsolution}
\end{oefening}
\end{oefeningen}

\subsection{De som als rij}

\begin{theorie}
De sommen $s_1,\ s_2,\ s_3,\ \ldots$ vormen zelf een rij. Daaruit vind je de termen terug:

\vspace{1mm}
\qquad $\formule{u_1 = s_1} \quad\text{en}\quad \formule{u_n = s_n - s_{n-1}} \quad\text{voor } n \geqslant 2$
\end{theorie}

\voorbeeld{}{Gegeven: $s_n = n^2 + 2n$. Bepaal $u_n$. Welke soort rij is dat?\label{vb:h3-uitsom}}{%
\begin{stappen}
u_1 = s_1 = 3 & Vervang $n$ door 1. \\
s_{n-1} = (n-1)^2 + 2(n-1) = n^2 - 1 & Vervang $n$ door $n - 1$. \\
u_n = n^2 + 2n - (n^2 - 1) = 2n + 1 & Voor $n \geqslant 2$. Ook $u_1 = 3$ past in dit voorschrift. \\
\text{rekenkundig met } v = 2 & Controle: $3 + 5 + 7 = 15 = s_3$.
\end{stappen}}

\begin{oefeningen}
\begin{oefening}
Gegeven: $s_n = n^2 + 4n$. Bepaal $u_n$.
\begin{solution}
$u_n=2n+3$
\end{solution}
\begin{fullsolution}
$s_{n-1}=(n-1)^2+4(n-1)=n^2+2n-3$, dus $u_n=s_n-s_{n-1}=2n+3$ voor $n\geqslant 2$. Ook $u_1=s_1=5$ past hierin.
\end{fullsolution}
\end{oefening}
\end{oefeningen}

%========================================================================================
\section{Sparen met vaste stortingen}

Je stort aan het einde van elk jaar een bedrag $a$ op een rekening met $p\%$ samengestelde intrest per jaar. Elke storting groeit apart aan, met $q = 1 + \tfrac{p}{100}$. De laatste storting brengt nog niets op.

\vspace{2mm}
\begin{tikzpicture}[x=2.25cm,font=\small]
  \draw[->,lijngrijs] (.55,0) -- (5.6,0) node[right,text=black] {tijd};
  \foreach \x/\t/\w in {1/{jaar 1}/{a\,q^{\,n-1}},2/{jaar 2}/{a\,q^{\,n-2}},4/{jaar $n-1$}/{a\,q},5/{jaar $n$}/{a}}{
    \fill[accent] (\x,0) circle (2pt);
    \node[above=1pt,font=\scriptsize,text=grijs] at (\x,0) {\t};
    \node[below=2pt] at (\x,0) {$\w$};}
  \node[fill=white,inner sep=1pt] at (3,0) {$\cdots$};
\end{tikzpicture}

\begin{theorie}
De eindwaarde na $n$ jaar is de som van een meetkundige rij met eerste term $a$: \quad $\formule{E = a \cdot \dfrac{q^{\,n} - 1}{q - 1}}$
\end{theorie}

\lkr{Dit deel past het best bij Economie-Wiskunde. Het leerplan noemt enkel samengestelde intrest als context: sla de opdrachten over sparen over als je dit deel niet behandelt. Leningen komen alleen voor in een opdracht op niveau 3.}

\voorbeeld{}{Je stort 6 jaar lang aan het einde van elk jaar \eur 1000 op een rekening met 3\% intrest. Hoeveel staat er na de laatste storting op? Hoeveel daarvan is intrest?\label{vb:h3-sparen}}{%
\begin{stappen}
q = 1 + \dfrac{3}{100} = 1{,}03 & \\[1.5ex]
E = 1000 \cdot \dfrac{1{,}03^6 - 1}{1{,}03 - 1} \approx 6468{,}41 & Na 6 jaar: ongeveer \textbf{\eur 6468,41}. \\[1.5ex]
6468{,}41 - 6 \cdot 1000 = 468{,}41 & \textbf{\eur 468,41} is intrest.
\end{stappen}}

\nujij{Je stort 10 jaar lang aan het einde van elk jaar \eur 500 op een rekening met 2\% intrest. Hoeveel staat er na de laatste storting op?}{$500 \cdot \tfrac{1{,}02^{10} - 1}{0{,}02} \approx 5474{,}86$ euro.}

\begin{oefeningen}
\begin{oefening}
Je stort 8 jaar lang aan het einde van elk jaar \eur 1500 op een rekening met 2\% intrest. Hoeveel staat er na de laatste storting op?
\begin{solution}
Ongeveer \eur 12\,874,45
\end{solution}
\begin{fullsolution}
$E=1500\cdot\dfrac{1{,}02^8-1}{0{,}02}\approx 12\,874{,}45$.
\end{fullsolution}
\end{oefening}
\end{oefeningen}

%----------------------------------------------------------------------------------------
% OPDRACHTEN: door elkaar. Niveau 1, dan 2, dan 3.
%----------------------------------------------------------------------------------------
\opdrachtenkop

\lkr{De opdrachten staan bewust door elkaar, met stijgende moeilijkheid. In deze versie staat bij elke opdracht het onderwerp, zodat je kunt kiezen.}

\begin{oefeningen}
\begin{oefening}[1]\label{oef:h3-soort}
\thema{gemengd: soort rij herkennen}Is de rij rekenkundig of meetkundig? Bereken dan de som.
\begin{enumerate}
  \item $4 + 9 + 14 + \cdots$ (12 termen)
  \item $3 + 6 + 12 + \cdots$ (7 termen)
  \item $20 + 17 + 14 + \cdots$ (10 termen)
  \item $2 - 6 + 18 - \cdots$ (6 termen)
\end{enumerate}
\begin{solution}
\textbf{a.}~rekenkundig, 378 \quad \textbf{b.}~meetkundig, 381 \quad \textbf{c.}~rekenkundig, 65 \quad \textbf{d.}~meetkundig, $-364$
\end{solution}
\begin{fullsolution}
\textbf{a.}~$v=5$, $u_{12}=59$, dus $12\cdot\dfrac{4+59}{2}=378$.\\
\textbf{b.}~$q=2$, dus $3\cdot(2^7-1)=381$.\\
\textbf{c.}~$v=-3$, $u_{10}=-7$, dus $10\cdot\dfrac{20-7}{2}=65$.\\
\textbf{d.}~$q=-3$, dus $2\cdot\dfrac{(-3)^6-1}{-3-1}=2\cdot\dfrac{728}{-4}=-364$.
\end{fullsolution}
\end{oefening}

\begin{oefening}[1]\label{oef:h3-rrsom}
\thema{rekenkundig: som}Bereken de som van de rekenkundige rij.
\begin{enumerate}
  \item $s_{15}$ als $u_1 = 8$ en $v = 3$
  \item $s_{11}$ als $u_1 = 40$ en $v = -6$
  \item $s_9$ als $u_1 = \tfrac12$ en $v = \tfrac14$
\end{enumerate}
\begin{solution}
\textbf{a.}~435 \quad \textbf{b.}~110 \quad \textbf{c.}~$\tfrac{27}{2}$
\end{solution}
\begin{fullsolution}
\textbf{a.}~$u_{15}=8+14\cdot 3=50$, dus $15\cdot\dfrac{8+50}{2}=435$.\\
\textbf{b.}~$u_{11}=40-60=-20$, dus $11\cdot\dfrac{40-20}{2}=110$.\\
\textbf{c.}~$u_9=\tfrac12+8\cdot\tfrac14=\tfrac52$, dus $9\cdot\dfrac{\tfrac12+\tfrac52}{2}=\tfrac{27}{2}$.
\end{fullsolution}
\end{oefening}

\begin{oefening}[1]\label{oef:h3-sigma}
\thema{somteken: uitschrijven}Schrijf uit en bereken.
\begin{reeks}[4]
  \item $\displaystyle\sum_{i=1}^{4} (2i + 1)$
  \item $\displaystyle\sum_{k=2}^{5} k^2$
  \item $\displaystyle\sum_{j=0}^{3} 2^{\,j}$
  \item $\displaystyle\sum_{i=1}^{5} (-1)^i$
\end{reeks}
\begin{solution}
\textbf{a.}~24 \quad \textbf{b.}~54 \quad \textbf{c.}~15 \quad \textbf{d.}~$-1$
\end{solution}
\begin{fullsolution}
\textbf{a.}~$3+5+7+9=24$.\\
\textbf{b.}~$4+9+16+25=54$.\\
\textbf{c.}~$1+2+4+8=15$.\\
\textbf{d.}~$-1+1-1+1-1=-1$.
\end{fullsolution}
\end{oefening}

\begin{oefening}[1]\label{oef:h3-mrsom}
\thema{meetkundig: som}Bereken de som van de meetkundige rij.
\begin{enumerate}
  \item $s_8$ als $u_1 = 5$ en $q = 2$
  \item $s_5$ als $u_1 = 81$ en $q = \tfrac13$
  \item $s_5$ als $u_1 = -2$ en $q = -3$
  \item $s_{20}$ als $u_1 = 7$ en $q = 1$
\end{enumerate}
\begin{solution}
\textbf{a.}~1275 \quad \textbf{b.}~121 \quad \textbf{c.}~$-122$ \quad \textbf{d.}~140
\end{solution}
\begin{fullsolution}
\textbf{a.}~$5\cdot(2^8-1)=1275$.\\
\textbf{b.}~$81\cdot\dfrac{(\tfrac13)^5-1}{\tfrac13-1}=81\cdot\dfrac{-\tfrac{242}{243}}{-\tfrac23}=121$.\\
\textbf{c.}~$-2\cdot\dfrac{(-3)^5-1}{-3-1}=-2\cdot\dfrac{-244}{-4}=-122$.\\
\textbf{d.}~$q=1$: alle termen zijn 7, dus $20\cdot 7=140$.
\end{fullsolution}
\end{oefening}

\begin{oefening}[1]\label{oef:h3-aantal}
\thema{aantal termen}Bereken de som.
\begin{enumerate}
  \item $11 + 14 + 17 + \cdots + 95$
  \item $6 + 12 + 24 + \cdots + 1536$
  \item $100 + 96 + 92 + \cdots + 4$
\end{enumerate}
\begin{solution}
\textbf{a.}~1537 \quad \textbf{b.}~3066 \quad \textbf{c.}~1300
\end{solution}
\begin{fullsolution}
\textbf{a.}~$v=3$: $11+(n-1)\cdot 3=95$ geeft $n=29$. Som: $29\cdot\dfrac{11+95}{2}=1537$.\\
\textbf{b.}~$q=2$: $6\cdot 2^{n-1}=1536$ geeft $2^{n-1}=256$, dus $n=9$. Som: $6\cdot(2^9-1)=3066$.\\
\textbf{c.}~$v=-4$: $100+(n-1)\cdot(-4)=4$ geeft $n=25$. Som: $25\cdot\dfrac{100+4}{2}=1300$.
\end{fullsolution}
\end{oefening}

\begin{oefening}[1]
\thema{fouten: exponent}Iemand schrijft: `$u_1 = 3$ en $q = 2$, dus $s_5 = 3 \cdot \dfrac{2^4 - 1}{2 - 1} = 45$.' Wat is er fout? Geef het juiste antwoord.
\begin{solution}
De exponent is $n$, niet $n-1$: $s_5=93$.
\end{solution}
\begin{fullsolution}
In de somformule staat de exponent $n$: $s_5=3\cdot(2^5-1)=93$. Controle: $3+6+12+24+48=93$.
\end{fullsolution}
\end{oefening}

\begin{oefening}[1]\label{oef:h3-noteer}
\thema{somteken: noteren}Noteer met het somteken.
\begin{enumerate}
  \item $5 + 10 + 15 + \cdots + 200$
  \item $1 + 4 + 9 + \cdots + 144$
  \item $2 + 4 + 8 + \cdots + 1024$
  \item $1 - 2 + 3 - 4 + \cdots - 20$
\end{enumerate}
\begin{solution}
\textbf{a.}~$\displaystyle\sum_{i=1}^{40} 5i$ \quad \textbf{b.}~$\displaystyle\sum_{i=1}^{12} i^2$ \quad \textbf{c.}~$\displaystyle\sum_{i=1}^{10} 2^{\,i}$ \quad \textbf{d.}~$\displaystyle\sum_{i=1}^{20} (-1)^{i+1} i$
\end{solution}
\begin{fullsolution}
\textbf{a.}~De veelvouden van 5, en $200=5\cdot 40$: $\displaystyle\sum_{i=1}^{40} 5i$.\\
\textbf{b.}~De kwadraten, en $144=12^2$: $\displaystyle\sum_{i=1}^{12} i^2$.\\
\textbf{c.}~De machten van 2, en $1024=2^{10}$: $\displaystyle\sum_{i=1}^{10} 2^{\,i}$.\\
\textbf{d.}~Positief bij oneven $i$: $\displaystyle\sum_{i=1}^{20} (-1)^{i+1} i$.\\
Andere schrijfwijzen kunnen ook, bijvoorbeeld met een andere beginwaarde.
\end{fullsolution}
\end{oefening}

\begin{oefening}[1]
\thema{fouten: aantal termen}Iemand schrijft: `$12 + 15 + \cdots + 60$ telt $\dfrac{60 - 12}{3} = 16$ termen.' Wat is er fout? Bereken de som.
\begin{solution}
Dat zijn 16 sprongen, dus 17 termen; de som is 612.
\end{solution}
\begin{fullsolution}
$\dfrac{60-12}{3}=16$ is het aantal sprongen. Er zijn dus 17 termen, en de som is $17\cdot\dfrac{12+60}{2}=612$.
\end{fullsolution}
\end{oefening}

\begin{oefening}[1]
\begin{metprent}{fiets}
\thema{context: rekenkundig}Een wielrenner fietst op de eerste dag van zijn trainingsperiode 20 km. Elke dag fietst hij 3 km meer dan de dag ervoor. Hoeveel kilometer fietst hij in 30 dagen?
\end{metprent}
\begin{solution}
1905 km
\end{solution}
\begin{fullsolution}
$u_{30}=20+29\cdot 3=107$, dus $s_{30}=30\cdot\dfrac{20+107}{2}=1905$.
\end{fullsolution}
\end{oefening}

\begin{oefening}[1]
\thema{de som als rij}Van een rij is $s_1 = 4$, $s_2 = 10$ en $s_3 = 18$. Bepaal $u_1$, $u_2$ en $u_3$.
\begin{solution}
4, 6, 8
\end{solution}
\begin{fullsolution}
$u_1=s_1=4$, $u_2=s_2-s_1=6$ en $u_3=s_3-s_2=8$.
\end{fullsolution}
\end{oefening}

\begin{oefening}[1]
\thema{somteken: berekenen}Bereken.
\begin{reeks}[2]
  \item $\displaystyle\sum_{i=1}^{30} (4i - 1)$
  \item $\displaystyle\sum_{i=1}^{6} 3 \cdot 2^{\,i}$
\end{reeks}
\begin{solution}
\textbf{a.}~1830 \quad \textbf{b.}~378
\end{solution}
\begin{fullsolution}
\textbf{a.}~$3+7+\cdots+119$: 30 termen, dus $30\cdot\dfrac{3+119}{2}=1830$.\\
\textbf{b.}~$6+12+\cdots+192$: $u_1=6$, $q=2$, dus $6\cdot(2^6-1)=378$.
\end{fullsolution}
\end{oefening}

\begin{oefening}[1]
\thema{rekenkundig: voorbeeld uit hoofdstuk 1}In hoofdstuk 1 telde trap $n$ precies $1 + 2 + \cdots + n$ blokjes. Toon met de somformule aan dat dit $\dfrac{n(n+1)}{2}$ is.
\begin{solution}
$u_1=1$, $u_n=n$, $n$ termen.
\end{solution}
\begin{fullsolution}
Dit is een rekenkundige rij met $u_1=1$, $u_n=n$ en $n$ termen, dus $s_n=n\cdot\dfrac{1+n}{2}=\dfrac{n(n+1)}{2}$.
\end{fullsolution}
\end{oefening}

\begin{oefening}[2]\label{oef:h3-somindex}
\thema{somteken: andere beginwaarde}Bereken.
\begin{reeks}
  \item $\displaystyle\sum_{k=5}^{24} (3k + 2)$
  \item $\displaystyle\sum_{j=3}^{8} 2^{\,j}$
  \item $\displaystyle\sum_{i=0}^{9} (10 - i)$
\end{reeks}
\begin{solution}
\textbf{a.}~910 \quad \textbf{b.}~504 \quad \textbf{c.}~55
\end{solution}
\begin{fullsolution}
\textbf{a.}~$17+20+\cdots+74$: $24-5+1=20$ termen, dus $20\cdot\dfrac{17+74}{2}=910$.\\
\textbf{b.}~$8+16+\cdots+256$: $8-3+1=6$ termen, dus $8\cdot(2^6-1)=504$.\\
\textbf{c.}~$10+9+\cdots+1$: $9-0+1=10$ termen, dus $10\cdot\dfrac{10+1}{2}=55$.
\end{fullsolution}
\end{oefening}

\begin{oefening}[2]\label{oef:h3-nuitsom}
\thema{aantal termen uit de som}Hoeveel termen van $2,\ 5,\ 8,\ \ldots$ moet je optellen om 155 te krijgen?
\begin{solution}
10
\end{solution}
\begin{fullsolution}
$u_n=3n-1$, dus $s_n=n\cdot\dfrac{2+3n-1}{2}=\dfrac{n(3n+1)}{2}$. Uit $3n^2+n-310=0$ volgt $n=10$ of $n=-\tfrac{31}{3}$. Enkel $n=10$ kan.
\end{fullsolution}
\end{oefening}

\begin{oefening}[2]\label{oef:h3-mrcontext}
\thema{context: meetkundig}Een filmpje krijgt op de eerste dag 200 weergaven. Elke dag zijn er 1,5 keer zoveel weergaven als de dag ervoor.
\begin{enumerate}
  \item Hoeveel weergaven zijn er in totaal na 10 dagen?
  \item Na hoeveel dagen is het totaal voor het eerst meer dan 100\,000?
\end{enumerate}
\begin{solution}
\textbf{a.}~ongeveer 22\,666 \quad \textbf{b.}~na 14 dagen
\end{solution}
\begin{fullsolution}
\textbf{a.}~$s_{10}=200\cdot\dfrac{1{,}5^{10}-1}{0{,}5}\approx 22\,666$.\\
\textbf{b.}~$400\,(1{,}5^{\,n}-1)>100\,000$ geeft $1{,}5^{\,n}>251$ en $n>13{,}63$, dus na 14 dagen.
\end{fullsolution}
\end{oefening}

\begin{oefening}[2]\label{oef:h3-uitsom}
\thema{de som als rij}Bepaal $u_n$. Is de rij rekenkundig, meetkundig of geen van beide?
\begin{enumerate}
  \item $s_n = 2n^2 + n$
  \item $s_n = 3^n - 1$
  \item $s_n = n^2 + 1$
\end{enumerate}
\begin{solution}
\textbf{a.}~$u_n=4n-1$, rekenkundig \quad \textbf{b.}~$u_n=2\cdot 3^{n-1}$, meetkundig \quad \textbf{c.}~$u_1=2$ en $u_n=2n-1$ voor $n\geqslant 2$: geen van beide
\end{solution}
\begin{fullsolution}
\textbf{a.}~$s_{n-1}=2n^2-3n+1$, dus $u_n=4n-1$. Ook $u_1=3$ past hierin.\\
\textbf{b.}~$u_n=3^n-3^{n-1}=3^{n-1}(3-1)=2\cdot 3^{n-1}$. Ook $u_1=2$ past hierin.\\
\textbf{c.}~$u_n=n^2+1-\bigl((n-1)^2+1\bigr)=2n-1$ voor $n\geqslant 2$. Maar $u_1=s_1=2$ past daar niet in: de rij is $2,\ 3,\ 5,\ 7,\ \ldots$
\end{fullsolution}
\end{oefening}

\begin{oefening}[2]\label{oef:h3-sparen}
\thema{sparen}Je stort 10 jaar lang aan het einde van elk jaar \eur 600 op een rekening met 2,5\% intrest.
\begin{enumerate}
  \item Hoeveel staat er na de laatste storting op?
  \item Hoeveel daarvan is intrest?
  \item Hoeveel had er gestaan als je de \eur 6000 meteen aan het begin had gestort?
\end{enumerate}
\begin{solution}
\textbf{a.}~ongeveer \eur 6722,03 \quad \textbf{b.}~ongeveer \eur 722,03 \quad \textbf{c.}~ongeveer \eur 7680,51
\end{solution}
\begin{fullsolution}
\textbf{a.}~$E=600\cdot\dfrac{1{,}025^{10}-1}{0{,}025}\approx 6722{,}03$.\\
\textbf{b.}~$6722{,}03-6000=722{,}03$.\\
\textbf{c.}~$6000\cdot 1{,}025^{10}\approx 7680{,}51$: het geld staat dan langer op de rekening.
\end{fullsolution}
\end{oefening}

\begin{oefening}[2]
\thema{rekenkundig: onbekende uit de som}Bepaal de onbekende van de rekenkundige rij.
\begin{enumerate}
  \item $v = 3$ en $s_{12} = 342$. Bepaal $u_1$.
  \item $u_1 = 5$ en $s_{10} = 185$. Bepaal $v$.
  \item $s_5 = 40$ en $s_{10} = 130$. Bepaal $u_1$ en $v$.
\end{enumerate}
\begin{solution}
\textbf{a.}~$u_1=12$ \quad \textbf{b.}~$v=3$ \quad \textbf{c.}~$u_1=4$, $v=2$
\end{solution}
\begin{fullsolution}
Gebruik $s_n=\tfrac{n}{2}\bigl(2u_1+(n-1)v\bigr)$.\\
\textbf{a.}~$6(2u_1+33)=342$ geeft $2u_1+33=57$, dus $u_1=12$.\\
\textbf{b.}~$5(10+9v)=185$ geeft $10+9v=37$, dus $v=3$.\\
\textbf{c.}~$\tfrac52(2u_1+4v)=40$ geeft $u_1+2v=8$. $5(2u_1+9v)=130$ geeft $2u_1+9v=26$. Dus $5v=10$, $v=2$ en $u_1=4$.
\end{fullsolution}
\end{oefening}

\begin{oefening}[2]
\begin{metprent}{klok}
\thema{context: rekenkundig}Een klok slaat elk heel uur het aantal uren, van 1 tot en met 12, en elk half uur één keer. Hoeveel slagen hoor je in één etmaal?
\end{metprent}
\begin{solution}
180
\end{solution}
\begin{fullsolution}
Elke twaalf uur: $1+2+\cdots+12=12\cdot\dfrac{1+12}{2}=78$ slagen. In een etmaal is dat $2\cdot 78=156$. Daar komen 24 halve uren bij: $156+24=180$.
\end{fullsolution}
\end{oefening}

\begin{oefening}[2]
\thema{meetkundig: onbekende uit de som}Bepaal de onbekende van de meetkundige rij.
\begin{enumerate}
  \item $q = 2$ en $s_6 = 189$. Bepaal $u_1$.
  \item $u_1 = 2$ en $s_3 = 14$. Bepaal $q$. Let op: er zijn twee oplossingen.
\end{enumerate}
\begin{solution}
\textbf{a.}~$u_1=3$ \quad \textbf{b.}~$q=2$ of $q=-3$
\end{solution}
\begin{fullsolution}
\textbf{a.}~$u_1\cdot(2^6-1)=63\,u_1=189$, dus $u_1=3$.\\
\textbf{b.}~$s_3=u_1+u_1q+u_1q^2$, dus $2+2q+2q^2=14$ en $q^2+q-6=0$. Dat geeft $q=2$ of $q=-3$. Controle: $2+4+8=14$ en $2-6+18=14$.
\end{fullsolution}
\end{oefening}

\begin{oefening}[2]
\thema{aantal termen: voorwaarde}Bereken de som van alle getallen van drie cijfers die bij deling door 7 als rest 2 geven.
\begin{solution}
70\,692
\end{solution}
\begin{fullsolution}
Het kleinste is $100=7\cdot 14+2$ en het grootste is $996=7\cdot 142+2$. Dat zijn $142-14+1=129$ termen, dus de som is $129\cdot\dfrac{100+996}{2}=70\,692$.
\end{fullsolution}
\end{oefening}

\begin{oefening}[2]
\thema{fouten: geen rekenkundige rij}Iemand berekent $1 + 4 + 9 + 16 + \cdots + 100$ als $10 \cdot \dfrac{1 + 100}{2} = 505$. Wat is er fout? Wat is de juiste uitkomst?
\begin{solution}
De rij is niet rekenkundig; de som is 385.
\end{solution}
\begin{fullsolution}
De verschillen 3, 5, 7 zijn niet gelijk, dus de rij is niet rekenkundig en de formule geldt niet. Optellen geeft $1+4+9+16+25+36+49+64+81+100=385$.
\end{fullsolution}
\end{oefening}

\begin{oefening}[2]
\thema{meetkundig: wisselend teken}Bereken $1 - \tfrac12 + \tfrac14 - \tfrac18 + \cdots - \tfrac{1}{128}$.
\begin{solution}
$\tfrac{85}{128}$
\end{solution}
\begin{fullsolution}
$q=-\tfrac12$ en $\tfrac{1}{128}=(\tfrac12)^7$, dus $n-1=7$ en $n=8$. De som is $\dfrac{(-\tfrac12)^8-1}{-\tfrac12-1}=\dfrac{-\tfrac{255}{256}}{-\tfrac32}=\dfrac{85}{128}$.
\end{fullsolution}
\end{oefening}

\begin{oefening}[2]
\thema{somteken: splitsen}Bereken $\displaystyle\sum_{i=1}^{10} \left(2i + 2^{\,i}\right)$.
\begin{solution}
2156
\end{solution}
\begin{fullsolution}
Splits de som: $\displaystyle\sum_{i=1}^{10} 2i=10\cdot\dfrac{2+20}{2}=110$ en $\displaystyle\sum_{i=1}^{10} 2^{\,i}=2\cdot(2^{10}-1)=2046$. Samen 2156.
\end{fullsolution}
\end{oefening}

\begin{oefening}[2]
\thema{context: meetkundig}Een bal valt van 3 m hoogte en stuit telkens terug tot 80\% van de vorige hoogte. Welke afstand heeft de bal afgelegd op het moment dat hij voor de zesde keer de grond raakt?
\begin{solution}
Ongeveer 19,14 m
\end{solution}
\begin{fullsolution}
De bal valt 3 m, en daarna gaat hij 5 keer omhoog en weer omlaag. De stuithoogtes zijn $2{,}4;\ 1{,}92;\ \ldots$ met $q=0{,}8$. Hun som is $2{,}4\cdot\dfrac{0{,}8^5-1}{0{,}8-1}=8{,}06784$. De afstand is $3+2\cdot 8{,}06784\approx 19{,}14$ m.
\end{fullsolution}
\end{oefening}

\begin{oefening}[2]
\thema{meetkundig: aantal termen uit de som}Hoeveel termen van $3,\ 6,\ 12,\ \ldots$ moet je minstens optellen om meer dan een miljoen te krijgen?
\begin{solution}
19
\end{solution}
\begin{fullsolution}
$3\,(2^{\,n}-1)>1\,000\,000$ geeft $2^{\,n}>333\,334{,}33\ldots$ en $n>18{,}35$, dus $n=19$.
\end{fullsolution}
\end{oefening}

\begin{oefening}[2]
\thema{rekenkundig: grootste som}Je telt de termen van $61,\ 57,\ 53,\ \ldots$ één voor één op. Bij hoeveel termen is de som het grootst? Hoe groot is ze dan?
\begin{solution}
Bij 16 termen; 496
\end{solution}
\begin{fullsolution}
De som groeit zolang de termen positief zijn. $u_n=65-4n>0$ geeft $n<16{,}25$, dus $u_{16}=1$ is de laatste positieve term. Dan is $s_{16}=16\cdot\dfrac{61+1}{2}=496$.
\end{fullsolution}
\end{oefening}

\begin{oefening}[2]
\thema{context: meetkundig}Je mag kiezen: \eur 100\,000 ineens, of 30 dagen lang geld volgens dit schema: 1 cent op dag 1, 2 cent op dag 2, 4 cent op dag 3, enzovoort. Wat kies je?
\begin{solution}
Het schema: \eur 10\,737\,418,23
\end{solution}
\begin{fullsolution}
$s_{30}=1\cdot(2^{30}-1)=1\,073\,741\,823$ cent, dat is \eur 10\,737\,418,23. Dat is ruim 100 keer meer. Na 20 dagen heb je nochtans maar $2^{20}-1$ cent, ongeveer \eur 10\,486.
\end{fullsolution}
\end{oefening}

\begin{oefening}[3]
\thema{rekenkundig: bewijs}Bewijs: als $s_n = an^2 + bn$ voor alle $n$, met $a$ en $b$ vaste getallen, dan is de rij rekenkundig. Wat is het verschil?
\begin{solution}
Het verschil is $2a$.
\end{solution}
\begin{fullsolution}
$s_{n-1}=a(n-1)^2+b(n-1)=an^2-2an+a+bn-b$, dus $u_n=s_n-s_{n-1}=2an-a+b$ voor $n\geqslant 2$. Ook $u_1=s_1=a+b$ past hierin. De coëfficiënt van $n$ is $2a$, dus $v=2a$.
\end{fullsolution}
\end{oefening}

\begin{oefening}[3]
\begin{metprent}{slak}
\thema{meetkundig: de som loopt niet op}Een slak kruipt op dag 1 een meter. Elke dag kruipt ze de helft van de dag ervoor.
\begin{enumerate}
  \item Toon aan dat de afstand na $n$ dagen $s_n = 2 - 2 \cdot (\tfrac12)^n$ meter is.
  \item Na hoeveel dagen heeft de slak voor het eerst meer dan 1,99 m afgelegd?
  \item Haalt de slak ooit 2 m? Leg uit.
\end{enumerate}
\end{metprent}
\begin{solution}
\textbf{b.}~na 8 dagen \quad \textbf{c.}~nee
\end{solution}
\begin{fullsolution}
\textbf{a.}~$s_n=\dfrac{(\tfrac12)^n-1}{\tfrac12-1}=-2\bigl((\tfrac12)^n-1\bigr)=2-2\cdot(\tfrac12)^n$.\\
\textbf{b.}~$2\cdot(\tfrac12)^n<0{,}01$ geeft $(\tfrac12)^n<0{,}005$ en $n>7{,}64$, dus na 8 dagen.\\
\textbf{c.}~Nee: $2\cdot(\tfrac12)^n$ is altijd positief, dus $s_n<2$. De som komt wel zo dicht bij 2 als je wilt. In hoofdstuk 4 noemen we dat een limiet.
\end{fullsolution}
\end{oefening}

\begin{oefening}[3]
\thema{rekenkundig: slim groeperen}Bereken $100^2 - 99^2 + 98^2 - 97^2 + \cdots + 2^2 - 1^2$ zonder rekentoestel.
\begin{solution}
5050
\end{solution}
\begin{fullsolution}
Gebruik $a^2-b^2=(a-b)(a+b)$: elk paar is $(a+b)$ met $a-b=1$. Dat geeft $199+195+\cdots+3$. Dit is een rekenkundige rij met 50 termen, dus de som is $50\cdot\dfrac{199+3}{2}=5050$.
\end{fullsolution}
\end{oefening}

\begin{oefening}[3]
\thema{lening}Je leent \eur 10\,000 aan 4\% per jaar, en betaalt 5 jaar lang aan het einde van elk jaar een vast bedrag $a$. Na de laatste betaling is de lening volledig terugbetaald.
\begin{enumerate}
  \item Waarom geldt $a \cdot \dfrac{1{,}04^5 - 1}{1{,}04 - 1} = 10\,000 \cdot 1{,}04^5$? Denk aan twee rekeningen die na 5 jaar evenveel waard moeten zijn.
  \item Bereken $a$. Hoeveel intrest betaal je in totaal?
\end{enumerate}
\begin{solution}
\textbf{b.}~$a\approx$ \eur 2246,27; ongeveer \eur 1231,36 intrest
\end{solution}
\begin{fullsolution}
\textbf{a.}~Links: de betalingen, elk aangegroeid tot het einde van jaar 5. Rechts: de schuld, die zonder betalingen tot dan aangroeit. Die twee moeten gelijk zijn.\\
\textbf{b.}~$a=10\,000\cdot 1{,}04^5\cdot\dfrac{0{,}04}{1{,}04^5-1}\approx 2246{,}27$. In totaal betaal je $5a\approx 11\,231{,}36$, dus ongeveer \eur 1231,36 intrest.
\end{fullsolution}
\end{oefening}

\begin{oefening}[3]
\thema{som in een tabel}Je telt alle getallen op in de tafels van vermenigvuldiging van 1 tot en met 10, dus de $10 \times 10$ producten van $1 \cdot 1$ tot en met $10 \cdot 10$.
\begin{enumerate}
  \item Hoeveel is de som van de tafel van $k$?
  \item Toon aan dat het totaal gelijk is aan $(1 + 2 + \cdots + 10)^2$ en bereken het.
\end{enumerate}
\begin{solution}
\textbf{a.}~$55k$ \quad \textbf{b.}~3025
\end{solution}
\begin{fullsolution}
\textbf{a.}~$k\cdot 1+k\cdot 2+\cdots+k\cdot 10=k\,(1+2+\cdots+10)=55k$.\\
\textbf{b.}~Het totaal is $55\cdot 1+55\cdot 2+\cdots+55\cdot 10=55\cdot(1+2+\cdots+10)=55^2=3025$.
\end{fullsolution}
\end{oefening}

\begin{oefening}[3]
\thema{bewijs: rekenkundige som}Bewijs: bij een rekenkundige rij is $s_n = \dfrac{n}{2}\bigl(2u_1 + (n-1)v\bigr)$. Toon daarmee aan dat de som van de eerste $n$ oneven getallen gelijk is aan $n^2$.
\begin{solution}
$u_1=1$, $v=2$ invullen geeft $n^2$.
\end{solution}
\begin{fullsolution}
Vul $u_n=u_1+(n-1)v$ in de somformule in: $s_n=n\cdot\dfrac{u_1+u_1+(n-1)v}{2}=\dfrac{n}{2}\bigl(2u_1+(n-1)v\bigr)$.\\
Voor de oneven getallen is $u_1=1$ en $v=2$, dus $s_n=\dfrac{n}{2}\bigl(2+2(n-1)\bigr)=\dfrac{n}{2}\cdot 2n=n^2$.
\end{fullsolution}
\end{oefening}
\end{oefeningen}

%----------------------------------------------------------------------------------------
% KOPPEL: verbanden met andere leerstof
%----------------------------------------------------------------------------------------
\opdrachtsectie{Koppel}

\lkr{Matrices en complexe getallen zijn gezien in het 5de jaar. Logica, kansrekenen en integralen staan ook in het 6de jaar: kijk na of ze al aan bod kwamen, of laat die oefeningen weg.}

\begin{oefeningen}
\begin{oefening}[3]\label{oef:h3-inductie}
\koppel{discreet}{Logica en bewijzen} Voor kwadraten bestaat geen rekenkundige of meetkundige som, maar wel een formule: $\displaystyle\sum_{i=1}^{n} i^2 = \dfrac{n(n+1)(2n+1)}{6}$.
\begin{enumerate}
  \item Controleer de formule voor $n = 10$.
  \item Bewijs de formule met volledige inductie.
\end{enumerate}
\begin{solution}
\textbf{a.}~385
\end{solution}
\begin{fullsolution}
\textbf{a.}~$\dfrac{10\cdot 11\cdot 21}{6}=385$.\\
\textbf{b.}~Basis: voor $n=1$ staat er $1=\dfrac{1\cdot 2\cdot 3}{6}$.\\
Stap: stel dat de formule geldt voor $n=k$. Dan is de som tot $k+1$ gelijk aan\\
$\dfrac{k(k+1)(2k+1)}{6}+(k+1)^2=\dfrac{(k+1)\bigl(2k^2+k+6k+6\bigr)}{6}=\dfrac{(k+1)(k+2)(2k+3)}{6}$.\\
Dat is de formule voor $n=k+1$.
\end{fullsolution}
\end{oefening}

\begin{oefening}[2]
\koppel{data}{Kansrekenen} Je gooit met een dobbelsteen tot je een zes gooit. $u_i$ is de kans dat de eerste zes valt bij worp $i$.
\begin{enumerate}
  \item Bereken met de somformule de kans dat de eerste zes valt in de eerste $n$ worpen.
  \item Bereken diezelfde kans via de complementaire gebeurtenis. Komt hetzelfde uit?
  \item Hoeveel worpen heb je nodig om minstens 90\% kans op een zes te hebben?
\end{enumerate}
\begin{solution}
\textbf{a.}~$1-(\tfrac56)^n$ \quad \textbf{b.}~ja \quad \textbf{c.}~13
\end{solution}
\begin{fullsolution}
\textbf{a.}~$u_i=\tfrac16\cdot(\tfrac56)^{i-1}$, dus $s_n=\tfrac16\cdot\dfrac{(\tfrac56)^n-1}{\tfrac56-1}=1-(\tfrac56)^n$.\\
\textbf{b.}~De kans op $n$ keer geen zes is $(\tfrac56)^n$, dus de gevraagde kans is $1-(\tfrac56)^n$.\\
\textbf{c.}~$(\tfrac56)^n\leqslant 0{,}1$ geeft $n\geqslant 12{,}63$, dus 13 worpen.
\end{fullsolution}
\end{oefening}

\begin{oefening}[2]
\koppel{algebra}{Complexe getallen} De somformule geldt ook als $q$ een complex getal is.
\begin{enumerate}
  \item Bereken $1 + \mathrm{i} + \mathrm{i}^2 + \cdots + \mathrm{i}^{10}$ met de somformule, en controleer door de termen uit te schrijven.
  \item Stel dat $w \neq 1$ en $w^n = 1$. Toon aan dat $1 + w + w^2 + \cdots + w^{n-1} = 0$. Wat betekent dat voor de $n$-de wortels uit 1?
\end{enumerate}
\begin{solution}
\textbf{a.}~$\mathrm{i}$ \quad \textbf{b.}~Hun som is 0.
\end{solution}
\begin{fullsolution}
\textbf{a.}~$\dfrac{\mathrm{i}^{11}-1}{\mathrm{i}-1}=\dfrac{-\mathrm{i}-1}{\mathrm{i}-1}=\dfrac{(-\mathrm{i}-1)(-\mathrm{i}-1)}{(\mathrm{i}-1)(-\mathrm{i}-1)}=\dfrac{2\mathrm{i}}{2}=\mathrm{i}$. Controle: per vier termen is de som 0, en $1+\mathrm{i}-1=\mathrm{i}$ blijft over.\\
\textbf{b.}~$\dfrac{w^n-1}{w-1}=\dfrac{0}{w-1}=0$. Neem voor $w$ de wortel met het kleinste positieve argument: dan zijn $1,\ w,\ \ldots,\ w^{n-1}$ precies alle $n$-de wortels uit 1. Hun som is dus 0 (voor $n\geqslant 2$).
\end{fullsolution}
\end{oefening}

\begin{oefening}[2]
\koppel{algebra}{Matrices} Gegeven: $A = \begin{pmatrix} 2 & 0 \\ 0 & 3 \end{pmatrix}$ en $I = \begin{pmatrix} 1 & 0 \\ 0 & 1 \end{pmatrix}$.
\begin{enumerate}
  \item Bereken $S = I + A + A^2 + A^3$. Welke twee sommen herken je?
  \item Toon aan: $(A - I)\,S = A^4 - I$. Vergelijk met het bewijs van de somformule.
\end{enumerate}
\begin{solution}
\textbf{a.}~$S=\begin{pmatrix} 15 & 0 \\ 0 & 40 \end{pmatrix}$
\end{solution}
\begin{fullsolution}
\textbf{a.}~$S=\begin{pmatrix} 15 & 0 \\ 0 & 40 \end{pmatrix}$: op de diagonaal staan $1+2+4+8$ en $1+3+9+27$.\\
\textbf{b.}~$(A-I)S=AS-S=A+A^2+A^3+A^4-I-A-A^2-A^3=A^4-I$. Dit is dezelfde redenering als $q\cdot s_n-s_n$. Controle: $\begin{pmatrix} 1 & 0 \\ 0 & 2 \end{pmatrix}\begin{pmatrix} 15 & 0 \\ 0 & 40 \end{pmatrix}=\begin{pmatrix} 15 & 0 \\ 0 & 80 \end{pmatrix}$.
\end{fullsolution}
\end{oefening}

\begin{oefening}[3]
\koppel{analyse}{Integralen} Verdeel $[0, 1]$ in $n$ gelijke stukken en teken onder de grafiek van $f(x) = x^2$ de rechthoeken met breedte $\tfrac1n$ en als hoogte de functiewaarde rechts in elk stuk.

\begin{tikzpicture}[x=2.6cm,y=1.7cm,font=\scriptsize]
  \foreach \i in {1,2,3,4,5}{
    \draw[accent,fill=accent!15] ({(\i-1)/5},0) rectangle ({\i/5},{(\i/5)^2});}
  \draw[->,lijngrijs] (-.05,0) -- (1.15,0) node[right,text=black] {$x$};
  \draw[->,lijngrijs] (0,-.05) -- (0,1.15) node[above,text=black] {$y$};
  \draw[thick,domain=0:1.04,samples=30] plot (\x,{\x*\x});
  \node[below=1pt,text=grijs] at (1,0) {1};
\end{tikzpicture}
\begin{enumerate}
  \item Toon aan dat de totale oppervlakte van de rechthoeken gelijk is aan $\dfrac{1}{n^3} \displaystyle\sum_{i=1}^{n} i^2$.
  \item Gebruik de formule uit oefening \ref{oef:h3-inductie} en bereken die oppervlakte voor $n = 10$ en $n = 100$. Welk getal benader je?
\end{enumerate}
\begin{solution}
\textbf{b.}~0,385 en 0,33835; de benadering gaat naar $\tfrac13$
\end{solution}
\begin{fullsolution}
\textbf{a.}~De $i$-de rechthoek heeft breedte $\tfrac1n$ en hoogte $(\tfrac{i}{n})^2$, dus oppervlakte $\dfrac{i^2}{n^3}$.\\
\textbf{b.}~$n=10$: $\dfrac{385}{1000}=0{,}385$. $n=100$: $\dfrac{338\,350}{1\,000\,000}=0{,}33835$. Het getal is $\displaystyle\int_0^1 x^2\,\mathrm{d}x=\tfrac13$.
\end{fullsolution}
\end{oefening}
\end{oefeningen}

%----------------------------------------------------------------------------------------
\opdrachtsectie{Test jezelf}

\begin{oefeningen}
\begin{oefening}[1]\label{oef:h3-t1}
Een rekenkundige rij heeft $u_1 = 6$ en $v = 5$. Bereken $s_{12}$.
\begin{solution}
402
\end{solution}
\begin{fullsolution}
$u_{12}=6+55=61$, dus $s_{12}=12\cdot\dfrac{6+61}{2}=402$.
\end{fullsolution}
\end{oefening}

\begin{oefening}[1]\label{oef:h3-t2}
Bereken $13 + 17 + 21 + \cdots + 93$.
\begin{solution}
1113
\end{solution}
\begin{fullsolution}
$13+(n-1)\cdot 4=93$ geeft $n=21$, dus de som is $21\cdot\dfrac{13+93}{2}=1113$.
\end{fullsolution}
\end{oefening}

\begin{oefening}[1]\label{oef:h3-t3}
Bereken $s_7$ van de rij $4,\ 12,\ 36,\ \ldots$
\begin{solution}
4372
\end{solution}
\begin{fullsolution}
$q=3$, dus $s_7=4\cdot\dfrac{3^7-1}{3-1}=4372$.
\end{fullsolution}
\end{oefening}

\begin{oefening}[2]\label{oef:h3-t4}
Bereken $\displaystyle\sum_{k=4}^{13} (3k - 2)$.
\begin{solution}
235
\end{solution}
\begin{fullsolution}
$10+13+\cdots+37$: $13-4+1=10$ termen, dus $10\cdot\dfrac{10+37}{2}=235$.
\end{fullsolution}
\end{oefening}

\begin{oefening}[1]\label{oef:h3-t5}
Noteer met het somteken: $7 + 14 + 21 + \cdots + 84$.
\begin{solution}
$\displaystyle\sum_{i=1}^{12} 7i$
\end{solution}
\begin{fullsolution}
De veelvouden van 7, en $84=7\cdot 12$: $\displaystyle\sum_{i=1}^{12} 7i$.
\end{fullsolution}
\end{oefening}

\begin{oefening}[2]\label{oef:h3-t6}
Hoeveel termen van $5,\ 9,\ 13,\ \ldots$ moet je optellen om 230 te krijgen?
\begin{solution}
10
\end{solution}
\begin{fullsolution}
$u_n=4n+1$, dus $s_n=n\cdot\dfrac{5+4n+1}{2}=n(2n+3)$. Uit $2n^2+3n-230=0$ volgt $n=10$ of $n=-\tfrac{23}{2}$. Enkel $n=10$ kan.
\end{fullsolution}
\end{oefening}

\begin{oefening}[2]\label{oef:h3-t7}
Gegeven: $s_n = 3n^2 - n$. Bepaal $u_n$.
\begin{solution}
$u_n=6n-4$
\end{solution}
\begin{fullsolution}
$s_{n-1}=3(n-1)^2-(n-1)=3n^2-7n+4$, dus $u_n=6n-4$ voor $n\geqslant 2$. Ook $u_1=s_1=2$ past hierin.
\end{fullsolution}
\end{oefening}

\begin{oefening}[2]\label{oef:h3-t8}
Je loopt op dag 1 een afstand van 2 km, en elke dag 10\% meer dan de dag ervoor. Hoeveel kilometer loop je in 7 dagen?
\begin{solution}
Ongeveer 18,97 km
\end{solution}
\begin{fullsolution}
$s_7=2\cdot\dfrac{1{,}1^7-1}{0{,}1}\approx 18{,}97$.
\end{fullsolution}
\end{oefening}

\begin{oefening}[2]\label{oef:h3-t9}
Je stort 5 jaar lang aan het einde van elk jaar \eur 800 op een rekening met 3\% intrest. Hoeveel staat er na de laatste storting op?
\begin{solution}
Ongeveer \eur 4247,31
\end{solution}
\begin{fullsolution}
$E=800\cdot\dfrac{1{,}03^5-1}{0{,}03}\approx 4247{,}31$.
\end{fullsolution}
\end{oefening}
\end{oefeningen}

\herhaalschema{%
  \fv{oef:h3-t1}{vb:h3-rr}{oef:h3-rrsom}%
  \fv{oef:h3-t2}{vb:h3-aantal}{oef:h3-aantal}%
  \fv{oef:h3-t3}{vb:h3-mr}{oef:h3-mrsom}%
  \fv{oef:h3-t4}{vb:h3-somteken}{oef:h3-somindex}%
  \fv{oef:h3-t5}{vb:h3-somteken}{oef:h3-noteer}%
  \fv{oef:h3-t6}{vb:h3-nuit}{oef:h3-nuitsom}%
  \fv{oef:h3-t7}{vb:h3-uitsom}{oef:h3-uitsom}%
  \fv{oef:h3-t8}{vb:h3-mr}{oef:h3-mrcontext}%
  \fv{oef:h3-t9}{vb:h3-sparen}{oef:h3-sparen}}

\ifzelfstudie
\section*{Maak je eigen samenvatting}
Sluit dit boek en schrijf in je schrift: de somformule van een rekenkundige en van een meetkundige rij, met het bewijs van elk, hoe je het aantal termen vindt, wat het somteken betekent, hoe je $u_n$ uit $s_n$ haalt, en welke fouten je wilt vermijden. Vergelijk daarna met de kaders.
\fi

%----------------------------------------------------------------------------------------
% LIMIET VAN EEN RIJ. Alle opgaven zijn eigen werk.
%
% Opbouw: Start, leerstof (kern, uitgewerkt voorbeeld, korte oefening), en daarna op een
% nieuwe pagina de Opdrachten (door elkaar, stijgende moeilijkheid), Koppel en Test jezelf.
% Labels van dit hoofdstuk beginnen met h4-. \minigrafiek staat in hoofdstuk 1.
%----------------------------------------------------------------------------------------

\mychapter{4}{Limiet van een rij}

\lkr{Leerplan B+S'': LPD 45 (formele definitie, met kwantoren, toegepast op een concreet voorbeeld) en LPD 46 (limieten bepalen, convergentie, rekenregels, insluitstelling). De meetkundige reeks is een suggestie bij LPD 44 en staat hier omdat ze de limiet nodig heeft. Voorkennis: limieten van functies op oneindig. Rode draad: je kunt de limiet op drie manieren vinden, uit de grafiek, met rekenregels en uit een vergelijking, maar elke manier heeft haar voorwaarden.}

\begin{oefeningen}[Start]
\begin{oefening}
\blokkop{Opwarmer}
\begin{enumerate}
  \item Bereken $\displaystyle\lim_{x\to+\infty} \frac{1}{x}$, $\displaystyle\lim_{x\to+\infty} \frac{2x+1}{x-3}$ en $\displaystyle\lim_{x\to+\infty} (1 - x^2)$.
  \item Schrijf $|x - 5| < 0{,}1$ als een open interval.
  \item Voor welke $n > 0$ is $\dfrac{1}{n} < 0{,}001$?
  \item Voor welke $n$ is $0{,}5^{\,n} < 0{,}01$?
\end{enumerate}
\begin{solution}
\textbf{a.}~0; 2; $-\infty$ \quad \textbf{b.}~$]4{,}9;\ 5{,}1[$ \quad \textbf{c.}~$n>1000$ \quad \textbf{d.}~$n>6{,}64$
\end{solution}
\begin{fullsolution}
\textbf{a.}~De noemer wordt onbegrensd groot: 0. Hoogste machten: $\tfrac{2x}{x}=2$. De term $-x^2$ beslist: $-\infty$.\\
\textbf{b.}~$-0{,}1<x-5<0{,}1$, dus $4{,}9<x<5{,}1$.\\
\textbf{c.}~Omgekeerde nemen, en dan keert het teken om: $n>1000$.\\
\textbf{d.}~$n>\log_{0{,}5} 0{,}01\approx 6{,}64$. Het teken keert om, want het grondtal is kleiner dan 1.
\end{fullsolution}
\end{oefening}

\begin{oefening}\label{oef:h4-verken}
\blokkop{Verken}\quad Gegeven: $u_n = \dfrac{2n+3}{n}$.
\begin{enumerate}
  \item Bereken $u_{10}$, $u_{100}$ en $u_{1000}$. Welk getal naderen de termen?
  \item Vanaf welk volgnummer ligt $u_n$ minder dan 0,01 van dat getal? En minder dan 0,0001?
  \item Kun je voor elke afstand, hoe klein ook, zo'n volgnummer vinden? Hoe?
\end{enumerate}
\begin{solution}
\textbf{a.}~2,3; 2,03; 2,003; het getal 2 \quad \textbf{b.}~vanaf 301; vanaf 30\,001 \quad \textbf{c.}~ja
\end{solution}
\begin{fullsolution}
\textbf{a.}~$u_n=2+\tfrac3n$, dus 2,3; 2,03 en 2,003. De termen naderen 2.\\
\textbf{b.}~$\tfrac3n<0{,}01$ geeft $n>300$, dus vanaf 301. $\tfrac3n<0{,}0001$ geeft $n>30\,000$, dus vanaf 30\,001.\\
\textbf{c.}~Ja: voor een afstand $\varepsilon>0$ geeft $\tfrac3n<\varepsilon$ dat $n>\tfrac{3}{\varepsilon}$.
\end{fullsolution}
\end{oefening}
\end{oefeningen}

\lkr{Laat oefening \ref{oef:h4-verken} maken vóór de uitleg. Vraag c is de kern van de formele definitie: die komt straks in sectie 2 terug. Stel ook de vraag wat `naderen' eigenlijk betekent.}

%========================================================================================
\section{Convergent en divergent}

\subsection{De begrippen}

\begin{theorie}
Een rij is \emph{convergent} als de termen willekeurig dicht bij een reëel getal $L$ komen en daar blijven. $L$ is de \emph{limiet}: \ $\formule{\displaystyle\lim_{n\to+\infty} u_n = L}$. Een rij die niet convergeert, is \emph{divergent}.

\vspace{1mm}
\noindent\minigrafiek{limiet 2}{-.5}{5.5}{2,4}{1/5,2/3.5,3/3,4/2.75,5/2.6,6/2.5}\hfill
\minigrafiek{limiet $+\infty$}{-.5}{10}{5,10}{1/.5,2/2,3/3,4/5,5/7,6/9.5}\hfill
\minigrafiek{geen limiet}{-3}{3}{-2,2}{1/-2,2/2,3/-2,4/2,5/-2,6/2}
\end{theorie}

De limiet hoeft niet bereikt te worden, maar het mag wel: de rij $5,\ 5,\ 5,\ \ldots$ heeft limiet 5.

\voorbeeld{}{Bepaal de limiet. Is de rij convergent of divergent?\\[.5ex]
\textbf{a.}~$u_n = 4 - \tfrac1n$ \qquad \textbf{b.}~$u_n = 3n - n^2$ \qquad \textbf{c.}~$u_n = 1 + (-1)^n$\label{vb:h4-soort}}{%
\begin{stappen}
\text{a. } \lim_{n\to+\infty} \left(4 - \dfrac1n\right) = 4 & $\tfrac1n$ wordt willekeurig klein: \textbf{convergent}. \\[1.5ex]
\text{b. } \lim_{n\to+\infty} \left(3n - n^2\right) = -\infty & De term $-n^2$ beslist: \textbf{divergent}. \\[1ex]
\text{c. } 0,\ 2,\ 0,\ 2,\ \ldots & De limiet bestaat niet: \textbf{divergent}.
\end{stappen}}

\lkr{Veel leerlingen denken dat een limiet nooit bereikt wordt, en dat divergent hetzelfde is als `naar oneindig gaan'. De constante rij en de rij $1 + (-1)^n$ weerleggen allebei.}

\nujij{Bepaal de limiet van $u_n = \dfrac{(-1)^n}{n}$ en van $u_n = (-1)^n \cdot n$.}{Eerste rij: 0, convergent. Tweede rij: de limiet bestaat niet, divergent.}

\begin{oefeningen}
\begin{oefening}
Bepaal de limiet. Is de rij convergent of divergent?
\begin{reeks}
  \item $u_n = 2 + \tfrac3n$
  \item $u_n = n^2 + 1$
  \item $u_n = 3 \cdot (-1)^n$
\end{reeks}
\begin{solution}
\textbf{a.}~2, convergent \quad \textbf{b.}~$+\infty$, divergent \quad \textbf{c.}~bestaat niet, divergent
\end{solution}
\begin{fullsolution}
\textbf{a.}~$\tfrac3n$ wordt willekeurig klein.\\
\textbf{b.}~$n^2$ groeit onbegrensd.\\
\textbf{c.}~$-3,\ 3,\ {-3},\ \ldots$: de termen naderen geen vast getal.
\end{fullsolution}
\end{oefening}
\end{oefeningen}

\subsection{Rekenkundige en meetkundige rijen}

\begin{theorie}
\begin{tabular}{@{}>{\RaggedRight\arraybackslash}p{6.2cm}>{\RaggedRight\arraybackslash}p{8.4cm}@{}}
\textbf{Rekenkundige rij} & \textbf{Meetkundige rij} ($u_1 \neq 0$) \\[1mm]
$v > 0$: limiet $+\infty$ & $|q| < 1$: limiet 0 \\
$v < 0$: limiet $-\infty$ & $q = 1$: limiet $u_1$ \\
$v = 0$: limiet $u_1$ & $q > 1$: limiet $+\infty$ of $-\infty$, volgens het teken van $u_1$ \\
 & $q \leqslant -1$: de limiet bestaat niet
\end{tabular}
\end{theorie}

\voorbeeld{}{Bepaal de limiet.\\[.5ex]
\textbf{a.}~$9,\ 7,\ 5,\ \ldots$ \qquad \textbf{b.}~$-3,\ {-6},\ {-12},\ \ldots$ \qquad \textbf{c.}~$8,\ {-4},\ 2,\ \ldots$ \qquad \textbf{d.}~$1,\ {-3},\ 9,\ \ldots$\label{vb:h4-rrmr}}{%
\begin{stappen}
\text{a. } v = -2 < 0 & Limiet $-\infty$. \\
\text{b. } q = 2 > 1 \text{ en } u_1 < 0 & Limiet $-\infty$. \\
\text{c. } q = -\tfrac12,\ |q| < 1 & Limiet 0, ook al wisselt het teken. \\
\text{d. } q = -3 \leqslant -1 & De limiet bestaat niet.
\end{stappen}}

\begin{oefeningen}
\begin{oefening}
Bepaal de limiet van \ \textbf{a.}~$2,\ 6,\ 18,\ \ldots$ \quad \textbf{b.}~$100,\ 10,\ 1,\ \ldots$ \quad \textbf{c.}~$5,\ {-5},\ 5,\ \ldots$
\begin{solution}
\textbf{a.}~$+\infty$ \quad \textbf{b.}~0 \quad \textbf{c.}~bestaat niet
\end{solution}
\begin{fullsolution}
\textbf{a.}~$q=3>1$ en $u_1>0$.\\
\textbf{b.}~$q=\tfrac{1}{10}$, dus $|q|<1$.\\
\textbf{c.}~$q=-1$.
\end{fullsolution}
\end{oefening}
\end{oefeningen}

%========================================================================================
\section{De formele definitie}

\begin{theorie}
\begin{minipage}[t]{.57\linewidth}
`Willekeurig dicht' betekent: hoe smal je een band rond $L$ ook kiest, vanaf een bepaald volgnummer liggen alle termen in die band.

\vspace{1.5mm}
$\formule{\displaystyle\lim_{n\to+\infty} u_n = L \ \Leftrightarrow}$

\vspace{.5mm}
$\formule{\forall\, \varepsilon > 0,\ \exists\, N \in \mathbb{N}: \ n > N \Rightarrow |u_n - L| < \varepsilon}$

\vspace{1.5mm}
$\formule{\displaystyle\lim_{n\to+\infty} u_n = +\infty \ \Leftrightarrow}$

\vspace{.5mm}
$\formule{\forall\, r > 0,\ \exists\, N \in \mathbb{N}: \ n > N \Rightarrow u_n > r}$

\vspace{1mm}
Bij $-\infty$ staat er op het einde $u_n < -r$.
\end{minipage}\hfill
\begin{minipage}[t]{.4\linewidth}
\vspace{-2mm}
\begin{tikzpicture}[x=.33cm,y=.72cm,>={Stealth[length=1.6mm]},font=\scriptsize]
  \fill[accent!18] (0,1.4) rectangle (13,2.6);
  \draw[lijngrijs,dashed] (0,2) -- (13,2);
  \draw[->,lijngrijs] (-.1,0) -- (13.5,0) node[right,text=black] {$n$};
  \draw[->,lijngrijs] (0,-.1) -- (0,5.4) node[above,text=black] {$u_n$};
  \draw[lijngrijs,dashed] (5,0) -- (5,2.6);
  \node[below=-1pt] at (5,0) {$N$};
  \node[left=-1pt] at (0,2) {$L$};
  \node[left=-1pt] at (0,2.75) {$L+\varepsilon$};
  \node[left=-1pt] at (0,1.25) {$L-\varepsilon$};
  \foreach \x/\y in {1/5,2/3.5,3/3,4/2.75,5/2.6,6/2.5,7/2.43,8/2.38,9/2.33,10/2.3,11/2.27,12/2.25}
    \fill[accent] (\x,\y) circle (1.5pt);
\end{tikzpicture}
\end{minipage}
\end{theorie}

\thuis{Lees de definitie als een spel. Je tegenstander kiest een $\varepsilon$, hoe klein ook. Jij moet dan een $N$ vinden zodat alle termen na $u_N$ minder dan $\varepsilon$ van $L$ liggen. Lukt dat voor elke $\varepsilon$, dan is $L$ de limiet. De volgorde telt: $N$ mag van $\varepsilon$ afhangen.}

\voorbeeld{}{Bewijs met de definitie: $\displaystyle\lim_{n\to+\infty} \dfrac{2n+3}{n} = 2$.\label{vb:h4-bewijs}}{%
\begin{stappen}
\left|\dfrac{2n+3}{n} - 2\right| = \dfrac{3}{n} & Neem een willekeurige $\varepsilon > 0$ en werk de afstand uit. \\[1.5ex]
\dfrac{3}{n} < \varepsilon \ \Leftrightarrow \ n > \dfrac{3}{\varepsilon} & Los op naar $n$. \\[1.5ex]
N \geqslant \dfrac{3}{\varepsilon},\ N \in \mathbb{N} & Kies zo'n $N$. Dan geldt voor elke $n > N$ dat $\left|u_n - 2\right| < \varepsilon$. \textbf{Bewezen.} \\[1.5ex]
\varepsilon = 0{,}01 \ \Rightarrow \ N = 300 & Controle met oefening \ref{oef:h4-verken}: vanaf $n = 301$.
\end{stappen}}

\lkr{Laat het verschil zien tussen het kladwerk (oplossen naar $n$) en het bewijs zelf (kies $N$, en dan volgt het). Een concreet voorbeeld volstaat volgens het leerplan; de oefeningen op niveau 3 gaan verder.}

\voorbeeld{}{Bewijs met de definitie: $\displaystyle\lim_{n\to+\infty} \sqrt{n} = +\infty$.\label{vb:h4-bewijsinf}}{%
\begin{stappen}
\sqrt{n} > r \ \Leftrightarrow \ n > r^2 & Neem een willekeurige $r > 0$ en los op naar $n$. \\
N \geqslant r^2,\ N \in \mathbb{N} & Kies zo'n $N$. Dan geldt voor elke $n > N$ dat $\sqrt{n} > r$. \textbf{Bewezen.}
\end{stappen}}

\nujij{Bewijs met de definitie: $\displaystyle\lim_{n\to+\infty} \dfrac{5}{n} = 0$.}{$\tfrac5n < \varepsilon$ geeft $n > \tfrac{5}{\varepsilon}$. Kies $N \geqslant \tfrac{5}{\varepsilon}$.}

\begin{oefeningen}
\begin{oefening}
Bewijs met de definitie: $\displaystyle\lim_{n\to+\infty} \dfrac{4n+1}{n} = 4$. Welke $N$ past bij $\varepsilon = 0{,}02$?
\begin{solution}
Kies $N\geqslant\tfrac{1}{\varepsilon}$; bij $\varepsilon=0{,}02$ is $N=50$.
\end{solution}
\begin{fullsolution}
Neem $\varepsilon>0$. $\left|\dfrac{4n+1}{n}-4\right|=\dfrac1n$ en $\dfrac1n<\varepsilon$ geeft $n>\dfrac{1}{\varepsilon}$. Kies $N\geqslant\dfrac{1}{\varepsilon}$: dan geldt voor elke $n>N$ dat $|u_n-4|<\varepsilon$. Bij $\varepsilon=0{,}02$ is $N=50$.
\end{fullsolution}
\end{oefening}
\end{oefeningen}

%========================================================================================
\section{Limieten berekenen}

\subsection{Rekenregels}

\begin{theorie}
Is $u_n = f(n)$ en $\displaystyle\lim_{x\to+\infty} f(x) = L$, dan is ook $\displaystyle\lim_{n\to+\infty} u_n = L$. De rekenregels voor limieten van functies mogen dus: limiet van een som, verschil, product en quotiënt. De rijen $\tfrac1n$ en $\tfrac{1}{n^2}$ hebben limiet 0.
\end{theorie}

Let op: het omgekeerde geldt niet. $u_n = \sin(n\pi)$ is de rij $0,\ 0,\ 0,\ \ldots$ met limiet 0, maar $\displaystyle\lim_{x\to+\infty} \sin(x\pi)$ bestaat niet. En $\infty - \infty$ is niet 0: reken eerst om.

\voorbeeld{}{Bereken \ \textbf{a.}~$\displaystyle\lim_{n\to+\infty} \dfrac{3n^2 - n}{n^2 + 2}$ \quad \textbf{b.}~$\displaystyle\lim_{n\to+\infty} \left(\sqrt{n+1} - \sqrt{n}\right)$\label{vb:h4-reken}}{%
\begin{stappen}
\text{a. } \lim_{n\to+\infty} \dfrac{3 - \tfrac1n}{1 + \tfrac{2}{n^2}} = \dfrac{3 - 0}{1 + 0} = 3 & Deel teller en noemer door de hoogste macht, $n^2$. \\[3ex]
\text{b. } \sqrt{n+1} - \sqrt{n} = \dfrac{1}{\sqrt{n+1} + \sqrt{n}} & Dit is $\infty - \infty$. Vermenigvuldig teller en noemer met de toegevoegde uitdrukking. \\[3ex]
\lim_{n\to+\infty} \dfrac{1}{\sqrt{n+1} + \sqrt{n}} = 0 & De noemer wordt onbegrensd groot.
\end{stappen}}

\nujij{Bereken $\displaystyle\lim_{n\to+\infty} \dfrac{2n - n^2}{4n^2 + 1}$.}{$-\tfrac14$: deel teller en noemer door $n^2$.}

\begin{oefeningen}
\begin{oefening}
Bereken.
\begin{reeks}[2]
  \item $\displaystyle\lim_{n\to+\infty} \dfrac{5n + 2}{3n - 1}$
  \item $\displaystyle\lim_{n\to+\infty} \left(\sqrt{n+4} - \sqrt{n}\right)$
\end{reeks}
\begin{solution}
\textbf{a.}~$\tfrac53$ \quad \textbf{b.}~0
\end{solution}
\begin{fullsolution}
\textbf{a.}~Deel door $n$: $\dfrac{5+\tfrac2n}{3-\tfrac1n}\to\dfrac53$.\\
\textbf{b.}~$\sqrt{n+4}-\sqrt{n}=\dfrac{4}{\sqrt{n+4}+\sqrt{n}}\to 0$.
\end{fullsolution}
\end{oefening}
\end{oefeningen}

\subsection{De insluitstelling}

\begin{theorie}
Als $t_n \leqslant u_n \leqslant v_n$ vanaf een bepaald volgnummer, en $\displaystyle\lim_{n\to+\infty} t_n = \lim_{n\to+\infty} v_n = L$, dan is ook $\displaystyle\lim_{n\to+\infty} u_n = L$.
\end{theorie}

\begin{minipage}[t]{.6\linewidth}
Handig bij een factor die blijft schommelen, zoals $\sin n$, $\cos n$ of $(-1)^n$: die laat zich niet met rekenregels aanpakken, maar is wel begrensd.
\end{minipage}\hfill
\begin{minipage}[t]{.37\linewidth}
\vspace{-1mm}
\begin{tikzpicture}[x=.56cm,y=1.05cm,>={Stealth[length=1.6mm]},font=\scriptsize]
  \draw[->,lijngrijs] (-.1,0) -- (8.8,0) node[right,text=black] {$n$};
  \draw[->,lijngrijs] (0,-1.1) -- (0,1.2) node[above,text=black] {$u_n$};
  \draw[lijngrijs,dashed,domain=.9:8.5,samples=40] plot (\x,{1/\x});
  \draw[lijngrijs,dashed,domain=.9:8.5,samples=40] plot (\x,{-1/\x});
  \node[right=-1pt] at (1.25,.82) {$\tfrac1n$};
  \node[right=-1pt] at (1.25,-.82) {$-\tfrac1n$};
  \foreach \x/\y in {1/.54,2/-.208,3/-.33,4/-.163,5/.057,6/.16,7/.108,8/-.018}
    \fill[accent] (\x,\y) circle (1.6pt);
\end{tikzpicture}
\end{minipage}

\voorbeeld{}{Bereken $\displaystyle\lim_{n\to+\infty} \dfrac{\cos n}{n}$. (Zie de figuur.)\label{vb:h4-insluit}}{%
\begin{stappen}
-1 \leqslant \cos n \leqslant 1 & De cosinus is begrensd. \\[1ex]
-\dfrac1n \leqslant \dfrac{\cos n}{n} \leqslant \dfrac1n & Delen door $n > 0$. \\[1.5ex]
\lim_{n\to+\infty} \dfrac{\cos n}{n} = 0 & Beide grenzen hebben limiet 0: de rij wordt ingesloten.
\end{stappen}}

\begin{oefeningen}
\begin{oefening}
Bereken $\displaystyle\lim_{n\to+\infty} \dfrac{\sin n}{n^2}$.
\begin{solution}
0
\end{solution}
\begin{fullsolution}
$-\dfrac{1}{n^2}\leqslant\dfrac{\sin n}{n^2}\leqslant\dfrac{1}{n^2}$, en beide grenzen hebben limiet 0.
\end{fullsolution}
\end{oefening}
\end{oefeningen}

%========================================================================================
\section{Stijgende en dalende rijen}

\begin{theorie}
Een rij is \emph{stijgend} als $u_{n+1} > u_n$ voor elke $n$, en \emph{dalend} als $u_{n+1} < u_n$ voor elke $n$. Een rij is \emph{naar boven begrensd} als er een $M$ is met $u_n \leqslant M$ voor elke $n$, en \emph{naar onder begrensd} als er een $m$ is met $u_n \geqslant m$ voor elke $n$. Begrensd is allebei.

\vspace{1.5mm}
\textbf{Eigenschap.} Een stijgende, naar boven begrensde rij is convergent. Een dalende, naar onder begrensde rij ook.
\end{theorie}

Eén voorwaarde volstaat niet: $u_n = n$ is stijgend en $u_n = (-1)^n$ is begrensd, maar beide zijn divergent.

\thuis{Stel je de grafiek voor: de punten gaan alleen maar omhoog, maar er is een plafond dat ze niet kunnen doorbreken. Dan moeten ze zich wel ophopen onder een bepaalde hoogte. Die hoogte is de limiet. Het bewijs gebruiken we hier niet.}

\voorbeeld{}{Toon aan dat $u_n = \dfrac{n}{n+2}$ stijgend en naar boven begrensd is. Wat is de limiet?\label{vb:h4-monotoon}}{%
\begin{stappen}
u_n = 1 - \dfrac{2}{n+2} & Herschrijf. Dan is $u_n < 1$ voor elke $n$: bovengrens 1. \\[2ex]
u_{n+1} - u_n = \dfrac{2}{n+2} - \dfrac{2}{n+3} & Bereken het verschil van twee opeenvolgende termen. \\[2ex]
\phantom{u_{n+1} - u_n} = \dfrac{2}{(n+2)(n+3)} > 0 & Positief voor elke $n$: stijgend, dus convergent. \\[2ex]
\lim_{n\to+\infty} u_n = 1 & De eigenschap zegt dat er een limiet is. De rekenregels geven hier de waarde.
\end{stappen}}

\begin{oefeningen}
\begin{oefening}
Toon aan dat $u_n = 4 + \dfrac{3}{n}$ dalend en naar onder begrensd is. Wat is de limiet?
\begin{solution}
Ondergrens 4; limiet 4
\end{solution}
\begin{fullsolution}
$u_n>4$ voor elke $n$, dus 4 is een ondergrens. $u_{n+1}-u_n=\dfrac{3}{n+1}-\dfrac3n=\dfrac{-3}{n(n+1)}<0$, dus de rij is dalend. De limiet is 4.
\end{fullsolution}
\end{oefening}
\end{oefeningen}

%========================================================================================
\section{Discrete dynamische processen}

\begin{theorie}
Een proces met $u_n = a\,u_{n-1} + b$ en $|a| < 1$ gaat naar een \emph{evenwicht}. Dat vind je door $u_n$ en $u_{n-1}$ gelijk te stellen aan de limiet $L$:

\vspace{1mm}
\qquad $\formule{L = a L + b}$
\end{theorie}

\begin{minipage}[t]{.6\linewidth}
Discreet: je kijkt met vaste tussenpozen. Dynamisch: er verandert iets. In het evenwicht blijft de hoeveelheid gelijk, maar er verdwijnt en er komt evenveel bij.
\end{minipage}\hfill
\begin{minipage}[t]{.37\linewidth}
\vspace{-1mm}
\begin{tikzpicture}[x=.39cm,y=.0095cm,>={Stealth[length=1.6mm]},font=\scriptsize]
  \draw[->,lijngrijs] (-.1,0) -- (11,0) node[right,text=black] {$n$};
  \draw[->,lijngrijs] (0,-10) -- (0,450) node[above,text=black] {$u_n$};
  \foreach \x in {2,4,...,10} \draw[lijngrijs] (\x,6) -- (\x,-6) node[below=-1pt,text=grijs] {\x};
  \foreach \y in {100,200,300,400} \draw[lijngrijs] (.07,\y) -- (-.07,\y) node[left=-1pt,text=grijs] {\y};
  \draw[lijngrijs,dashed] (0,250) -- (10.6,250) node[below=-1pt,pos=.9,text=black] {250};
  \foreach \x/\y in {0/400,1/370,2/346,3/326.8,4/311.4,5/299.2,6/289.3,7/281.5,8/275.2,9/270.1,10/266.1}
    \fill[accent] (\x,\y) circle (1.6pt);
\end{tikzpicture}
\end{minipage}

\voorbeeld{}{Het zwembad uit hoofdstuk 1: $u_0 = 400$ en $u_n = 0{,}8\,u_{n-1} + 50$. Bepaal het evenwicht $L$. Toon aan dat de rij $v_n = u_n - L$ meetkundig is, en leid daaruit af dat de limiet echt $L$ is.\label{vb:h4-dyn}}{%
\begin{stappen}
L = 0{,}8\,L + 50 \ \Rightarrow \ L = 250 & Het evenwicht: 250 g chloor. \\
v_n = u_n - 250 = 0{,}8\,u_{n-1} - 200 & Vervang $u_n$ door het voorschrift. \\
\phantom{v_n} = 0{,}8\,(u_{n-1} - 250) = 0{,}8\,v_{n-1} & Meetkundig met $q = 0{,}8$ en $v_0 = 150$. \\
u_n = 250 + 150 \cdot 0{,}8^{\,n} & Controle: $u_1 = 250 + 120 = 370$. \\
\lim_{n\to+\infty} u_n = 250 & Want $|q| < 1$, dus $0{,}8^{\,n} \to 0$.
\end{stappen}}

\lkr{De vergelijking $L = aL + b$ veronderstelt dat de limiet bestaat. De hulprij $v_n$ bewijst dat, en sluit aan bij de hulprij uit hoofdstuk 1. Bij $|a| \geqslant 1$ levert de vergelijking vaak nog een getal op, ook al is er geen limiet: zie de opdrachten.}

\nujij{$u_0 = 90$ en $u_n = 0{,}5\,u_{n-1} + 10$. Bepaal het evenwicht en een expliciet voorschrift.}{$L = 20$ en $u_n = 20 + 70 \cdot 0{,}5^{\,n}$.}

\begin{oefeningen}
\begin{oefening}
Een bedrijf heeft 300 werknemers. Elk jaar vertrekt 10\% en komen er 20 nieuwe bij. Geef een recursief voorschrift. Hoeveel werknemers zijn er op lange termijn?
\begin{solution}
$u_n=0{,}9\,u_{n-1}+20$, $u_0=300$; 200
\end{solution}
\begin{fullsolution}
Na het vertrek blijft 90\% over, dus $u_n=0{,}9\,u_{n-1}+20$ met $u_0=300$. $L=0{,}9\,L+20$ geeft $L=200$.
\end{fullsolution}
\end{oefening}
\end{oefeningen}

%========================================================================================
\section{De meetkundige reeks}

\begin{theorie}
Een \emph{reeks} is een oneindige som. Haar waarde is de limiet van de sommen $s_n$. Bij een meetkundige rij met $|q| < 1$ gaat $q^{\,n}$ naar 0, dus:

\vspace{1mm}
\qquad $u_1 + u_1 q + u_1 q^2 + \cdots = \displaystyle\lim_{n\to+\infty} u_1 \cdot \frac{q^{\,n} - 1}{q - 1} = \formule{\dfrac{u_1}{1 - q}}$

\vspace{1.5mm}
Bij $|q| \geqslant 1$ en $u_1 \neq 0$ heeft de reeks geen eindige waarde: ze \emph{divergeert}.
\end{theorie}

\voorbeeld{}{Bereken $8 + 2 + \tfrac12 + \cdots$ en schrijf $0{,}272727\ldots$ als een breuk.\label{vb:h4-reeks}}{%
\begin{stappen}
q = \dfrac14,\quad \dfrac{8}{1 - \tfrac14} = \dfrac{32}{3} & $|q| < 1$: de formule mag. \\[2ex]
0{,}27 + 0{,}0027 + 0{,}000027 + \cdots & Schrijf het getal als een som: $u_1 = 0{,}27$ en $q = 0{,}01$. \\[1ex]
\dfrac{0{,}27}{1 - 0{,}01} = \dfrac{27}{99} = \dfrac{3}{11} & Controle met het rekentoestel: $3 : 11 = 0{,}2727\ldots$
\end{stappen}}

\nujij{Bereken $9 - 3 + 1 - \cdots$ en schrijf $0{,}1515\ldots$ als een breuk.}{$\tfrac{9}{1 + 1/3} = \tfrac{27}{4}$ en $\tfrac{15}{99} = \tfrac{5}{33}$.}

\lkr{De slak uit hoofdstuk 3 past hier: $1 + \tfrac12 + \tfrac14 + \cdots = 2$. De reeks is een suggestie in het leerplan, geen doel.}

\begin{oefeningen}
\begin{oefening}
Bereken $12 + 4 + \tfrac43 + \cdots$ en schrijf $0{,}888\ldots$ als een breuk.
\begin{solution}
18 en $\tfrac89$
\end{solution}
\begin{fullsolution}
$q=\tfrac13$, dus $\dfrac{12}{1-\tfrac13}=18$. Verder is $0{,}8+0{,}08+\cdots=\dfrac{0{,}8}{1-0{,}1}=\dfrac89$.
\end{fullsolution}
\end{oefening}
\end{oefeningen}

%----------------------------------------------------------------------------------------
% OPDRACHTEN: door elkaar. Niveau 1, dan 2, dan 3.
%----------------------------------------------------------------------------------------
\opdrachtenkop

\lkr{De opdrachten staan bewust door elkaar, met stijgende moeilijkheid. In deze versie staat bij elke opdracht het onderwerp, zodat je kunt kiezen. Vraag bij elke limiet ook: convergent of divergent?}

\begin{oefeningen}
\begin{oefening}[1]\label{oef:h4-soort}
\thema{convergent en divergent}Bepaal de limiet. Is de rij convergent of divergent?
\begin{reeks}
  \item $u_n = 5 + \tfrac2n$
  \item $u_n = n^2 - 3$
  \item $u_n = 4 \cdot (-1)^n$
  \item $u_n = 7$
  \item $u_n = \tfrac{(-1)^n}{n}$
  \item $u_n = 1 - n$
\end{reeks}
\begin{solution}
\textbf{a.}~5, conv. \quad \textbf{b.}~$+\infty$, div. \quad \textbf{c.}~bestaat niet, div. \quad \textbf{d.}~7, conv. \quad \textbf{e.}~0, conv. \quad \textbf{f.}~$-\infty$, div.
\end{solution}
\begin{fullsolution}
\textbf{a.}~$\tfrac2n\to 0$.\\
\textbf{b.}~$n^2$ groeit onbegrensd.\\
\textbf{c.}~$-4,\ 4,\ {-4},\ \ldots$: geen vast getal.\\
\textbf{d.}~Een constante rij heeft die constante als limiet.\\
\textbf{e.}~Het teken wisselt, maar de termen komen willekeurig dicht bij 0.\\
\textbf{f.}~De termen dalen onbegrensd.
\end{fullsolution}
\end{oefening}

\begin{oefening}[1]\label{oef:h4-reeks}
\thema{meetkundige reeks}Bereken, of zeg dat de reeks divergeert.
\begin{reeks}[2]
  \item $6 + 3 + \tfrac32 + \cdots$
  \item $1 - \tfrac12 + \tfrac14 - \cdots$
  \item $1 + 3 + 9 + \cdots$
  \item $0{,}444\ldots$
\end{reeks}
\begin{solution}
\textbf{a.}~12 \quad \textbf{b.}~$\tfrac23$ \quad \textbf{c.}~divergeert \quad \textbf{d.}~$\tfrac49$
\end{solution}
\begin{fullsolution}
\textbf{a.}~$q=\tfrac12$: $\dfrac{6}{1-\tfrac12}=12$.\\
\textbf{b.}~$q=-\tfrac12$: $\dfrac{1}{1+\tfrac12}=\dfrac23$.\\
\textbf{c.}~$q=3$: de formule mag niet.\\
\textbf{d.}~$u_1=0{,}4$ en $q=0{,}1$: $\dfrac{0{,}4}{0{,}9}=\dfrac49$.
\end{fullsolution}
\end{oefening}

\begin{oefening}[1]\label{oef:h4-ratio}
\thema{rekenregels}Bereken.
\begin{reeks}[2]
  \item $\displaystyle\lim_{n\to+\infty} \dfrac{5n - 1}{2n + 3}$
  \item $\displaystyle\lim_{n\to+\infty} \dfrac{n^2 + 1}{n^3 - 2}$
  \item $\displaystyle\lim_{n\to+\infty} \dfrac{4 - n^2}{n^2 + n}$
  \item $\displaystyle\lim_{n\to+\infty} \dfrac{n^2 + 3}{n + 1}$
\end{reeks}
\begin{solution}
\textbf{a.}~$\tfrac52$ \quad \textbf{b.}~0 \quad \textbf{c.}~$-1$ \quad \textbf{d.}~$+\infty$
\end{solution}
\begin{fullsolution}
Vergelijk de hoogste machten in teller en noemer.\\
\textbf{a.}~$\dfrac{5n}{2n}=\dfrac52$.\\
\textbf{b.}~De noemer heeft de hoogste macht: 0.\\
\textbf{c.}~$\dfrac{-n^2}{n^2}=-1$.\\
\textbf{d.}~De teller heeft de hoogste macht, en beide zijn positief: $+\infty$.
\end{fullsolution}
\end{oefening}

\begin{oefening}[1]\label{oef:h4-rrmr}
\thema{rekenkundige en meetkundige rijen}Bepaal de limiet.
\begin{reeks}
  \item $2,\ 5,\ 8,\ \ldots$
  \item $100,\ 50,\ 25,\ \ldots$
  \item $3,\ {-6},\ 12,\ \ldots$
  \item $-1,\ {-3},\ {-9},\ \ldots$
  \item $9,\ 7,\ 5,\ \ldots$
  \item $8,\ {-4},\ 2,\ \ldots$
\end{reeks}
\begin{solution}
\textbf{a.}~$+\infty$ \quad \textbf{b.}~0 \quad \textbf{c.}~bestaat niet \quad \textbf{d.}~$-\infty$ \quad \textbf{e.}~$-\infty$ \quad \textbf{f.}~0
\end{solution}
\begin{fullsolution}
\textbf{a.}~$v=3>0$.\\
\textbf{b.}~$q=\tfrac12$.\\
\textbf{c.}~$q=-2$.\\
\textbf{d.}~$q=3$ en $u_1<0$.\\
\textbf{e.}~$v=-2<0$.\\
\textbf{f.}~$q=-\tfrac12$.
\end{fullsolution}
\end{oefening}

\begin{oefening}[1]
\thema{fouten: limiet bereiken}Iemand schrijft: `Een limiet wordt nooit bereikt, dus de rij $5,\ 5,\ 5,\ \ldots$ heeft geen limiet.' Wat is er fout?
\begin{solution}
De limiet is 5: de limiet mag bereikt worden.
\end{solution}
\begin{fullsolution}
Voor elke $\varepsilon>0$ is $|u_n-5|=0<\varepsilon$, dus de limiet is 5. De definitie vraagt niet dat de termen van $L$ verschillen.
\end{fullsolution}
\end{oefening}

\begin{oefening}[1]\label{oef:h4-monotoon}
\thema{stijgend, dalend, begrensd}Is de rij stijgend, dalend of geen van beide? Is ze naar boven begrensd? Naar onder? Is ze convergent?
\begin{reeks}[4]
  \item $u_n = \tfrac1n$
  \item $u_n = n^2$
  \item $u_n = (-1)^n$
  \item $u_n = 5 - \tfrac1n$
\end{reeks}
\begin{solution}
\textbf{a.}~dalend; boven 1, onder 0; conv. \quad \textbf{b.}~stijgend; enkel onder (1); div. \quad \textbf{c.}~geen van beide; boven 1, onder $-1$; div. \quad \textbf{d.}~stijgend; boven 5, onder 4; conv.
\end{solution}
\begin{fullsolution}
\textbf{a.}~Dalend, met $0<u_n\leqslant 1$. Limiet 0.\\
\textbf{b.}~Stijgend, met $u_n\geqslant 1$, en geen bovengrens. Limiet $+\infty$.\\
\textbf{c.}~Wisselt, met $-1\leqslant u_n\leqslant 1$. Begrensd, maar de limiet bestaat niet.\\
\textbf{d.}~Stijgend, met $4\leqslant u_n<5$. Limiet 5.
\end{fullsolution}
\end{oefening}

\begin{oefening}[1]\label{oef:h4-evenwicht}
\thema{dynamisch proces}Bepaal het evenwicht.
\begin{reeks}
  \item $u_n = 0{,}5\,u_{n-1} + 6$
  \item $u_n = 0{,}9\,u_{n-1} + 20$
  \item $u_n = 0{,}75\,u_{n-1} - 3$
\end{reeks}
\begin{solution}
\textbf{a.}~12 \quad \textbf{b.}~200 \quad \textbf{c.}~$-12$
\end{solution}
\begin{fullsolution}
\textbf{a.}~$0{,}5\,L=6$.\\
\textbf{b.}~$0{,}1\,L=20$.\\
\textbf{c.}~$0{,}25\,L=-3$.
\end{fullsolution}
\end{oefening}

\begin{oefening}[1]\label{oef:h4-epsgetal}
\thema{definitie: met getallen}Vanaf welk volgnummer geldt dit?
\begin{enumerate}
  \item $u_n = \dfrac{1}{n^2}$ ligt minder dan 0,0001 van 0.
  \item $u_n = \dfrac{n+1}{n}$ ligt minder dan 0,001 van 1.
  \item $u_n = \dfrac{(-1)^n}{n}$ ligt minder dan 0,05 van 0.
\end{enumerate}
\begin{solution}
\textbf{a.}~101 \quad \textbf{b.}~1001 \quad \textbf{c.}~21
\end{solution}
\begin{fullsolution}
\textbf{a.}~$\dfrac{1}{n^2}<0{,}0001$ geeft $n^2>10\,000$, dus $n>100$.\\
\textbf{b.}~$\left|\dfrac{n+1}{n}-1\right|=\dfrac1n<0{,}001$ geeft $n>1000$.\\
\textbf{c.}~$\left|\dfrac{(-1)^n}{n}\right|=\dfrac1n<0{,}05$ geeft $n>20$.
\end{fullsolution}
\end{oefening}

\begin{oefening}[2]\label{oef:h4-bewijs}
\thema{definitie: bewijs}Bewijs met de definitie.
\begin{reeks}
  \item $\displaystyle\lim_{n\to+\infty} \dfrac{7}{n} = 0$
  \item $\displaystyle\lim_{n\to+\infty} \dfrac{4n+1}{2n} = 2$
  \item $\displaystyle\lim_{n\to+\infty} \dfrac{n-1}{n+3} = 1$
\end{reeks}
\begin{solution}
\textbf{a.}~$N\geqslant\tfrac{7}{\varepsilon}$ \quad \textbf{b.}~$N\geqslant\tfrac{1}{2\varepsilon}$ \quad \textbf{c.}~$N\geqslant\tfrac{4}{\varepsilon}$
\end{solution}
\begin{fullsolution}
Neem telkens een willekeurige $\varepsilon>0$.\\
\textbf{a.}~$\tfrac7n<\varepsilon$ geeft $n>\tfrac{7}{\varepsilon}$. Kies $N\geqslant\tfrac{7}{\varepsilon}$.\\
\textbf{b.}~$\left|\dfrac{4n+1}{2n}-2\right|=\dfrac{1}{2n}<\varepsilon$ geeft $n>\dfrac{1}{2\varepsilon}$. Kies $N\geqslant\dfrac{1}{2\varepsilon}$.\\
\textbf{c.}~$\left|\dfrac{n-1}{n+3}-1\right|=\dfrac{4}{n+3}<\varepsilon$ geeft $n>\dfrac{4}{\varepsilon}-3$. Elke $N\geqslant\dfrac{4}{\varepsilon}$ is goed, want dan is ook $n>\dfrac{4}{\varepsilon}-3$.
\end{fullsolution}
\end{oefening}

\begin{oefening}[2]
\thema{context: dynamisch proces}Een sportclub telt 800 leden. Elk jaar stopt 25\% en sluiten er 150 nieuwe leden aan.
\begin{enumerate}
  \item Geef een recursief voorschrift. Hoeveel leden zijn er op lange termijn?
  \item Toon aan dat $u_n = 600 + 200 \cdot 0{,}75^{\,n}$.
  \item Na hoeveel jaar zijn er voor het eerst minder dan 650 leden?
  \item Hoeveel nieuwe leden per jaar zijn nodig om op lange termijn 1000 leden te hebben?
\end{enumerate}
\begin{solution}
\textbf{a.}~$u_n=0{,}75\,u_{n-1}+150$, $u_0=800$; 600 \quad \textbf{c.}~na 5 jaar \quad \textbf{d.}~250
\end{solution}
\begin{fullsolution}
\textbf{a.}~$L=0{,}75\,L+150$ geeft $L=600$.\\
\textbf{b.}~$v_n=u_n-600=0{,}75\,u_{n-1}-450=0{,}75\,v_{n-1}$ met $v_0=200$, dus $v_n=200\cdot 0{,}75^{\,n}$.\\
\textbf{c.}~$200\cdot 0{,}75^{\,n}<50$ geeft $0{,}75^{\,n}<0{,}25$ en $n>4{,}82$, dus na 5 jaar.\\
\textbf{d.}~$1000=750+b$ geeft $b=250$.
\end{fullsolution}
\end{oefening}

\begin{oefening}[2]
\thema{definitie: oneindig}Bewijs met de definitie.
\begin{reeks}[2]
  \item $\displaystyle\lim_{n\to+\infty} (2n + 1) = +\infty$
  \item $\displaystyle\lim_{n\to+\infty} (5 - n^2) = -\infty$
\end{reeks}
\begin{solution}
\textbf{a.}~$N\geqslant\tfrac{r-1}{2}$ \quad \textbf{b.}~$N\geqslant\sqrt{r+5}$
\end{solution}
\begin{fullsolution}
\textbf{a.}~Neem $r>0$. $2n+1>r$ geeft $n>\dfrac{r-1}{2}$. Kies $N\in\mathbb{N}$ met $N\geqslant\dfrac{r-1}{2}$.\\
\textbf{b.}~Neem $r>0$. $5-n^2<-r$ geeft $n^2>r+5$, dus $n>\sqrt{r+5}$. Kies $N\in\mathbb{N}$ met $N\geqslant\sqrt{r+5}$.
\end{fullsolution}
\end{oefening}

\begin{oefening}[2]\label{oef:h4-insluit}
\thema{insluitstelling}Bereken.
\begin{reeks}
  \item $\displaystyle\lim_{n\to+\infty} \dfrac{(-1)^n}{\sqrt{n}}$
  \item $\displaystyle\lim_{n\to+\infty} \dfrac{3 + \cos n}{n}$
  \item $\displaystyle\lim_{n\to+\infty} \left(2 + \dfrac{\sin n}{n}\right)$
\end{reeks}
\begin{solution}
\textbf{a.}~0 \quad \textbf{b.}~0 \quad \textbf{c.}~2
\end{solution}
\begin{fullsolution}
\textbf{a.}~$-\dfrac{1}{\sqrt n}\leqslant u_n\leqslant\dfrac{1}{\sqrt n}$, en beide grenzen gaan naar 0.\\
\textbf{b.}~$2\leqslant 3+\cos n\leqslant 4$, dus $\dfrac2n\leqslant u_n\leqslant\dfrac4n$.\\
\textbf{c.}~$2-\dfrac1n\leqslant u_n\leqslant 2+\dfrac1n$, en beide grenzen gaan naar 2.
\end{fullsolution}
\end{oefening}

\begin{oefening}[2]
\thema{context: reeks}Een bal valt van 3 m hoogte en stuit telkens terug tot 80\% van de vorige hoogte. Welke afstand legt de bal in totaal af, als hij blijft stuiten tot hij stil ligt?
\begin{solution}
27 m
\end{solution}
\begin{fullsolution}
De stuithoogtes $2{,}4;\ 1{,}92;\ \ldots$ tel je dubbel, omhoog en omlaag. De som is $\dfrac{2{,}4}{1-0{,}8}=12$, dus de afstand is $3+2\cdot 12=27$ m.
\end{fullsolution}
\end{oefening}

\begin{oefening}[2]
\thema{rekenregels: machten}Bereken.
\begin{reeks}
  \item $\displaystyle\lim_{n\to+\infty} \dfrac{3^n + 1}{3^n}$
  \item $\displaystyle\lim_{n\to+\infty} \dfrac{2^n - 1}{3^n}$
  \item $\displaystyle\lim_{n\to+\infty} \dfrac{2^{n+1} + 3^n}{3^{n+1}}$
\end{reeks}
\begin{solution}
\textbf{a.}~1 \quad \textbf{b.}~0 \quad \textbf{c.}~$\tfrac13$
\end{solution}
\begin{fullsolution}
\textbf{a.}~$1+(\tfrac13)^n\to 1$.\\
\textbf{b.}~$(\tfrac23)^n-(\tfrac13)^n\to 0$.\\
\textbf{c.}~$\tfrac23\cdot(\tfrac23)^n+\tfrac13\to\tfrac13$.
\end{fullsolution}
\end{oefening}

\begin{oefening}[2]
\thema{functie en rij}Gegeven: $u_n = \sin(2n\pi)$ en $f(x) = \sin(2x\pi)$.
\begin{enumerate}
  \item Bereken $u_1$, $u_2$ en $u_3$. Wat is de limiet van de rij?
  \item Bestaat $\displaystyle\lim_{x\to+\infty} f(x)$? Wat leer je hieruit?
\end{enumerate}
\begin{solution}
\textbf{a.}~0, 0, 0; limiet 0 \quad \textbf{b.}~Nee.
\end{solution}
\begin{fullsolution}
\textbf{a.}~$\sin(2n\pi)=0$ voor elk natuurlijk getal $n$, dus de limiet is 0.\\
\textbf{b.}~$f$ blijft schommelen tussen $-1$ en 1, dus die limiet bestaat niet. Uit de limiet van de rij volgt dus niets over de limiet van de functie: alleen het omgekeerde geldt.
\end{fullsolution}
\end{oefening}

\begin{oefening}[2]
\thema{dynamisch proces: niet stijgend of dalend}Gegeven: $u_0 = 10$ en $u_n = -0{,}5\,u_{n-1} + 6$.
\begin{enumerate}
  \item Bereken $u_1$, $u_2$ en $u_3$. Is de rij stijgend of dalend?
  \item Bepaal het evenwicht $L$ en een expliciet voorschrift.
  \item De rij is niet stijgend of dalend en toch convergent. Is dat in strijd met de eigenschap uit sectie 4?
\end{enumerate}
\begin{solution}
\textbf{a.}~1; 5,5; 3,25; geen van beide \quad \textbf{b.}~$L=4$, $u_n=4+6\cdot(-0{,}5)^n$ \quad \textbf{c.}~nee
\end{solution}
\begin{fullsolution}
\textbf{a.}~$-5+6=1$, $-0{,}5+6=5{,}5$, $-2{,}75+6=3{,}25$. De termen gaan op en neer.\\
\textbf{b.}~$L=-0{,}5\,L+6$ geeft $L=4$. Met $v_n=u_n-4$ is $v_n=-0{,}5\,v_{n-1}$ en $v_0=6$.\\
\textbf{c.}~Nee: de eigenschap zegt dat stijgend en begrensd voldoende is om te convergeren, niet dat het nodig is.
\end{fullsolution}
\end{oefening}

\begin{oefening}[2]
\thema{reeks: repeterende decimalen}Schrijf als een breuk.
\begin{reeks}
  \item $0{,}121212\ldots$
  \item $1{,}2555\ldots$
  \item $0{,}999\ldots$
\end{reeks}
\begin{solution}
\textbf{a.}~$\tfrac{4}{33}$ \quad \textbf{b.}~$\tfrac{113}{90}$ \quad \textbf{c.}~1
\end{solution}
\begin{fullsolution}
\textbf{a.}~$\dfrac{0{,}12}{1-0{,}01}=\dfrac{12}{99}=\dfrac{4}{33}$.\\
\textbf{b.}~$1{,}2+\dfrac{0{,}05}{1-0{,}1}=\dfrac65+\dfrac{1}{18}=\dfrac{113}{90}$.\\
\textbf{c.}~$\dfrac{0{,}9}{1-0{,}1}=1$. Dus $0{,}999\ldots$ en 1 zijn hetzelfde getal.
\end{fullsolution}
\end{oefening}

\begin{oefening}[2]
\thema{rekenregels: toegevoegde}Bereken.
\begin{reeks}[2]
  \item $\displaystyle\lim_{n\to+\infty} \left(\sqrt{n^2 + n} - n\right)$
  \item $\displaystyle\lim_{n\to+\infty} \left(n^2 - n\right)$
\end{reeks}
\begin{solution}
\textbf{a.}~$\tfrac12$ \quad \textbf{b.}~$+\infty$
\end{solution}
\begin{fullsolution}
\textbf{a.}~$\sqrt{n^2+n}-n=\dfrac{n}{\sqrt{n^2+n}+n}=\dfrac{1}{\sqrt{1+\tfrac1n}+1}\to\dfrac12$.\\
\textbf{b.}~$n(n-1)$: een product van twee factoren die onbegrensd groeien. Het antwoord 0, uit $\infty-\infty$, is fout.
\end{fullsolution}
\end{oefening}

\begin{oefening}[2]
\begin{metprent}{schildpad}
\thema{context: reeks}Een loper start 100 m achter een schildpad. De loper loopt 10 m/s, de schildpad 1 m/s. Als de loper op de startplaats van de schildpad is, is die 10 m verder. Als de loper daar is, is ze nog 1 m verder, enzovoort.
\begin{enumerate}
  \item Schrijf de tijd tot de loper de schildpad inhaalt als een oneindige som. Bereken ze.
  \item Controleer met een vergelijking.
\end{enumerate}
\end{metprent}
\begin{solution}
\textbf{a.}~$10+1+0{,}1+\cdots=\tfrac{100}{9}$ s, ongeveer 11,11 s
\end{solution}
\begin{fullsolution}
\textbf{a.}~De stukken duren 10 s, 1 s, 0,1 s, \ldots, dus $\dfrac{10}{1-0{,}1}=\dfrac{100}{9}$ s.\\
\textbf{b.}~$10t=100+t$ geeft $t=\dfrac{100}{9}$. Een oneindig aantal stukken kan dus samen een eindige tijd duren.
\end{fullsolution}
\end{oefening}

\begin{oefening}[2]
\thema{reeks: parameter}Voor welke $x \neq 1$ convergeert $1 + (x-1) + (x-1)^2 + \cdots$? Wat is dan de som?
\begin{solution}
$0<x<2$ met $x\neq 1$; de som is $\dfrac{1}{2-x}$
\end{solution}
\begin{fullsolution}
$q=x-1$, en $|x-1|<1$ geeft $0<x<2$. De som is $\dfrac{1}{1-(x-1)}=\dfrac{1}{2-x}$.
\end{fullsolution}
\end{oefening}

\begin{oefening}[2]
\thema{rekenregels: parameter}Bepaal $k$ zodat $\displaystyle\lim_{n\to+\infty} \dfrac{k\,n^2 - 1}{4n^2 + n} = 3$.
\begin{solution}
$k=12$
\end{solution}
\begin{fullsolution}
De hoogste machten geven $\dfrac{k}{4}=3$, dus $k=12$.
\end{fullsolution}
\end{oefening}

\begin{oefening}[3]
\thema{dynamisch proces: voorwaarde}Bekijk $u_n = a\,u_{n-1} + b$ met $a \neq 1$.
\begin{enumerate}
  \item Stel $L = \dfrac{b}{1-a}$ en $v_n = u_n - L$. Toon aan dat $v_n = a\,v_{n-1}$.
  \item Voor welke $a$ is de rij convergent, als $u_0 \neq L$?
  \item Neem $u_0 = 1$ en $u_n = 2u_{n-1} - 3$. De vergelijking $L = 2L - 3$ geeft $L = 3$. Is 3 de limiet?
\end{enumerate}
\begin{solution}
\textbf{b.}~$|a|<1$ \quad \textbf{c.}~nee: limiet $-\infty$
\end{solution}
\begin{fullsolution}
\textbf{a.}~Uit $L=aL+b$ volgt $v_n=au_{n-1}+b-L=au_{n-1}-aL=a\,v_{n-1}$.\\
\textbf{b.}~$v_n=v_0\cdot a^n$ met $v_0\neq 0$. Dat gaat enkel naar een eindige limiet als $|a|<1$ (want $a\neq 1$).\\
\textbf{c.}~Nee: $1,\ {-1},\ {-5},\ {-13},\ \ldots$ Hier is $u_n=3-2\cdot 2^n$, met limiet $-\infty$. De vergelijking vindt alleen de waarde die de limiet zou moeten hebben, als hij bestaat.
\end{fullsolution}
\end{oefening}

\begin{oefening}[3]
\thema{stijgend en begrensd: inductie}Gegeven: $u_1 = 1$ en $u_n = \sqrt{2 + u_{n-1}}$.
\begin{enumerate}
  \item Bereken $u_2$, $u_3$ en $u_4$ op drie decimalen.
  \item Bewijs met volledige inductie dat $0 < u_n < 2$ voor elke $n$.
  \item Toon aan dat de rij stijgend is. Kijk daarvoor naar $u_{n+1}^{\,2} - u_n^{\,2}$.
  \item Waarom is de rij convergent? Bereken de limiet.
\end{enumerate}
\begin{solution}
\textbf{a.}~1,732; 1,932; 1,983 \quad \textbf{d.}~2
\end{solution}
\begin{fullsolution}
\textbf{a.}~$\sqrt3\approx 1{,}732$, $\sqrt{3{,}732}\approx 1{,}932$, $\sqrt{3{,}932}\approx 1{,}983$.\\
\textbf{b.}~Basis: $0<1<2$. Stap: uit $0<u_k<2$ volgt $0<\sqrt{2+u_k}<\sqrt4=2$.\\
\textbf{c.}~$u_{n+1}^{\,2}-u_n^{\,2}=2+u_n-u_n^{\,2}=(2-u_n)(1+u_n)>0$, en alle termen zijn positief.\\
\textbf{d.}~Stijgend en naar boven begrensd. Voor de limiet geldt $L=\sqrt{2+L}$, dus $L^2-L-2=0$ en $L=2$ of $L=-1$. Omdat de termen positief zijn, is $L=2$.
\end{fullsolution}
\end{oefening}

\begin{oefening}[3]
\thema{definitie: rekenregel}Stel dat $\displaystyle\lim_{n\to+\infty} u_n = L$ en $k \neq 0$. Bewijs met de definitie dat $\displaystyle\lim_{n\to+\infty} k\,u_n = kL$.
\begin{solution}
Gebruik de $N$ die bij $\tfrac{\varepsilon}{|k|}$ hoort.
\end{solution}
\begin{fullsolution}
Neem $\varepsilon>0$. Dan is ook $\dfrac{\varepsilon}{|k|}>0$, dus er is een $N$ met $|u_n-L|<\dfrac{\varepsilon}{|k|}$ voor elke $n>N$. Voor die $n$ geldt $|ku_n-kL|=|k|\cdot|u_n-L|<\varepsilon$.
\end{fullsolution}
\end{oefening}

\begin{oefening}[3]
\thema{waar of niet waar}Waar of niet waar? Bewijs kort, of geef een tegenvoorbeeld.
\begin{enumerate}
  \item Elke begrensde rij is convergent.
  \item Elke stijgende rij is divergent.
  \item Elke convergente rij is begrensd.
  \item Als $u_n + v_n$ convergent is, dan zijn $u_n$ en $v_n$ allebei convergent.
  \item Als $\displaystyle\lim_{n\to+\infty} u_n = 0$, dan is $\displaystyle\lim_{n\to+\infty} u_n v_n = 0$.
\end{enumerate}
\begin{solution}
\textbf{a.}~niet waar \quad \textbf{b.}~niet waar \quad \textbf{c.}~waar \quad \textbf{d.}~niet waar \quad \textbf{e.}~niet waar
\end{solution}
\begin{fullsolution}
\textbf{a.}~Niet waar: $(-1)^n$.\\
\textbf{b.}~Niet waar: $1-\tfrac1n$.\\
\textbf{c.}~Waar: neem $\varepsilon=1$. Vanaf een bepaalde $N$ liggen alle termen tussen $L-1$ en $L+1$, en er blijven maar eindig veel andere termen over.\\
\textbf{d.}~Niet waar: $u_n=(-1)^n$ en $v_n=-(-1)^n$ geven som 0.\\
\textbf{e.}~Niet waar: $u_n=\tfrac1n$ en $v_n=n$ geven product 1.
\end{fullsolution}
\end{oefening}

\begin{oefening}[3]
\thema{reeks: telescoperen}Gegeven: $u_i = \dfrac{1}{i(i+1)}$.
\begin{enumerate}
  \item Toon aan dat $u_i = \dfrac1i - \dfrac{1}{i+1}$.
  \item Schrijf $s_n = u_1 + u_2 + \cdots + u_n$ uit en vereenvoudig.
  \item Bereken $\tfrac12 + \tfrac16 + \tfrac{1}{12} + \tfrac{1}{20} + \cdots$
\end{enumerate}
\begin{solution}
\textbf{b.}~$s_n=1-\dfrac{1}{n+1}$ \quad \textbf{c.}~1
\end{solution}
\begin{fullsolution}
\textbf{a.}~$\dfrac1i-\dfrac{1}{i+1}=\dfrac{i+1-i}{i(i+1)}=\dfrac{1}{i(i+1)}$.\\
\textbf{b.}~$s_n=\left(1-\tfrac12\right)+\left(\tfrac12-\tfrac13\right)+\cdots+\left(\tfrac1n-\tfrac{1}{n+1}\right)=1-\dfrac{1}{n+1}$: bijna alles valt weg.\\
\textbf{c.}~De rij is $u_i$, en de som is de limiet van $s_n$: 1. Deze reeks is niet meetkundig, en toch heeft ze een som.
\end{fullsolution}
\end{oefening}
\end{oefeningen}

%----------------------------------------------------------------------------------------
% KOPPEL: verbanden met andere leerstof
%----------------------------------------------------------------------------------------
\opdrachtsectie{Koppel}

\lkr{Exponentiële functies, complexe getallen en matrices zijn gezien in het 5de jaar. Kijk na of het getal $e$ al ingevoerd is. Logica en kansrekenen staan ook in het 6de jaar: kijk na of ze al aan bod kwamen, of laat die oefeningen weg.}

\begin{oefeningen}
\begin{oefening}[3]
\koppel{discreet}{Logica en bewijzen} De definitie van de limiet heeft de vorm $\forall\, \varepsilon > 0,\ \exists\, N \in \mathbb{N},\ \forall\, n > N: \ |u_n - L| < \varepsilon$.
\begin{enumerate}
  \item Schrijf de ontkenning met kwantoren. Wat betekent ze in woorden?
  \item Bewijs daarmee dat 1 niet de limiet is van $u_n = (-1)^n$.
\end{enumerate}
\begin{solution}
\textbf{a.}~$\exists\,\varepsilon>0,\ \forall\,N\in\mathbb{N},\ \exists\,n>N:\ |u_n-L|\geqslant\varepsilon$
\end{solution}
\begin{fullsolution}
\textbf{a.}~$\exists\,\varepsilon>0,\ \forall\,N\in\mathbb{N},\ \exists\,n>N:\ |u_n-L|\geqslant\varepsilon$. In woorden: er is een band waar steeds opnieuw termen buiten vallen, hoe ver je ook gaat.\\
\textbf{b.}~Neem $\varepsilon=1$. Voor elke $N$ is er een oneven $n>N$, en dan is $|u_n-1|=|-1-1|=2\geqslant 1$.
\end{fullsolution}
\end{oefening}

\begin{oefening}[2]
\koppel{analyse}{Exponentiële functies} Je belegt \eur 1 aan 100\% per jaar. Je krijgt de intrest niet één keer, maar $n$ keer per jaar, telkens $\tfrac{100}{n}$\%. Na een jaar heb je $u_n = \left(1 + \tfrac1n\right)^n$.
\begin{enumerate}
  \item Bereken $u_1$, $u_{12}$, $u_{365}$ en $u_{1\,000\,000}$ op vier decimalen.
  \item Word je zo onbeperkt rijk? Welk getal herken je?
\end{enumerate}
\begin{solution}
\textbf{a.}~2; 2,6130; 2,7146; 2,7183 \quad \textbf{b.}~nee: de limiet is $\mathrm{e}$
\end{solution}
\begin{fullsolution}
\textbf{a.}~$2^1=2$, $u_{12}\approx 2{,}6130$, $u_{365}\approx 2{,}7146$, $u_{1\,000\,000}\approx 2{,}7183$.\\
\textbf{b.}~Nee: de waarden stijgen, maar komen nooit boven $\mathrm{e}\approx 2{,}71828$. De rij is stijgend en naar boven begrensd.
\end{fullsolution}
\end{oefening}

\begin{oefening}[2]
\koppel{algebra}{Complexe getallen} Gegeven: $z = \tfrac12 + \tfrac12\mathrm{i}$.
\begin{enumerate}
  \item Bereken $|z|$. Wat is de limiet van $|z^n|$? Wat betekent dat voor $z^n$ in het complexe vlak?
  \item Bereken $1 + z + z^2 + \cdots$ met de formule van de meetkundige reeks.
\end{enumerate}
\begin{solution}
\textbf{a.}~$\tfrac{\sqrt2}{2}$; 0; de punten spiraleren naar de oorsprong \quad \textbf{b.}~$1+\mathrm{i}$
\end{solution}
\begin{fullsolution}
\textbf{a.}~$|z|=\sqrt{\tfrac14+\tfrac14}=\tfrac{\sqrt2}{2}<1$, dus $|z^n|=|z|^n\to 0$. De punten $z^n$ draaien telkens 45° en komen steeds dichter bij de oorsprong.\\
\textbf{b.}~$\dfrac{1}{1-z}=\dfrac{1}{\tfrac12-\tfrac12\mathrm{i}}=\dfrac{2}{1-\mathrm{i}}=\dfrac{2(1+\mathrm{i})}{2}=1+\mathrm{i}$.
\end{fullsolution}
\end{oefening}

\begin{oefening}[2]
\begin{metprent}{bakker}
\koppel{algebra}{Matrices} In een dorp zijn er twee bakkers. Elke maand stapt 10\% van de klanten van A over naar B, en 20\% van de klanten van B naar A. De aantallen klanten zijn $X_n = \begin{pmatrix} a_n \\ b_n \end{pmatrix}$, met $X_{n} = M X_{n-1}$ en $M = \begin{pmatrix} 0{,}9 & 0{,}2 \\ 0{,}1 & 0{,}8 \end{pmatrix}$.
\begin{enumerate}
  \item Bereken $X_1$ en $X_2$ als $X_0 = \begin{pmatrix} 300 \\ 0 \end{pmatrix}$.
  \item Zoek het evenwicht: los $X = MX$ op met $a + b = 300$. Vergelijk met $L = aL + b$.
\end{enumerate}
\end{metprent}
\begin{solution}
\textbf{a.}~$\begin{pmatrix} 270 \\ 30 \end{pmatrix}$, $\begin{pmatrix} 249 \\ 51 \end{pmatrix}$ \quad \textbf{b.}~200 en 100
\end{solution}
\begin{fullsolution}
\textbf{a.}~$X_1=\begin{pmatrix} 270 \\ 30 \end{pmatrix}$ en $X_2=\begin{pmatrix} 243+6 \\ 27+24 \end{pmatrix}=\begin{pmatrix} 249 \\ 51 \end{pmatrix}$.\\
\textbf{b.}~$a=0{,}9\,a+0{,}2\,b$ geeft $0{,}1\,a=0{,}2\,b$, dus $a=2b$. Met $a+b=300$ is $a=200$ en $b=100$. Het is dezelfde redenering, met een matrix in plaats van een getal: de limiet verandert niet meer bij nog een stap.
\end{fullsolution}
\end{oefening}

\begin{oefening}[2]
\koppel{data}{Kansrekenen} Twee spelers gooien om beurt met een dobbelsteen. Wie eerst een zes gooit, wint. Speler A begint.
\begin{enumerate}
  \item Bereken de kans dat A wint in zijn eerste, tweede en derde beurt. Welke soort rij is dat?
  \item Bereken de kans dat A het spel wint.
  \item Is het voordeel van beginnen groot?
\end{enumerate}
\begin{solution}
\textbf{a.}~$\tfrac16$, $\tfrac{25}{216}$, $\tfrac{625}{7776}$; meetkundig, $q=\tfrac{25}{36}$ \quad \textbf{b.}~$\tfrac{6}{11}$ \quad \textbf{c.}~ongeveer 54,5\% tegen 45,5\%
\end{solution}
\begin{fullsolution}
\textbf{a.}~Voor elke volgende beurt moeten beide spelers eerst geen zes gooien: telkens maal $(\tfrac56)^2=\tfrac{25}{36}$.\\
\textbf{b.}~$\dfrac{\tfrac16}{1-\tfrac{25}{36}}=\dfrac16\cdot\dfrac{36}{11}=\dfrac{6}{11}$.\\
\textbf{c.}~A wint met kans $\tfrac{6}{11}\approx 54{,}5\%$ en B met $\tfrac{5}{11}\approx 45{,}5\%$.
\end{fullsolution}
\end{oefening}
\end{oefeningen}

%----------------------------------------------------------------------------------------
\opdrachtsectie{Test jezelf}

\begin{oefeningen}
\begin{oefening}[1]\label{oef:h4-t1}
Bepaal de limiet. Is de rij convergent of divergent?
\begin{reeks}
  \item $u_n = 3 - \tfrac4n$
  \item $u_n = (-1)^n \cdot n$
  \item $u_n = 2n + 1$
\end{reeks}
\begin{solution}
\textbf{a.}~3, conv. \quad \textbf{b.}~bestaat niet, div. \quad \textbf{c.}~$+\infty$, div.
\end{solution}
\begin{fullsolution}
\textbf{a.}~$\tfrac4n\to 0$.\\
\textbf{b.}~De termen wisselen van teken en worden steeds groter in absolute waarde.\\
\textbf{c.}~De termen stijgen onbegrensd.
\end{fullsolution}
\end{oefening}

\begin{oefening}[1]\label{oef:h4-t2}
Bepaal de limiet.
\begin{reeks}
  \item $5,\ 10,\ 20,\ \ldots$
  \item $16,\ {-8},\ 4,\ \ldots$
  \item $7,\ 4,\ 1,\ \ldots$
\end{reeks}
\begin{solution}
\textbf{a.}~$+\infty$ \quad \textbf{b.}~0 \quad \textbf{c.}~$-\infty$
\end{solution}
\begin{fullsolution}
\textbf{a.}~$q=2$ en $u_1>0$.\\
\textbf{b.}~$q=-\tfrac12$, dus $|q|<1$.\\
\textbf{c.}~$v=-3$.
\end{fullsolution}
\end{oefening}

\begin{oefening}[2]\label{oef:h4-t3}
Bewijs met de definitie: $\displaystyle\lim_{n\to+\infty} \dfrac{3n+2}{n} = 3$.
\begin{solution}
Kies $N\geqslant\tfrac{2}{\varepsilon}$.
\end{solution}
\begin{fullsolution}
Neem $\varepsilon>0$. $\left|\dfrac{3n+2}{n}-3\right|=\dfrac2n<\varepsilon$ geeft $n>\dfrac{2}{\varepsilon}$. Kies $N\geqslant\dfrac{2}{\varepsilon}$: dan geldt voor elke $n>N$ dat $|u_n-3|<\varepsilon$.
\end{fullsolution}
\end{oefening}

\begin{oefening}[1]\label{oef:h4-t4}
Bereken $\displaystyle\lim_{n\to+\infty} \dfrac{2n^2 - n}{5n^2 + 3}$.
\begin{solution}
$\tfrac25$
\end{solution}
\begin{fullsolution}
Deel teller en noemer door $n^2$: $\dfrac{2-\tfrac1n}{5+\tfrac{3}{n^2}}\to\dfrac25$.
\end{fullsolution}
\end{oefening}

\begin{oefening}[2]\label{oef:h4-t5}
Bereken $\displaystyle\lim_{n\to+\infty} \dfrac{2 + \sin n}{n^2}$.
\begin{solution}
0
\end{solution}
\begin{fullsolution}
$1\leqslant 2+\sin n\leqslant 3$, dus $\dfrac{1}{n^2}\leqslant u_n\leqslant\dfrac{3}{n^2}$. Beide grenzen hebben limiet 0.
\end{fullsolution}
\end{oefening}

\begin{oefening}[2]\label{oef:h4-t6}
Toon aan dat $u_n = \dfrac{3n}{n+1}$ stijgend en naar boven begrensd is. Wat is de limiet?
\begin{solution}
Bovengrens 3; limiet 3
\end{solution}
\begin{fullsolution}
$u_n=3-\dfrac{3}{n+1}<3$. $u_{n+1}-u_n=\dfrac{3}{n+1}-\dfrac{3}{n+2}=\dfrac{3}{(n+1)(n+2)}>0$. De limiet is 3.
\end{fullsolution}
\end{oefening}

\begin{oefening}[1]\label{oef:h4-t7}
Gegeven: $u_0 = 100$ en $u_n = 0{,}6\,u_{n-1} + 20$. Bepaal het evenwicht.
\begin{solution}
50
\end{solution}
\begin{fullsolution}
$L=0{,}6\,L+20$ geeft $0{,}4\,L=20$, dus $L=50$.
\end{fullsolution}
\end{oefening}

\begin{oefening}[1]\label{oef:h4-t8}
Bereken $20 + 5 + \tfrac54 + \cdots$ en schrijf $0{,}363636\ldots$ als een breuk.
\begin{solution}
$\tfrac{80}{3}$ en $\tfrac{4}{11}$
\end{solution}
\begin{fullsolution}
$q=\tfrac14$, dus $\dfrac{20}{1-\tfrac14}=\dfrac{80}{3}$. Verder is $\dfrac{0{,}36}{1-0{,}01}=\dfrac{36}{99}=\dfrac{4}{11}$.
\end{fullsolution}
\end{oefening}
\end{oefeningen}

\herhaalschema{%
  \fv{oef:h4-t1}{vb:h4-soort}{oef:h4-soort}%
  \fv{oef:h4-t2}{vb:h4-rrmr}{oef:h4-rrmr}%
  \fv{oef:h4-t3}{vb:h4-bewijs}{oef:h4-bewijs}%
  \fv{oef:h4-t4}{vb:h4-reken}{oef:h4-ratio}%
  \fv{oef:h4-t5}{vb:h4-insluit}{oef:h4-insluit}%
  \fv{oef:h4-t6}{vb:h4-monotoon}{oef:h4-monotoon}%
  \fv{oef:h4-t7}{vb:h4-dyn}{oef:h4-evenwicht}%
  \fv{oef:h4-t8}{vb:h4-reeks}{oef:h4-reeks}}

\ifzelfstudie
\section*{Maak je eigen samenvatting}
Sluit dit boek en schrijf in je schrift: wat convergent en divergent betekent, de formele definitie, wanneer een rekenkundige en een meetkundige rij convergeren, hoe je een limiet berekent, de insluitstelling, de eigenschap van stijgende en begrensde rijen, hoe je een evenwicht vindt, en de som van een meetkundige reeks. Vergelijk daarna met de kaders.
\fi


\Closesolutionfile{oplossingen}
\Closesolutionfile{uitgewerkteoplossingen}

% De leerkrachten- en zelfstudieversie krijgen de antwoorden meteen achteraan
\ifuitwerking
\clearpage\global\opdrachtdeelfalse
\section*{Antwoorden}
\begin{multicols}{2}\footnotesize
\input{oplossingen}
\end{multicols}
\fi

\else

\ifdefstring{\uitvoer}{numeriek}{%
\chapter*{Numerieke oplossingen}
\begin{multicols}{2}
\input{oplossingen}
\end{multicols}
}{%
\chapter*{Uitgewerkte oplossingen}
\input{uitgewerkteoplossingen}
}

\fi

\end{document}
